\documentclass[11pt,a4paper]{report}

% ============================================================
% PACKAGES
% ============================================================
\usepackage[T1]{fontenc}
\usepackage[utf8]{inputenc}
\usepackage{graphicx}
% parskip.sty replacement (only sets paragraph spacing; inlined to avoid
% the dependency):
\setlength{\parindent}{0pt}
\setlength{\parskip}{6pt plus 2pt minus 1pt}
\usepackage{latexsym}
\usepackage{amsmath, amssymb, amsthm, bm}
\usepackage{booktabs}
\usepackage{longtable}
\usepackage{array}
\usepackage{enumitem}
\usepackage{xcolor}
\usepackage{tcolorbox}
\tcbuselibrary{skins,breakable}
\usepackage[a4paper,left=2.8cm,right=2.8cm,top=2.8cm,bottom=2.8cm]{geometry}
\usepackage{zi4} % monospace for code
\usepackage[
    colorlinks=true,
    linkcolor=blue!50!black,
    citecolor=blue!50!black,
    urlcolor=blue!50!black,
    bookmarksnumbered=true,
    bookmarksopen=true
]{hyperref}
\usepackage{cleveref}

% ============================================================
% THEOREM ENVIRONMENTS
% ============================================================
\newtheorem{theorem}{Theorem}[chapter]
\newtheorem{proposition}[theorem]{Proposition}
\newtheorem{lemma}[theorem]{Lemma}
\newtheorem{corollary}[theorem]{Corollary}
\theoremstyle{definition}
\newtheorem{definition}[theorem]{Definition}
\newtheorem{example}[theorem]{Example}
\theoremstyle{remark}
\newtheorem{remark}[theorem]{Remark}
\newtheorem{notation}[theorem]{Notation}
\newtheorem{caution}[theorem]{Caution}

% ============================================================
% CODE-CORRESPONDENCE BOX
% ============================================================
% Usage: \codebox{Short title}{Math object}{File/function}{Shape}{Interpretation}
\newtcolorbox{codebox}[1][]{
  breakable, enhanced, colback=gray!5, colframe=gray!60!black,
  boxrule=0.4pt, arc=1mm, left=2mm, right=2mm, top=1mm, bottom=1mm,
  title={\small\textbf{Code correspondence}#1},
  fonttitle=\small\bfseries, coltitle=black,
  attach boxed title to top left={yshift=-2mm, xshift=2mm},
  boxed title style={colback=gray!5, boxrule=0pt},
}
% (code excerpts use plain \texttt{}/verbatim, not the listings package,
% so no \lstset configuration is needed)

% ============================================================
% OPERATORS AND MACROS
% ============================================================
\DeclareMathOperator{\Exp}{Exp}
\DeclareMathOperator{\Log}{Log}
\DeclareMathOperator{\so}{\mathfrak{so}}
\DeclareMathOperator{\se}{\mathfrak{se}}
\DeclareMathOperator{\tr}{tr}
\DeclareMathOperator{\diag}{diag}
\DeclareMathOperator{\Var}{Var}
\DeclareMathOperator{\Cov}{Cov}
\DeclareMathOperator*{\argmin}{arg\,min}
\DeclareMathOperator*{\argmax}{arg\,max}
\DeclareMathOperator{\Attn}{Attn}
\DeclareMathOperator{\softmax}{softmax}
\DeclareMathOperator{\LN}{LN}
\DeclareMathOperator{\MLP}{MLP}
\DeclareMathOperator{\GELU}{GELU}
\DeclareMathOperator{\SiLU}{SiLU}
\DeclareMathOperator{\Unif}{Unif}
\DeclareMathOperator{\KL}{KL}
\DeclareMathOperator{\supp}{supp}
\newcommand{\dd}{\mathrm{d}}
\newcommand{\RR}{\mathbb R}
\newcommand{\EE}{\mathbb E}
\newcommand{\NN}{\mathbb N}
\newcommand{\ZZ}{\mathbb Z}
\newcommand{\CC}{\mathbb C}
\newcommand{\su}{\mathfrak{su}}
\newcommand{\PP}{\mathbb P}
\newcommand{\ind}{\mathbf 1}
\newcommand{\hatop}{\,\widehat{\phantom{x}}\,}
\newcommand{\eps}{\varepsilon}
\newcommand{\file}[1]{\texttt{#1}}
\newcommand{\pyfunc}[1]{\texttt{#1}}
\newcommand{\shape}[1]{\texttt{[#1]}}

\title{SigmaDock and SigmaFlow: A Self-Contained Mathematical and
Computational Monograph}
\author{Technical reference document --- companion to the SigmaFlow project}
\date{}

\begin{document}

% ============================================================
% FRONT MATTER
% ============================================================
\begin{titlepage}
\centering
\vspace*{2cm}
{\Huge\bfseries SigmaDock and SigmaFlow\par}
\vspace{0.8cm}
{\Large A Self-Contained Mathematical and Computational Monograph\par}
\vspace{0.4cm}
{\large From probability theory and Riemannian geometry to \\
$SE(3)$-equivariant diffusion and Riemannian flow matching \\
for protein--ligand docking\par}
\vspace{2.5cm}
{\large Technical reference document\par}
\vspace{0.3cm}
{\large companion to the SigmaFlow research project\par}
\vfill
{\large \today\par}
\end{titlepage}

\clearpage
\begin{abstract}
\noindent
This monograph is a from-scratch, self-contained derivation and code-verified
description of \textbf{SigmaDock}, an $SE(3)$-equivariant diffusion model for
protein--ligand pose generation, and \textbf{SigmaFlow}, its Riemannian
flow-matching counterpart. It assumes only real analysis, linear algebra,
elementary probability theory, and general mathematical maturity. Every other
concept --- Riemannian manifolds, Lie groups, Brownian motion on manifolds,
score-based diffusion, flow matching, spherical harmonics, Wigner-$D$
matrices, Clebsch--Gordan coefficients, and the Equiformer/EquiformerV2
architecture --- is introduced from first principles, motivated, and derived
in full. Every mathematical construction is, wherever it is used by
SigmaDock or SigmaFlow, tied explicitly to the concrete implementation: file,
class, function, tensor shape, and formula. The document distinguishes
throughout between (i) idealised theory, (ii) what the SigmaDock method
description states, (iii) what the SigmaDock reference implementation
actually computes, and (iv) what SigmaFlow changes, and flags explicitly any
point at which these four do not coincide or could not be verified with
certainty from the available source.
\end{abstract}

\clearpage
\tableofcontents
\clearpage

% ============================================================
\chapter*{How to read this document}
\addcontentsline{toc}{chapter}{How to read this document}
% ============================================================

This document is organised as a linear course: later chapters depend on
earlier ones, and no chapter uses a concept before that concept has been
defined \emph{in this document}. If you already know a topic (say, Lie
groups, or diffusion models), you may skip the corresponding chapter, but
you should still skim its \textbf{Notation} boxes, since notation is fixed
globally and is not repeated in later chapters.

Every new mathematical construction is introduced in six steps, repeated as
a pattern throughout: \emph{(1) motivation} --- why do we need this object at
all; \emph{(2) intuition} --- an informal picture; \emph{(3) formal
definition}; \emph{(4) derivation} of its key properties; \emph{(5) a small
worked example}; \emph{(6) application} to SigmaDock or SigmaFlow. Results
that are used later are stated as numbered theorems, propositions, or
lemmas; all proofs that are not immediate are carried out in full, not
sketched or deferred to a citation, \emph{unless} the result is a classical
one whose proof would require substantially more background than this
document provides (e.g.\ existence/uniqueness for SDEs), in which case this
is stated explicitly and a precise reference is given rather than an
unstated leap.

Whenever a mathematical object is realised in the SigmaDock/SigmaFlow
codebase, a grey \textbf{Code correspondence} box follows immediately, of
the form
\begin{codebox}
\textbf{Mathematical object:} $v_t(x)\in T_x\mathcal M$ \\
\textbf{Implementation:} \file{sigma\_flow\_generator.py}::\pyfunc{calc\_vector\_field} \\
\textbf{Tensor shape:} \shape{B x F, 3} \\
\textbf{Interpretation:} translational velocity, one $\RR^3$ vector per ligand
fragment in the batch.
\end{codebox}
These boxes are the load-bearing content of the document in the sense that
they are what makes the mathematics falsifiable against the actual program
that runs on real data. Every box has been checked against the source files
named in it, at the time of writing; file paths are given relative to the
repository roots \file{SigmaDock/} (the read-only original diffusion
reference) and \file{SigmaFlow\_Development/} (the flow-matching
conversion), or the checked-out reproduction clone
\file{SigmaDock\_Reproduction\_JulianMueller/sigmadock/} where the two
happen to coincide.

\begin{caution}[On uncertainty]
Some implementation details are genuinely ambiguous from static reading of
the code (for instance, a numerical constant whose provenance is not
commented, or a design choice for which no docstring states the intent).
Wherever this occurs, the text says so explicitly, using the phrase
\emph{``not determinable from the code as written; interpretation:''}
followed by the most plausible reading. Nothing in this document is
invented to fill such a gap silently.
\end{caution}

\section*{Global notation}
\addcontentsline{toc}{section}{Global notation}

Notation is introduced at first use and collected here for reference; this
table is updated as new symbols are introduced and should be treated as
living, not exhaustive at first read.

\begin{center}
\begin{longtable}{p{3.6cm}p{10.5cm}}
\toprule
Symbol & Meaning \\
\midrule
\endfirsthead
\toprule
Symbol & Meaning \\
\midrule
\endhead
$\RR^d$ & real Euclidean space of dimension $d$ \\
$\mathcal M$ & smooth (Riemannian) manifold \\
$T_x\mathcal M$ & tangent space of $\mathcal M$ at $x$ \\
$T\mathcal M$ & tangent bundle, $\bigsqcup_{x} T_x\mathcal M$ \\
$g_x(\cdot,\cdot)$, $\|\cdot\|_g$ & Riemannian metric and induced norm at $x$ \\
$\exp_x$, $\log_x$ & Riemannian exponential map and its (local) inverse at $x$ \\
$G$, $\mathfrak g$ & a Lie group and its Lie algebra $\mathfrak g = T_eG$ \\
$SO(3)$, $\so(3)$ & the rotation group and its Lie algebra (skew-symmetric $3\times3$ matrices) \\
$SE(3)$ & the group of rigid motions of $\RR^3$ \\
$(\cdot)^{\wedge}$, $(\cdot)^{\vee}$ & hat and vee maps, $\RR^3 \leftrightarrow \so(3)$ \\
$\mathcal M_N = SE(3)^N$ & product pose manifold of $N$ rigid fragments \\
$p_{\mathrm{data}}$, $p_{\mathrm{init}}$ & data distribution and prior (source) distribution \\
$p_t$, $p_t(\cdot\mid z)$ & marginal / conditional probability path at time $t\in[0,1]$ \\
$u_t$, $u_t(\cdot\mid z)$ & marginal / conditional (velocity) vector field \\
$u_t^\theta$, $s_\theta$ & learned vector field network / learned score network \\
$\nabla_x \log p(x)$ & (Riemannian) score of a density $p$ \\
$W_t$ & standard Wiener process (Brownian motion in $\RR^d$) \\
$B_t^{\mathcal M}$ & Brownian motion on a Riemannian manifold $\mathcal M$ \\
$Y^{\ell}_m$ & real spherical harmonic of degree $\ell$, order $m$ \\
$D^{\ell}(Q)$ & Wigner-$D$ matrix of degree $\ell$ for $Q \in SO(3)$ \\
$C^{(\ell,m)}_{(\ell_1,m_1)(\ell_2,m_2)}$ & Clebsch--Gordan coefficient \\
$\otimes_{\mathrm{cg}}$ & Clebsch--Gordan (steerable) tensor product \\
$\mathcal N(\mu,\Sigma)$ & Gaussian/normal distribution with mean $\mu$, covariance $\Sigma$ \\
$\mathrm{IGSO}(3)(R_0,\sigma^2)$ & isotropic Gaussian distribution on $SO(3)$, centred at $R_0$, scale $\sigma$ \\
$\mathcal U(\mathcal M)$ & uniform (Haar, when $\mathcal M$ is a compact group) measure on $\mathcal M$ \\
$B,F,N,A$ & batch size, number of fragments, number of atoms/nodes, generic index \\
$t$ & generation time, $t\in[0,1]$, with convention $p_0 = p_{\mathrm{init}}$, $p_1 = p_{\mathrm{data}}$ \\
$s = 1-t$ & noising time, used only when discussing the forward/noising process \\
\bottomrule
\end{longtable}
\end{center}

\begin{remark}[Time convention, fixed once and for all]
\label{rem:time-convention-global}
Throughout this document, $t \in [0,1]$ denotes \emph{generation time}: $t=0$
is the source/prior distribution ($p_{\mathrm{init}}$, e.g.\ pure noise), and
$t=1$ is the data distribution ($p_{\mathrm{data}}$, e.g.\ the true bound
pose). All sampling procedures therefore integrate \emph{forward} in $t$,
from $t=0$ to $t=1$. Where a \emph{noising} (forward diffusion) process is
discussed, it runs in the reversed variable $s = 1-t$, i.e.\ it starts at
data ($s=0$) and ends at noise ($s=1$). This convention is fixed because the
diffusion and flow-matching literatures disagree with each other (and
sometimes with themselves) on which direction is ``forward'', and because
the SigmaDock and SigmaFlow codebases must each be checked individually
against it -- this is done explicitly in \Cref{chap:sigmadock,chap:sigmaflow}.
\end{remark}

\clearpage

% ================================================================
\part{Preliminaries}
% ================================================================

\chapter{The Protein--Ligand Docking Problem}
\label{chap:docking}

Before any mathematics is introduced, we fix precisely what SigmaDock and
SigmaFlow are trying to compute. This chapter is deliberately non-technical
in the sense that it uses no probability theory or differential geometry
yet; its purpose is to pin down the \emph{objects} (proteins, ligands,
poses, degrees of freedom) and the \emph{task} (what must be predicted, and
how success is measured) so that every later mathematical construction has
an unambiguous referent.

\section{Proteins and ligands}

\begin{definition}[Protein, for docking purposes]
\label{def:protein}
For the purposes of this document, a \emph{protein} is a rigid, known
three-dimensional structure: a finite set of atoms $\{(y_j, \zeta_j)\}_{j=1}^{M}$,
where $y_j \in \RR^3$ is the atom's position and $\zeta_j$ collects its
chemical identity (element, residue type, atomic partial charge if used).
We do not model protein flexibility: $y_j$ is held fixed throughout.
\end{definition}

This is a modelling choice, not a physical fact --- real proteins move --- and
it corresponds to the standard \emph{rigid-receptor docking} setting used by
essentially all methods discussed in this document, including SigmaDock. The
\emph{binding pocket} is the subset of protein atoms within some distance of
the ligand; SigmaDock and SigmaFlow both assume the pocket (equivalently, a
point defining its centre) is known in advance, so that the generative model
only has to search over ligand poses \emph{within} a neighbourhood of a given
location, not over the protein surface.

\begin{definition}[Ligand, for docking purposes]
\label{def:ligand}
A \emph{ligand} is a small molecule, represented as a molecular graph
$\mathcal G^{2D} = (\mathcal V, \mathcal B)$: atoms $\mathcal V$ with chemical
labels (element, formal charge, hybridisation, aromaticity, chirality tag),
and bonds $\mathcal B \subseteq \mathcal V \times \mathcal V$ with labels
(bond order, aromaticity, whether the bond lies in a ring, stereochemistry).
A \emph{conformer} of the ligand is an assignment of 3D coordinates
$x = (x_1,\dots,x_n) \in \RR^{n\times3}$, $n=|\mathcal V|$, to its atoms.
\end{definition}

Crucially, the molecular graph $\mathcal G^{2D}$ is \emph{given} (it is
computed once, deterministically, from the ligand's chemical structure,
e.g.\ from a SMILES string), whereas the 3D coordinates $x$ are what a
docking model must predict. This is the sense in which docking is a
\emph{structure prediction} problem: the input already specifies which atoms
exist and how they are chemically connected; the output is where they sit in
space.

\section{Poses and degrees of freedom}
\label{sec:pose-dof}

\begin{definition}[Pose]
\label{def:pose}
A \emph{pose} of a ligand with $n$ atoms, relative to a fixed protein, is an
assignment of 3D coordinates $x \in \RR^{n \times 3}$ to its atoms, in the
same coordinate frame as the protein.
\end{definition}

Naively, a pose has $3n$ real degrees of freedom (three coordinates per
atom). This naive count is almost never the right one to reason about,
because it ignores that the ligand is a \emph{chemical object}, not an
arbitrary point cloud: bond lengths and bond angles are, to excellent
approximation, fixed by chemistry (they fluctuate by hundredths of an
\r{A}ngstr\"om around equilibrium values, an order of magnitude below the
accuracy that any docking method targets), and only a small number of
internal coordinates vary substantially. This observation is what licenses
essentially every design decision reviewed in \Cref{chap:sigmadock}, so we
make it precise now.

The internal geometry of a molecule can be decomposed into:
\begin{itemize}
\item \textbf{Bond lengths} $\|x_i - x_j\|$ for bonded pairs $(i,j)\in\mathcal B$:
essentially fixed by chemistry (covalent bond lengths are determined, to
first approximation, by the identities of the two bonded atoms and the bond
order).
\item \textbf{Bond angles} $\angle(x_i,x_j,x_k)$ for two bonds sharing atom
$j$: also close to fixed, determined by the hybridisation of atom $j$.
\item \textbf{Dihedral (torsion) angles} $\phi_{ijkl}$, defined for four
atoms $i$--$j$--$k$--$l$ connected by three consecutive bonds, as the angle
between the plane through $(i,j,k)$ and the plane through $(j,k,l)$: these
are, for single (non-ring) bonds, comparatively free to rotate, since
rotation about a single bond changes no bond length and no bond angle.
\end{itemize}

A bond about which the dihedral angle can vary freely without breaking the
molecule (i.e.\ a single, non-ring, non-terminal bond) is called a
\emph{rotatable bond}. If a ligand has $k$ rotatable bonds, and bond lengths
and angles are treated as exactly fixed, then the \emph{internal}
conformation of the ligand is fully described by $k$ angles
$\bm\phi = (\phi_1,\dots,\phi_k) \in \mathbb T^k$, where
$\mathbb T = \RR/2\pi\mathbb Z$ is the circle (the space of angles); we
define and use $\mathbb T^k$ formally in \Cref{sec:torus}. The
\emph{external} pose of the (now internally fixed) ligand is then a single
rigid motion --- a rotation and a translation --- applied to the whole
molecule, contributing $3+3=6$ further degrees of freedom (formalised as an
element of the group $SE(3)$ in \Cref{chap:groups}).

\begin{center}
\begin{tcolorbox}[colback=blue!4,colframe=blue!40!black,title=\textbf{Degrees of freedom of a ligand pose},boxrule=0.4pt]
\[
\underbrace{3}_{\text{translation}}
+
\underbrace{3}_{\text{global rotation}}
+
\underbrace{k}_{\text{torsions}}
\;=\; k+6
\]
against a naive count of $3n$, where typically $n \gg k$: for a
drug-like ligand with $n\approx 30$ heavy atoms one might have $k\approx 5$,
so the chemically-informed parameterisation has roughly $11$ degrees of
freedom against a naive $90$.
\end{tcolorbox}
\end{center}

This factorisation --- translation, rotation, torsions --- is the single most
important structural fact about the docking problem, and essentially every
design choice in SigmaDock (and hence SigmaFlow) can be read as a decision
about \emph{how literally to take it}. \Cref{sec:sigmadock-state-space} works
out three concrete choices (all-atom, torsional, fragment-based) and the
consequences of each; the mathematics needed to state those consequences
precisely (Riemannian volume forms, product measures) is developed in
\Cref{chap:groups}.

\begin{remark}[What is \emph{not} varied]
The receptor is rigid (\Cref{def:protein}). Bond lengths and bond angles are
either held exactly fixed, or --- as we will see is SigmaDock's actual
choice --- allowed to fluctuate only implicitly, by resampling a fresh
low-energy conformer rather than by being explicit generative degrees of
freedom. No docking model discussed here generates bond lengths or angles
directly.
\end{remark}

\section{What a docking model predicts, and what it does not}
\label{sec:docking-tasks}

The docking literature distinguishes several related but distinct tasks,
and conflating them is a common source of confusion when reading papers or
comparing methods.

\begin{description}
\item[Pose generation.] Given $\mathcal G_{\mathrm{dock}}$ (ligand graph,
protein structure, binding pocket), produce one or more candidate poses $x$.
This is the task SigmaDock and SigmaFlow solve; it is a \emph{generative}
task (a probabilistic model is used to \emph{sample} candidates), not a
search or optimisation task, though it can be built on top of either.
\item[Pose scoring.] Given a candidate pose $x$, output a scalar estimate of
its quality (e.g.\ predicted binding free energy, or a proxy such as a
learned or physics-based score). This is used to \emph{rank} multiple
generated poses (see the ranking heuristic in
\Cref{sec:sigmadock-ranking}), but it does not by itself produce new poses.
\item[Ranking.] Given a set of candidate poses (e.g.\ multiple i.i.d.\ draws
from a generative model, each called a \emph{seed}), order them by predicted
quality so that a downstream user can inspect the top-$k$.
\item[Confidence estimation.] A calibrated estimate of \emph{how likely a
specific prediction is to be correct}, as opposed to a relative score used
only for ranking among a fixed candidate set. \Cref{chap:sigmaflow} discusses
to what extent SigmaFlow's flow-matching formulation admits a
principled confidence signal via the ODE log-likelihood
(\Cref{sec:likelihood-confidence}), and states explicitly which parts of
this are implemented in the current codebase and which are conceptual only.
\end{description}

SigmaDock (\Cref{chap:sigmadock}) performs pose generation and a simple,
non-learned ranking heuristic; it explicitly does \emph{not} include a
learned confidence model or post-hoc energy minimisation
(\Cref{sec:sigmadock-ranking}), a deliberate design choice intended to
isolate how much of its accuracy is attributable to the generative model
itself, as opposed to a downstream correction step. SigmaFlow inherits this
scope unless stated otherwise.

\section{Evaluation: RMSD and structural validity}
\label{sec:evaluation}

Two families of metrics are used throughout this document to talk about
``how good'' a predicted pose is.

\begin{definition}[Root-mean-square deviation, RMSD]
\label{def:rmsd}
For a predicted pose $\hat x \in \RR^{n\times3}$ and a reference (experimentally
determined, ``ground truth'') pose $x^\star \in \RR^{n\times3}$ of the same
ligand, with atoms in the same order,
\[
\mathrm{RMSD}(\hat x, x^\star)
=
\left(
\frac1n \sum_{i=1}^n \|\hat x_i - x^\star_i\|_2^2
\right)^{1/2} .
\]
A pose is conventionally called a \emph{success} (in the re-docking
setting of this document) if $\mathrm{RMSD}(\hat x, x^\star) < 2\,$\r A after
optimal rigid alignment is \emph{not} applied (i.e.\ RMSD is computed
directly in the shared binding-site frame) --- we return to this point,
including the alternative of Kabsch-aligned RMSD, when we discuss the
specific evaluation used for SigmaDock/SigmaFlow comparisons.
\end{definition}

RMSD is a single number and, being an average, is insensitive to
\emph{local} chemical implausibility: a pose can have low RMSD to the
reference while containing an impossibly short bond or an overlapping ring,
if the error is averaged out over the rest of the molecule. This motivates a
second, complementary family of metrics.

\begin{definition}[PoseBusters-style validity checks, informal]
\label{def:posebusters}
A pose $\hat x$ is subjected to a battery of deterministic, chemistry-based
checks --- e.g.\ are all bond lengths within a tolerance of their expected
values, are all bond angles chemically reasonable, is stereochemistry
preserved, do any atoms sterically clash with each other or with the
protein, does the ligand's internal energy (under a cheap force field)
exceed a threshold --- each of which returns pass/fail. A pose
\emph{passes PoseBusters} if it passes every check in the battery.
\end{definition}

We will use the phrase \emph{PoseBusters pass rate} for the fraction of
generated poses (over some fixed evaluation set) that pass every check, and
we will occasionally need per-check pass rates (e.g.\ specifically the
\texttt{bond\_lengths} or \texttt{bond\_angles} check) when discussing
fragment-boundary artefacts; these are named explicitly whenever used.

\begin{remark}[Why both metrics matter, and why neither suffices alone]
RMSD measures \emph{closeness to the correct answer}; PoseBusters measures
\emph{physical plausibility}, independent of whether the pose is close to
any particular reference. A model can score well on one and poorly on the
other: for instance, a model whose fragment-level rigid-body placement is
accurate (low RMSD) can still produce chemically broken bonds precisely at
the boundaries between rigid fragments if the relative placement of two
fragments is even slightly off, since a rigid transformation exactly
preserves lengths \emph{within} a fragment but says nothing about the
distance \emph{between} fragments (see the discussion of triangulation in
\Cref{sec:triangulation-precise}, and the analysis of transition-bond lengths at
fragment boundaries in \Cref{chap:sigmaflow}). PoseBusters, being a binary
per-molecule pass/fail criterion, is additionally insensitive to
\emph{how far} a check is failed by, a limitation we discuss quantitatively
when it becomes relevant.
\end{remark}

This completes the non-mathematical scaffolding: we now know what a pose
is, which of its coordinates are genuinely free, what a docking model is
asked to produce, and how its output is judged. \Cref{chap:prob-foundations}
turns to the probability theory required to make ``a generative model for
poses'' a precise mathematical statement.

\clearpage

% ================================================================
\chapter{Probability and Generative Modelling Foundations}
\label{chap:prob-foundations}
% ================================================================

This chapter collects the probabilistic vocabulary used everywhere in the
rest of the document. None of it is specific to docking; it is the common
language in which \Cref{chap:diffusion,chap:flow-matching} will phrase
diffusion and flow matching. Readers comfortable with densities,
change-of-variables, KL divergence and maximum likelihood may skip to
\Cref{sec:generative-modelling-intro}, returning here only if a later
citation of a numbered result is unclear.

\section{Densities, conditional densities, and pushforward}

We work throughout with random variables taking values in $\RR^d$ (this
chapter) or, from \Cref{chap:groups} onward, in a Riemannian manifold
$\mathcal M$; every definition below is stated for $\RR^d$ and is repeated,
without change to the underlying idea, for manifolds once the necessary
geometric vocabulary (volume forms) is available.

\begin{definition}[Density]
\label{def:density}
A random variable $X$ taking values in $\RR^d$ has \emph{density} $p:\RR^d\to[0,\infty)$
with respect to Lebesgue measure if
\[
\PP(X\in A) = \int_A p(x)\,dx
\qquad\text{for every (measurable) } A\subseteq\RR^d,
\]
in particular $\int_{\RR^d}p(x)\,dx=1$.
\end{definition}

\begin{definition}[Conditional density]
\label{def:conditional-density}
For a pair $(X,Z)$ with joint density $p(x,z)$ and marginal $p(z)=\int p(x,z)\,dx>0$,
the \emph{conditional density} of $X$ given $Z=z$ is $p(x\mid z)=p(x,z)/p(z)$.
\end{definition}

The identity that we will use over and over, sometimes without comment, is
the following elementary but consequential rearrangement of
\Cref{def:conditional-density}.

\begin{proposition}[Marginalisation by conditioning]
\label{prop:marginalisation}
If $Z\sim p_{\mathrm{data}}$ and, given $Z=z$, $X\sim p(\cdot\mid z)$, then $X$
marginally has density
\[
p(x) = \int p(x\mid z)\,p_{\mathrm{data}}(z)\,dz .
\]
\end{proposition}

\begin{proof}
Immediate from $p(x,z)=p(x\mid z)p_{\mathrm{data}}(z)$ and integrating $z$ out.
\end{proof}

\Cref{prop:marginalisation} looks almost too simple to name, but it is the
single mechanism by which both diffusion models
(\Cref{chap:diffusion}) and flow matching (\Cref{chap:flow-matching})
convert an \emph{intractable} learning problem (regress against a marginal
quantity that is only available as an integral over all of
$p_{\mathrm{data}}$) into a \emph{tractable} one (regress against a
\emph{conditional} quantity that is available in closed form for one data
point $z$ at a time): both frameworks first construct a family
$p_t(\cdot\mid z)$ that is easy to sample and differentiate, then use
\Cref{prop:marginalisation} to reason about the (unavailable) marginal
$p_t$, and finally prove --- this is the content of
\Cref{thm:cfm-equivalence-scratch,thm:dsm-equivalence} below --- that a regression
loss built from the conditional quantity has the \emph{same minimiser} as
one built from the marginal quantity, sidestepping the need to ever
evaluate the intractable integral.

\begin{definition}[Pushforward measure]
\label{def:pushforward}
Let $X$ have density $p$ on $\RR^d$ and let $\phi:\RR^d\to\RR^d$ be a
(measurable) map. The \emph{pushforward} of $p$ under $\phi$, written
$\phi_{\#}p$, is the distribution of $Y=\phi(X)$:
\[
\PP(Y\in A) = \PP(\phi(X)\in A) = \PP\bigl(X\in\phi^{-1}(A)\bigr)
\qquad\text{for measurable } A.
\]
\end{definition}

When $\phi$ is a diffeomorphism (smooth, invertible, smooth inverse), the
pushforward has an explicit density, given by the classical
change-of-variables formula.

\begin{theorem}[Change of variables]
\label{thm:change-of-variables}
Let $X$ have density $p_X$ and let $\phi$ be a diffeomorphism of an open
subset of $\RR^d$ onto its image. Then $Y=\phi(X)$ has density
\[
p_Y(y) = p_X\bigl(\phi^{-1}(y)\bigr)\,
\left|\det D\phi^{-1}(y)\right|
=
\frac{p_X\bigl(\phi^{-1}(y)\bigr)}{\left|\det D\phi(\phi^{-1}(y))\right|},
\]
where $D\phi$ is the Jacobian matrix of $\phi$.
\end{theorem}

We take \Cref{thm:change-of-variables} as known from multivariable calculus
and do not reprove it; what matters for us is how it is used. In
\Cref{sec:cnf} we will apply it to a map $\phi$ that is itself the
time-$t$ flow of an ordinary differential equation, and derive from it the
\emph{continuity equation} that governs how a probability density evolves
under such a flow --- the central object of flow matching.

\begin{example}[A one-dimensional sanity check]
Let $X\sim\mathcal N(0,1)$ and $\phi(x)=2x$ (so $Y=2X$). Then
$\phi^{-1}(y)=y/2$, $D\phi(x)=2$, and \Cref{thm:change-of-variables} gives
$p_Y(y) = p_X(y/2)/2 = \tfrac1{2\sqrt{2\pi}}e^{-y^2/8}$, i.e.\ $Y\sim\mathcal N(0,4)$,
as expected from elementary Gaussian scaling.
\end{example}

\section{Score functions}
\label{sec:score-intro}

\begin{definition}[Score]
\label{def:score}
The \emph{score} of a density $p$ with $p(x)>0$ on its support is
\[
\nabla_x \log p(x) \;\in\;\RR^d,
\]
the gradient (with respect to $x$) of the log-density.
\end{definition}

The score is the direction in which $\log p$, and hence $p$ itself,
increases fastest; equivalently it is the vector field that would define
gradient ascent on $\log p$. Two properties of the score matter
disproportionately for what follows.

\begin{remark}[The score does not need a normalising constant]
\label{rem:score-no-const}
If $p(x)=\tilde p(x)/Z$ for an unnormalised $\tilde p$ and a normalising
constant $Z=\int \tilde p$, then $\nabla_x\log p(x)=\nabla_x\log\tilde p(x)$,
since $\nabla_x\log Z = 0$. The score is therefore computable from an
unnormalised density, which ordinary likelihood evaluation is not. This is
one of the two reasons (the other being \Cref{thm:langevin} below) that
score-based methods can be trained without ever evaluating $Z$.
\end{remark}

\begin{remark}[The score of a Gaussian, used repeatedly]
\label{rem:gaussian-score}
For $p(x)=\mathcal N(x\mid\mu,\sigma^2I)$,
\[
\log p(x) = -\frac{\|x-\mu\|^2}{2\sigma^2} + \text{const},
\qquad
\nabla_x\log p(x) = -\frac{x-\mu}{\sigma^2} .
\]
This single formula is, after substitution of the appropriate $\mu,\sigma$,
the target of every Euclidean denoising-score-matching loss used in
\Cref{chap:diffusion}.
\end{remark}

\section{Why the score suffices for sampling: Langevin dynamics}
\label{sec:langevin}

It is not obvious \emph{a priori} that knowing $\nabla_x\log p$ everywhere
is enough to sample from $p$ --- after all, the score is a strictly weaker
piece of information than $p$ itself (\Cref{rem:score-no-const}). The
classical result that resolves this is Langevin dynamics, and understanding
why it works is what motivates the entire diffusion-model construction of
\Cref{chap:diffusion}.

We first need Brownian motion in $\RR^d$; a manifold version is given in
\Cref{chap:diffusion} once the necessary geometry is available, but the flat
case can be defined immediately.

\begin{definition}[Standard Wiener process]
\label{def:wiener}
A \emph{standard Wiener process} (Brownian motion) $(W_\tau)_{\tau\ge0}$ in
$\RR^d$ is a continuous-time stochastic process with $W_0=0$, independent
increments, and $W_{\tau+h}-W_\tau\sim\mathcal N(0,hI_d)$ for all $\tau,h\ge0$.
\end{definition}

\begin{definition}[Stochastic differential equation, informal]
\label{def:sde}
Given a drift $f:\RR^d\times\RR_{\ge0}\to\RR^d$ and a diffusion coefficient
$g:\RR_{\ge0}\to\RR_{\ge0}$, the SDE
\[
dX_\tau = f(X_\tau,\tau)\,d\tau + g(\tau)\,dW_\tau
\]
denotes, informally, the continuous-time limit of the discrete update
$X_{\tau+h} = X_\tau + f(X_\tau,\tau)h + g(\tau)\sqrt h\,z$,
$z\sim\mathcal N(0,I_d)$, as $h\to0$. We use SDEs only as a organising
language for such limits and as a source of update rules to discretise; we
do not develop It\^o calculus, and every fact about SDEs used below is
either a discretisation recipe (which needs no further theory) or is quoted,
with a precise reference, as a classical result.
\end{definition}

\begin{theorem}[Langevin dynamics; invariance of the target density]
\label{thm:langevin}
Let $p(x)\propto e^{-U(x)}$ for a potential $U$ with $\nabla U$ Lipschitz and
$\int e^{-U}<\infty$ (sufficient regularity to guarantee existence,
uniqueness and ergodicity of the SDE below; see, e.g., Pavliotis,
\emph{Stochastic Processes and Applications}, 2014, Ch.~4). Then the
\emph{overdamped Langevin SDE}
\[
dX_\tau = -\nabla U(X_\tau)\,d\tau + \sqrt2\,dW_\tau
\]
has $p$ as its unique stationary distribution: if $X_0\sim p$ then $X_\tau\sim p$
for all $\tau\ge0$, and $\mathrm{Law}(X_\tau)\to p$ as $\tau\to\infty$ from
any initial distribution with support contained in that of $p$, under mild
additional conditions.
\end{theorem}

We do not reprove \Cref{thm:langevin} (it is a statement about the
stationary distribution of a diffusion process via its Fokker--Planck
equation, i.e.\ it requires exactly the machinery we introduce, in the
special case needed for our models, in \Cref{sec:sde-perspective}); we take
it as a classical fact and note only the substitution that makes it useful
here. Setting $U=-\log p_{\mathrm{data}}$ gives $-\nabla U = \nabla_x\log p_{\mathrm{data}}$,
so \Cref{thm:langevin} says that
\begin{equation}
\label{eq:langevin-target}
dX_\tau = \nabla_x\log p_{\mathrm{data}}(X_\tau)\,d\tau + \sqrt2\,dW_\tau
\end{equation}
has $p_{\mathrm{data}}$ as its stationary distribution. Discretising
\eqref{eq:langevin-target} with step size $\epsilon$ gives the
\emph{unadjusted Langevin algorithm},
\begin{equation}
\label{eq:ula}
x^{(k+1)} = x^{(k)} + \epsilon\,\nabla_x\log p_{\mathrm{data}}\bigl(x^{(k)}\bigr) + \sqrt{2\epsilon}\,z^{(k)},
\qquad z^{(k)}\overset{\mathrm{iid}}\sim\mathcal N(0,I_d),
\end{equation}
which, run for enough steps $k$ from an arbitrary $x^{(0)}$, produces (an
approximate) sample from $p_{\mathrm{data}}$ using \emph{only} the score,
never the density itself. This is the precise sense in which ``the score
suffices for sampling'', and it is why score-based diffusion models train a
score network rather than a density network:
\Cref{sec:score-matching-intro} shows how to fit that network without ever
observing $p_{\mathrm{data}}$ in closed form.

\section{Kullback--Leibler divergence and maximum likelihood}
\label{sec:kl-mle}

We record two further standard tools that recur when we discuss what a
training objective is actually minimising.

\begin{definition}[Kullback--Leibler divergence]
\label{def:kl}
For densities $p,q$ on the same space with $p$ absolutely continuous with
respect to $q$,
\[
\KL(p\,\|\,q) = \int p(x)\log\frac{p(x)}{q(x)}\,dx = \EE_{x\sim p}\Bigl[\log\frac{p(x)}{q(x)}\Bigr] .
\]
\end{definition}

\begin{proposition}[Basic properties of KL]
\label{prop:kl-properties}
$\KL(p\,\|\,q)\ge0$, with equality iff $p=q$ ($q$-almost everywhere). In
general $\KL(p\,\|\,q)\ne\KL(q\,\|\,p)$.
\end{proposition}

\begin{proof}
By Jensen's inequality applied to the strictly convex function $-\log$,
$\KL(p\|q) = \EE_{p}[-\log(q(x)/p(x))] \ge -\log\EE_{p}[q(x)/p(x)] = -\log\int q(x)\,dx = -\log1 = 0$,
with equality iff $q(x)/p(x)$ is $p$-almost-surely constant, i.e.\ (since
both integrate to $1$) equal to $1$ almost everywhere. Asymmetry is
immediate from an example, e.g.\ $p,q$ both Gaussian with different
variances.
\end{proof}

\begin{definition}[Maximum likelihood]
\label{def:mle}
Given i.i.d.\ samples $x_1,\dots,x_n\sim p_{\mathrm{data}}$ and a
parametric family $\{p_\theta\}$, the \emph{maximum likelihood estimator} is
$\hat\theta = \argmax_\theta \sum_{i=1}^n\log p_\theta(x_i)$.
\end{definition}

\begin{proposition}[MLE minimises KL to the data distribution]
\label{prop:mle-kl}
As $n\to\infty$, $\tfrac1n\sum_i \log p_\theta(x_i) \to \EE_{x\sim p_{\mathrm{data}}}[\log p_\theta(x)]$
almost surely (law of large numbers), and
\[
\argmax_\theta \EE_{p_{\mathrm{data}}}[\log p_\theta(x)]
=
\argmin_\theta \KL\bigl(p_{\mathrm{data}}\,\|\,p_\theta\bigr) .
\]
\end{proposition}

\begin{proof}
$\KL(p_{\mathrm{data}}\|p_\theta) = \EE_{p_{\mathrm{data}}}[\log p_{\mathrm{data}}(x)] - \EE_{p_{\mathrm{data}}}[\log p_\theta(x)]$;
the first term does not depend on $\theta$, so minimising the KL divergence
over $\theta$ is the same optimisation as maximising
$\EE_{p_{\mathrm{data}}}[\log p_\theta(x)]$.
\end{proof}

\Cref{prop:mle-kl} is why ``fit the model by maximum likelihood'' and ``fit
the model to minimise KL divergence to the data'' are the same instruction,
a fact we use once, in \Cref{sec:likelihood-confidence}, when we discuss to
what extent a flow-matching model's ODE admits an exact likelihood and what
that would be useful for.

\section{Generative modelling as distribution transport}
\label{sec:generative-modelling-intro}

We can now state, precisely, the problem that both diffusion models and
flow matching solve, and why casting protein--ligand docking as an instance
of it is the right move.

\begin{quote}
\textbf{The generative modelling problem.} Given i.i.d.\ samples
$x_1,\dots,x_n\sim p_{\mathrm{data}}$, produce a mechanism that generates new
samples $\hat x\sim p_{\mathrm{data}}$ (or a good approximation thereof),
\emph{without} an explicit closed-form, evaluable expression for
$p_{\mathrm{data}}$ ever being available.
\end{quote}

The docking instance of this problem, following \Cref{chap:docking}, sets
$p_{\mathrm{data}}$ to be the (conditional, given the protein pocket and
ligand graph) distribution of experimentally observed bound poses --- a
distribution about which we have samples (crystal structures) but no
formula.

Rejection sampling and classical Markov chain Monte Carlo both require
pointwise evaluation of $p_{\mathrm{data}}$ up to a known constant, which we
do not have; \Cref{sec:langevin} showed a way around this for MCMC-style
methods (only the score is needed), but that still leaves the problem of
\emph{learning} the score from samples alone
(\Cref{sec:score-matching-intro}). The alternative pursued by both halves of
this document is:
\begin{enumerate}
\item Fix a \emph{tractable} distribution $p_{\mathrm{init}}$ (called the
\emph{source} or \emph{prior}) that we can sample from directly --- for
docking, e.g.\ an isotropic Gaussian for translations and the uniform
(Haar) distribution for rotations (both defined precisely in
\Cref{chap:groups}).
\item Learn a \emph{transport mechanism} --- a stochastic process
(\Cref{chap:diffusion}) or a deterministic flow (\Cref{chap:flow-matching})
--- that turns a sample from $p_{\mathrm{init}}$ into a sample from
$p_{\mathrm{data}}$.
\end{enumerate}
The entire technical content of \Cref{chap:diffusion,chap:flow-matching} is
two different, and ultimately closely related (\Cref{sec:pf-ode-scratch}), answers
to step 2. Both answers share one structural device, which is worth
isolating before either is developed in full: both construct a
\emph{probability path}, a time-indexed family of distributions
interpolating between the two endpoints.

\begin{definition}[Probability path]
\label{def:prob-path}
A \emph{probability path} is a family of densities $(p_t)_{t\in[0,1]}$ on
$\RR^d$ (or, later, on a manifold $\mathcal M$) satisfying the global
convention of \Cref{rem:time-convention-global}: $p_0 = p_{\mathrm{init}}$
and $p_1=p_{\mathrm{data}}$.
\end{definition}

Diffusion models (\Cref{chap:diffusion}) construct $(p_t)$ implicitly, as
the law of a stochastic process at time $t$, and characterise it via its
score, $\nabla_x\log p_t(x)$; sampling then simulates a stochastic
differential equation. Flow matching (\Cref{chap:flow-matching}) constructs
$(p_t)$ explicitly, via a chosen interpolation between a source point and a
data point, and characterises it via a vector field $u_t$ satisfying the
continuity equation of \Cref{thm:continuity}; sampling then integrates a
deterministic ordinary differential equation.
\Cref{sec:pf-ode-scratch} proves that these are not two disjoint frameworks but two
different parameterisations of exactly the same underlying object, a fact
that will let us describe the conversion from SigmaDock to SigmaFlow
(\Cref{chap:sigmaflow}) as a single, controlled substitution rather than a
change of model class.

\clearpage

% ================================================================
\part{Geometry}
\label{part:geometry}
% ================================================================

\chapter{Manifolds, Groups, and Rigid Motion}
\label{chap:groups}

\Cref{chap:docking} established that the natural coordinates of a ligand
pose are not $3n$ free Cartesian numbers but a translation, a rotation, and
(if fragments are used) a handful of torsion angles. None of translations,
rotations, and torsion angles are elements of a vector space in the way
that is needed for the calculus of \Cref{chap:prob-foundations} to apply
directly: a rotation is not closed under addition (the sum of two rotation
matrices is generally not a rotation matrix), and an angle is only defined
modulo $2\pi$. This chapter develops exactly the amount of differential
geometry needed to do calculus, probability, and eventually diffusion and
flow matching correctly on such spaces.

\section{Manifolds and tangent spaces}
\label{sec:manifolds-scratch}

\paragraph{Motivation.} Ordinary multivariable calculus is calculus on
$\RR^n$: derivatives are linear maps $\RR^n\to\RR^m$, gradients are vectors
in $\RR^n$, and one can always add two points or scale a point by a real
number. None of this is true, globally, on the set of rotation matrices or
on a circle of angles. What \emph{is} true is that these sets look like
$\RR^n$ \emph{locally} --- close enough together, two rotations differ by
``a small rotation'', and this small difference behaves, to first order,
exactly like a vector. Formalising ``looks like $\RR^n$ locally, glued
together globally'' is the definition of a manifold.

\paragraph{Intuition.} Picture the surface of a sphere $S^2\subset\RR^3$.
Around any point, a small enough patch of the sphere looks flat --- this is
why maps of small regions of the Earth can be drawn on flat paper with
negligible distortion, even though the Earth as a whole cannot. The
``flat approximation at a point'' is the tangent \emph{plane} at that point;
the entire curved sphere is not itself a vector space (you cannot add two
points on a sphere and expect the sum to be on the sphere), but each tangent
plane is.

\begin{definition}[Smooth manifold, informal]
\label{def:manifold-scratch}
An $n$-dimensional \emph{smooth manifold} $\mathcal M$ is a set that can be
covered by patches (\emph{charts}), each in smooth, invertible correspondence
with an open subset of $\RR^n$, such that on overlaps the two competing
correspondences are related by a smooth map. (We will only ever need two
concrete examples, $SO(3)$ and $\mathbb T=\RR/2\pi\ZZ$, both realised
concretely below as subsets of a Euclidean space; a fully general,
chart-based treatment is standard differential geometry and is not needed
beyond this informal level for anything in this document.)
\end{definition}

\begin{definition}[Tangent space, informal]
\label{def:tangent-scratch}
For $x\in\mathcal M$, the \emph{tangent space} $T_x\mathcal M$ is the vector
space of all velocity vectors $\dot\gamma(0)$ of smooth curves
$\gamma:(-\eps,\eps)\to\mathcal M$ with $\gamma(0)=x$. Concretely, if
$\mathcal M\subset\RR^N$ is realised as a level set (as both $SO(3)$ and
$\mathbb T$ will be), $T_x\mathcal M$ is the set of directions in $\RR^N$
tangent to $\mathcal M$ at $x$, and it \emph{is} a genuine linear subspace of
$\RR^N$: sums and scalar multiples of tangent vectors are again tangent
vectors, even though sums of points on $\mathcal M$ are generally not points
on $\mathcal M$.
\end{definition}

\begin{example}
Let $\mathcal M=S^1=\{x\in\RR^2:\|x\|=1\}$, the unit circle, and $x=(1,0)$.
A curve $\gamma(\tau)=(\cos\tau,\sin\tau)$ has $\gamma(0)=x$ and
$\dot\gamma(0)=(0,1)$. Every tangent curve at $x$ has velocity a multiple of
$(0,1)$ (the direction perpendicular to $x$), so
$T_xS^1=\{(0,c):c\in\RR\}\cong\RR$: a one-dimensional vector space, matching
$\dim S^1=1$, even though $S^1$ itself is not a vector space (e.g.\
$(1,0)+(1,0)=(2,0)\notin S^1$).
\end{example}

\begin{definition}[Riemannian metric, geodesic; restated]
\label{def:metric-scratch}
A \emph{Riemannian metric} on $\mathcal M$ is a smoothly varying family of
inner products $g_x:T_x\mathcal M\times T_x\mathcal M\to\RR$
($x\in\mathcal M$), inducing the norm $\|v\|_g=\sqrt{g_x(v,v)}$. The length
of a smooth curve $\gamma:[0,1]\to\mathcal M$ is
$L(\gamma)=\int_0^1\|\dot\gamma(\tau)\|_{g}\,d\tau$, and a \emph{geodesic} is
a curve that locally minimises this length, traversed at constant speed.
\end{definition}

\begin{definition}[Exponential and logarithmic map; restated]
\label{def:explog-scratch}
For $x\in\mathcal M$, the \emph{exponential map} $\exp_x:T_x\mathcal M\to\mathcal M$
sends $v$ to $\gamma_v(1)$, where $\gamma_v$ is the unique geodesic with
$\gamma_v(0)=x$, $\dot\gamma_v(0)=v$. Its local inverse, defined near $x$, is
the \emph{logarithmic map} $\log_x$.
\end{definition}

The pair $(\exp_x,\log_x)$ is, throughout this document, the mechanism that
replaces ordinary vector addition and subtraction: $\exp_x(v)$ plays the
role of $x+v$, and $\log_x(y)$ plays the role of $y-x$, valid even though
$\mathcal M$ has no addition of its own. Diffusion models
(\Cref{chap:diffusion}) will use $\exp_x$ to push a randomly sampled
tangent-space perturbation onto the manifold; flow matching
(\Cref{chap:flow-matching}) will use $\log_x$ to construct a
geodesic interpolant between two points and $\exp_x$ to integrate the
learned velocity field one step at a time.

\section{Groups and Lie groups}
\label{sec:groups-scratch}

\paragraph{Motivation.} A manifold alone tells us how to do calculus on a
curved space; it says nothing about how to \emph{compose} two elements
(``first rotate by $Q_1$, then by $Q_2$''). Rigid motions have exactly this
compositional structure --- applying two rotations in sequence is again a
rotation --- and capturing it precisely is the job of group theory.

\begin{definition}[Group]
\label{def:group}
A \emph{group} is a set $G$ with a binary operation
$\circ:G\times G\to G$ such that: (i) \emph{associativity},
$(g_1\circ g_2)\circ g_3=g_1\circ(g_2\circ g_3)$; (ii) there is an
\emph{identity} $e\in G$ with $e\circ g=g\circ e=g$ for all $g$; (iii) every
$g$ has an \emph{inverse} $g^{-1}$ with $g\circ g^{-1}=g^{-1}\circ g=e$.
\end{definition}

\begin{example}
$(\RR^d,+)$ is a group (identity $0$, inverse $-x$); it is commutative
($g_1\circ g_2=g_2\circ g_1$), i.e.\ \emph{abelian}. The set of invertible
$n\times n$ matrices under matrix multiplication, $GL(n,\RR)$, is a
(non-abelian, for $n\ge2$) group.
\end{example}

\begin{definition}[Lie group; restated]
\label{def:liegroup-scratch}
A \emph{Lie group} is a group $G$ that is simultaneously a smooth manifold,
such that the group operations $\circ:G\times G\to G$ and
$(\cdot)^{-1}:G\to G$ are smooth maps.
\end{definition}

The point of requiring smoothness of the group operation is that it lets us
transport tangent vectors from one point of the group to another using group
multiplication itself --- a device with no analogue on a generic manifold,
and the reason Lie groups admit a canonical, group-wide notion of ``the''
tangent space, rather than merely a separate tangent space at every point.

\begin{definition}[Lie algebra; restated]
\label{def:liealgebra-scratch}
Let $G$ be a Lie group with identity $e$. Its \emph{Lie algebra} is
$\mathfrak g := T_eG$, the tangent space at the identity.
\end{definition}

\begin{definition}[Left translation and left-invariant vector fields]
\label{def:left-translation}
For $g\in G$, \emph{left translation by $g$} is the smooth map
$L_g:G\to G$, $L_g(h)=g\circ h$. Its differential at the identity,
$dL_g|_e:\mathfrak g=T_eG\to T_gG$, is a linear isomorphism (it has smooth
inverse $dL_{g^{-1}}|_g$), so every $\xi\in\mathfrak g$ determines a tangent
vector $dL_g|_e(\xi)\in T_gG$ at \emph{every} point $g$ simultaneously.
\end{definition}

\Cref{def:left-translation} is the mechanism, promised above, by which a
single vector space $\mathfrak g$ suffices to describe tangent vectors
everywhere on $G$: any $\xi\in\mathfrak g$ can be moved to any point $g$ by
$v = dL_g|_e(\xi) \in T_gG$. This is called \emph{left trivialisation}, and
it is not the only choice --- \emph{right} translation
$R_g(h)=h\circ g$ gives an equally valid but generally \emph{different}
identification, called \emph{right trivialisation}. Which one a given piece
of code or a given formula uses is a genuine convention, not a matter of
taste, because $dL_g|_e(\xi)$ and $dR_g|_e(\xi)$ are different tangent
vectors at $g$ whenever $G$ is non-abelian; \Cref{sec:so3-trivialisation}
below makes this concrete for $SO(3)$, and \Cref{chap:sigmaflow} traces the
convention actually used by the SigmaFlow implementation with care, because
--- as flagged already in the introduction to this document --- the two
codebases are not visibly consistent with each other or with the existing
thesis draft on this exact point, and getting it wrong silently flips signs
in the trained model without necessarily crashing anything.

\section{The rotation group $SO(3)$}
\label{sec:so3-scratch}

\subsection{Definition and dimension}

\begin{definition}[Orthogonal and special orthogonal group]
\label{def:so3-def}
\[
O(3) = \{R\in\RR^{3\times3} : R^\top R = I\},
\qquad
SO(3) = \{R\in O(3) : \det R = 1\} .
\]
\end{definition}

Every $R\in O(3)$ satisfies $\det(R^\top R)=\det(R)^2=\det(I)=1$, so
$\det R=\pm1$; $SO(3)$ is the subset with $\det R=+1$, and these are exactly
the \emph{orientation-preserving} orthogonal maps --- rotations, as opposed
to reflections ($\det R=-1$). $O(3)$ has two disjoint pieces (\emph{connected
components}), separated by the sign of the determinant, and $SO(3)$ is
precisely the piece containing the identity, i.e.\ the piece reachable from
$I$ by a continuous path within $O(3)$. Physically: a continuous, physical
rotation of a rigid body can never turn it into its mirror image, which is
why $SO(3)$, not $O(3)$, is the correct group for orientations of physical
objects, and why the distinction resurfaces, non-optionally, when we discuss
chirality of molecules in \Cref{rem:parity-scratch}.

\begin{proposition}[Rotations preserve length and orientation]
\label{prop:so3-isometry}
For $R\in SO(3)$ and $x,y\in\RR^3$: $\|Rx\|=\|x\|$, $\langle Rx,Ry\rangle=\langle x,y\rangle$,
and $(Rx)\times(Ry)=R(x\times y)$.
\end{proposition}

\begin{proof}
$\langle Rx,Ry\rangle = x^\top R^\top R y = x^\top y=\langle x,y\rangle$
using $R^\top R=I$; taking $y=x$ gives $\|Rx\|^2=\|x\|^2$. For the cross
product identity, use the general fact $(Ax)\times(Ay)=\det(A)\,A^{-\top}(x\times y)$
for invertible $A$ (a standard identity, provable by checking it on the
standard basis and using multilinearity/antisymmetry of both sides), and
specialise to $A=R$: $R^{-\top}=R$ (since $R^\top=R^{-1}$) and $\det R=1$.
\end{proof}

\begin{proposition}[Dimension of $SO(3)$]
\label{prop:so3-dim}
$SO(3)$ is a $3$-dimensional manifold.
\end{proposition}

\begin{proof}
The constraint $R^\top R=I$ is a system of equations in the $9$ entries of
$R$; the matrix $R^\top R$ is automatically symmetric, so this constraint
has only $\binom{3+1}{2}=6$ independent scalar components (the entries on
and above the diagonal of a $3\times3$ symmetric matrix). One checks (via
the implicit function theorem, using that the differential of
$R\mapsto R^\top R$ has full rank $6$ at every $R\in O(3)$) that this cuts
out a smooth $9-6=3$-dimensional manifold, of which $SO(3)$ is one
(open-and-closed, hence a full-dimensional) connected component; the extra
condition $\det R=1$ removes no further dimension, it only selects a
component. We take the rank computation as a standard fact from
differential topology and do not reprove it in full generality.
\end{proof}

Three degrees of freedom is exactly the number we expect: an orientation in
$3$-dimensional space can be described by, for instance, two angles fixing
an axis on the sphere $S^2$ (a $2$-dimensional choice) and one further angle
of rotation about that axis, and this count recurs throughout the rest of
this chapter and \Cref{chap:groups} onward.

\subsection{The Lie algebra $\so(3)$ and the hat/vee maps}
\label{sec:so3-liealgebra}

\begin{proposition}[$\so(3)$ is the skew-symmetric matrices]
\label{prop:so3-liealgebra}
$T_ISO(3) = \{\Omega\in\RR^{3\times3} : \Omega^\top=-\Omega\}=:\so(3)$.
\end{proposition}

\begin{proof}
Let $\gamma(\tau)$ be a smooth curve in $SO(3)$ with $\gamma(0)=I$. Since
$\gamma(\tau)^\top\gamma(\tau)=I$ for all $\tau$, differentiating at $\tau=0$
using the product rule gives
$\dot\gamma(0)^\top I + I^\top\dot\gamma(0) = 0$, i.e.\
$\dot\gamma(0)^\top=-\dot\gamma(0)$: every tangent vector at $I$ is
skew-symmetric. Conversely, the space of skew-symmetric $3\times3$ matrices
is $3$-dimensional (three free entries, above the diagonal, with the rest
determined by antisymmetry and a zero diagonal), matching
\Cref{prop:so3-dim}, so this necessary condition is also sufficient (both
sides are $3$-dimensional subspaces, one containing the other, hence equal).
\end{proof}

\begin{definition}[Hat and vee operators; restated precisely]
\label{def:hatvee-scratch}
\[
(\cdot)^{\wedge}:\RR^3\to\so(3),
\qquad
\widehat\omega
=
\begin{pmatrix}
0&-\omega_3&\omega_2\\
\omega_3&0&-\omega_1\\
-\omega_2&\omega_1&0
\end{pmatrix},
\qquad
\omega=(\omega_1,\omega_2,\omega_3)^\top,
\]
and $(\cdot)^{\vee}:\so(3)\to\RR^3$ is its (linear, bijective) inverse.
\end{definition}

\begin{proposition}[Hat map intertwines cross product and matrix commutator]
\label{prop:hat-cross}
For $\omega,v\in\RR^3$: $\widehat\omega\, v = \omega\times v$. Consequently,
for $\omega_1,\omega_2\in\RR^3$, the matrix commutator satisfies
$[\widehat{\omega_1},\widehat{\omega_2}] := \widehat{\omega_1}\widehat{\omega_2}-\widehat{\omega_2}\widehat{\omega_1}
= \widehat{\omega_1\times\omega_2}$.
\end{proposition}

\begin{proof}
Direct computation: $\widehat\omega v =
(\omega_2v_3-\omega_3v_2,\ \omega_3v_1-\omega_1v_3,\ \omega_1v_2-\omega_2v_1)^\top
= \omega\times v$ by the definition of the cross product. The
commutator identity follows by applying both sides to an arbitrary $v$ and
using the bilinearity and antisymmetry of $\times$ together with the
vector identity $a\times(b\times v) - b\times(a\times v) = (a\times b)\times v$
(itself checkable directly, or via the BAC-CAB rule
$a\times(b\times c)=b(a\cdot c)-c(a\cdot b)$).
\end{proof}

\begin{codebox}
\textbf{Mathematical object:} $(\cdot)^{\wedge}:\RR^3\to\so(3)$, $(\cdot)^{\vee}:\so(3)\to\RR^3$ \\
\textbf{Implementation:} \file{so3\_utils.py}::\pyfunc{hat}, \pyfunc{vee}
(identical in \file{SigmaDock/} and \file{SigmaFlow\_Development/}) \\
\textbf{Tensor shape:} \shape{..., 3} $\leftrightarrow$ \shape{..., 3, 3} \\
\textbf{Interpretation:} \pyfunc{hat} builds the skew-symmetric matrix by
constructing only the strict upper triangle and then anti-symmetrising
(\texttt{hat\_v + -hat\_v.transpose(-1,-2)}); \pyfunc{vee} inverts it by
reading off $(-A_{23},A_{13},-A_{12})$ and only \emph{prints} a warning
(rather than raising) if the input is not skew-symmetric within a loose
tolerance --- this is not a hard safety check, and a caller that
accidentally passes a non-skew matrix will not be stopped by \pyfunc{vee}
itself.
\end{codebox}

\subsection{The exponential map: Rodrigues' formula}
\label{sec:rodrigues-scratch}

\paragraph{Motivation.} We now want an explicit, closed-form recipe for
$\exp_I:\so(3)\to SO(3)$ (\Cref{def:explog-scratch} specialised to $G=SO(3)$,
$x=I$): given a small rotation $\widehat\omega\in\so(3)$, what is the actual
rotation matrix obtained by ``following the corresponding geodesic for unit
time''? On a Lie group, this question has a beautifully uniform answer.

\begin{proposition}[Geodesics through the identity of a bi-invariant Lie group are one-parameter subgroups]
\label{prop:geodesic-oneparam}
Equip $SO(3)$ with the metric $g_I(\Omega_1,\Omega_2) = \tfrac12\tr(\Omega_1^\top\Omega_2)$
at the identity, extended to every other point by left translation
(\Cref{def:left-translation}); one checks this metric is also invariant
under \emph{right} translation (\emph{bi-invariant}), which is a special,
convenient feature of compact groups such as $SO(3)$ and is what makes the
following true. The geodesic through $I$ with initial velocity
$\widehat\omega\in\so(3)$ is $\gamma(\tau) = \exp(\tau\widehat\omega)$, the
matrix exponential.
\end{proposition}

We do not prove \Cref{prop:geodesic-oneparam} in full generality (it is a
standard fact about bi-invariant metrics on Lie groups, following from the
observation that one-parameter subgroups $\tau\mapsto\exp(\tau\Omega)$ of a
bi-invariant Lie group have constant speed and locally minimise length; see,
e.g., do Carmo, \emph{Riemannian Geometry}, 1992, Ch.~11); we use it only to
justify that $\exp_I=\exp$, the ordinary matrix exponential
$\exp(A)=\sum_{k=0}^\infty A^k/k!$, restricted to $\so(3)$. What remains is
to evaluate this infinite series in closed form.

\begin{theorem}[Rodrigues' rotation formula]
\label{thm:rodrigues-scratch}
For $\omega\in\RR^3$ with $\theta=\|\omega\|$,
\[
\exp(\widehat\omega)
=
I + \frac{\sin\theta}{\theta}\,\widehat\omega + \frac{1-\cos\theta}{\theta^2}\,\widehat\omega^{\,2}
\qquad(\theta\ne0),
\]
and $\exp(\widehat\omega)=I$ when $\theta=0$ (continuously extending the
formula, since both fractions have removable singularities at $\theta=0$).
The resulting matrix is the rotation by angle $\theta$ about the axis
$\omega/\theta$, in the right-hand sense.
\end{theorem}

\begin{proof}
Write $n=\omega/\theta$ (unit vector) so $\widehat\omega=\theta\widehat n$.
The key algebraic fact is the following closed recursion for powers of
$\widehat n$, provable directly from \Cref{prop:hat-cross} (with $n$ a unit
vector) via the BAC-CAB identity:
\[
\widehat n^{\,2} = nn^\top - I,
\qquad
\widehat n^{\,3} = \widehat n\,\widehat n^{\,2} = \widehat n(nn^\top-I) = -\widehat n,
\]
using $\widehat n\,n=n\times n=0$ (so $\widehat n\,nn^\top=0$). Hence
$\widehat n^{\,3}=-\widehat n$, so every higher power reduces:
$\widehat n^{4}=-\widehat n^{2}$, $\widehat n^{5}=\widehat n$, etc., with
period $4$ up to sign, i.e.\ odd powers are $\pm\widehat n$ and even powers
(beyond $0$) are $\pm\widehat n^2$. Substituting into the exponential series
and grouping by parity,
\begin{align*}
\exp(\theta\widehat n)
&= I + \Bigl(\theta - \frac{\theta^3}{3!}+\frac{\theta^5}{5!}-\cdots\Bigr)\widehat n
+ \Bigl(\frac{\theta^2}{2!}-\frac{\theta^4}{4!}+\frac{\theta^6}{6!}-\cdots\Bigr)\widehat n^{\,2}\\
&= I + \sin(\theta)\,\widehat n + \bigl(1-\cos(\theta)\bigr)\,\widehat n^{\,2} ,
\end{align*}
recognising the Taylor series of $\sin$ and $1-\cos$. Substituting back
$\widehat n=\widehat\omega/\theta$ gives the claimed formula. That the
result is a rotation by $\theta$ about axis $n$ can be checked directly:
$\exp(\theta\widehat n)\,n = n$ (since $\widehat n n=0$, so both correction
terms vanish on $n$), so $n$ is fixed, and restricted to the plane
orthogonal to $n$ the formula reduces (by direct computation using
$\widehat n^2 v=-v$ for $v\perp n$) to the standard $2\times2$ rotation-by-$\theta$
matrix in an orthonormal basis of that plane.
\end{proof}

\begin{codebox}
\textbf{Mathematical object:} $\exp:\so(3)\to SO(3)$, $\Exp:\RR^3\to SO(3)$ \\
\textbf{Implementation:} \file{so3\_utils.py}::\pyfunc{exp}, \pyfunc{Exp} \\
\textbf{Tensor shape:} \shape{..., 3, 3} $\to$ \shape{..., 3, 3} (\pyfunc{exp}); \shape{..., 3} $\to$ \shape{..., 3, 3} (\pyfunc{Exp}, defined as \texttt{exp(hat(v))}) \\
\textbf{Interpretation:} \textbf{Neither codebase uses the closed-form
Rodrigues formula of \Cref{thm:rodrigues-scratch} directly.} \pyfunc{exp} is
implemented as a literal call to \texttt{torch.matrix\_exp}, PyTorch's
generic (Pad\'e-approximation-based) matrix exponential, with a CPU
round-trip fallback on Apple-Silicon (MPS) backends where
\texttt{torch.matrix\_exp} is unsupported. This is mathematically equivalent
to \Cref{thm:rodrigues-scratch} (both compute the same matrix exponential;
$\so(3)$'s exponential series converges absolutely for every finite
$\theta$, so there is no issue of principle), but it is worth recording
explicitly as a discrepancy between how this document derives the map and
how the shipped code evaluates it: \Cref{thm:rodrigues-scratch}'s closed
form is what makes the derivations of \Cref{sec:so3-log-scratch} and
\Cref{chap:diffusion,chap:flow-matching} tractable by hand, but the
executed program never evaluates $\sin\theta/\theta$ directly at this step.
\end{codebox}

\subsection{The logarithmic map, and its numerically dangerous points}
\label{sec:so3-log-scratch}

\paragraph{Motivation.} $\log:SO(3)\to\so(3)$ (near $I$; more on the
domain of validity below) is the local inverse of $\exp$, and it is what
lets us measure ``how different are two rotations'' as a single tangent
vector, exactly the operation flow matching needs to build an interpolant
between a source rotation and a target rotation
(\Cref{sec:fm-so3-scratch}).

\begin{theorem}[Logarithm of a rotation matrix]
\label{thm:so3-log-scratch}
For $R\in SO(3)$ with rotation angle $\theta=\arccos\bigl(\tfrac{\tr(R)-1}2\bigr)\in[0,\pi]$
and $\theta\notin\{0,\pi\}$,
\[
\log(R) = \frac{\theta}{2\sin\theta}\,(R-R^\top)\ \in\ \so(3) .
\]
\end{theorem}

\begin{proof}
By \Cref{thm:rodrigues-scratch}, $R=I+\sin\theta\,\widehat n+(1-\cos\theta)\widehat n^{\,2}$
for the (unknown) axis $n$ and angle $\theta$ of $R$. Since $\widehat n$ is
skew-symmetric and $\widehat n^{\,2}=nn^\top-I$ is symmetric, the
skew-symmetric part of $R$ is exactly
$\tfrac12(R-R^\top)=\sin\theta\,\widehat n$. Solving for $\widehat n$ (valid
whenever $\sin\theta\ne0$, i.e.\ $\theta\notin\{0,\pi\}$) gives
$\widehat n = (R-R^\top)/(2\sin\theta)$, and hence
$\log(R)=\theta\widehat n = \dfrac{\theta}{2\sin\theta}(R-R^\top)$. The
formula for $\theta$ itself follows by taking the trace of Rodrigues'
formula: $\tr(R)=3+0+(1-\cos\theta)\tr(\widehat n^{\,2})=3+(1-\cos\theta)(-2)=1+2\cos\theta$,
using $\tr(\widehat n^{\,2})=\tr(nn^\top-I)=\|n\|^2-3=-2$; solving gives
$\cos\theta=(\tr(R)-1)/2$.
\end{proof}

The two excluded points $\theta=0$ and $\theta=\pi$ are not edge cases that
can be ignored: $\theta=0$ (the identity rotation) is reached with
probability zero under any continuous sampling scheme but is approached
arbitrarily closely whenever two rotations being compared are nearly equal,
which happens \emph{routinely} near convergence of an iterative sampler; and
$\theta=\pi$ (a half-turn) has positive probability density under the
uniform distribution on $SO(3)$ derived in \Cref{prop:haar-so3} below, so it
\emph{will} be sampled during training and inference at a non-negligible
rate. Both require separate formulas.

\begin{proposition}[Small-angle behaviour of $\log$]
\label{prop:log-smallangle}
As $\theta\to0$,
\[
\log(R) = \frac12\Bigl(1+\frac{\theta^2}{6}+O(\theta^4)\Bigr)(R-R^\top) .
\]
\end{proposition}

\begin{proof}
Taylor-expand $\theta/(2\sin\theta) = \tfrac12\bigl(1-\theta^2/6+O(\theta^4)\bigr)^{-1}
= \tfrac12\bigl(1+\theta^2/6+O(\theta^4)\bigr)$, using $\sin\theta=\theta-\theta^3/6+O(\theta^5)$.
\end{proof}

\begin{proposition}[Behaviour of $\log$ near $\theta=\pi$]
\label{prop:log-nearpi}
As $\theta\to\pi$, $R\to I+2nn^\top-2I$ for the rotation axis $n$ (up to
sign), so $\tfrac12(R+I)\to nn^\top$; consequently $|n_i|$ can be recovered,
up to the sign ambiguity $n\leftrightarrow-n$ (which does not affect
$\widehat n$, since $\widehat n=\widehat{-n}$ composed with $\theta\to-\theta$
represents the same rotation and the formula below is used together with
$\theta=\pi$), as $|n_i|=\sqrt{\bigl(\tfrac12(R+I)\bigr)_{ii}}$, with signs
fixed by comparing off-diagonal entries of $\tfrac12(R+I)=nn^\top$ against
$n_in_j$.
\end{proposition}

\begin{proof}
At $\theta=\pi$: $\sin\theta=0,\ \cos\theta=-1$, so Rodrigues' formula gives
$R = I + 0 + 2\widehat n^{\,2} = I + 2(nn^\top-I) = 2nn^\top - I$, hence
$\tfrac12(R+I)=nn^\top$, a rank-one symmetric matrix with unit-norm columns
along $\pm n$. Reading off $|n_i|$ from the diagonal and signs from the
off-diagonal entries is then direct linear algebra on the outer product
$nn^\top$.
\end{proof}

\begin{codebox}
\textbf{Mathematical object:} $\log:SO(3)\to\so(3)$ near a generic rotation;
its numerically safe extension to $\theta\to0$ and $\theta\to\pi$ \\
\textbf{Implementation:} \file{so3\_utils.py}::\pyfunc{rotation\_vector\_from\_matrix}
(called by \pyfunc{Log}, whose output is then \pyfunc{hat}-lifted by
\pyfunc{log}) \\
\textbf{Tensor shape:} \shape{..., 3, 3} $\to$ \shape{..., 3} (vee-map
convention: the vector form $v=\theta n$, not the $\so(3)$ matrix) \\
\textbf{Interpretation:} the code computes \emph{all three} regimes ---
generic (\Cref{thm:so3-log-scratch}), near-zero
(\Cref{prop:log-smallangle}, using the exact Taylor truncation derived
above), and near-$\pi$ (\Cref{prop:log-nearpi}, with sign recovered from
whichever row of $\tfrac12(R+I)$ has the largest norm) --- for \emph{every}
input, and blends them with continuous $0/1$ masks built from
\pyfunc{torch.isclose} (tolerance $\mathrm{atol}=10^{-8}$ near $\theta=0$,
the much looser $\mathrm{atol}=10^{-2}$ near $\theta=\pi$), rather than a
hard \texttt{if}/\texttt{else} branch. This is a deliberate autodiff-friendly
pattern: computing all three branches unconditionally avoids
\texttt{NaN}-producing gradients from a branch that was not selected but
would otherwise still be backpropagated through, at the cost of some
redundant computation.
\end{codebox}

\begin{caution}[The angle-clamp bug, and why it mattered only for flow matching]
\label{caution:omega-clamp}
The function that computes $\theta=\arccos\bigl(\tfrac{\tr(R)-1}2\bigr)$
(\pyfunc{Omega} in \file{so3\_utils.py}) must clamp its argument to
$[-1,1]$ before calling \pyfunc{arccos}, since floating-point round-off can
push $\tfrac{\tr(R)-1}2$ fractionally outside this interval even for a
valid $R\in SO(3)$. The original SigmaDock code clamps to
$[-0.99,0.99]$, i.e.\ $\theta\in[\arccos(0.99),\arccos(-0.99)]\approx[8.1^\circ,172^\circ]$
--- \emph{any true angle beyond $172^\circ$ is silently rounded down to
$172^\circ$}. For diffusion this is close to harmless: $\theta$ enters only
through the smooth score function of \Cref{sec:diff-so3}, degrading
gracefully under a slightly wrong clamp. For flow matching it is not
harmless: the geodesic interpolant of \Cref{sec:fm-so3-scratch} needs
$\log$ to be an \emph{exact} inverse of $\exp$ (an exact round trip
$\exp(\log(R))=R$), which the $172^\circ$ clamp breaks for roughly the
$\bigl(180-172\bigr)/180\approx4.4\%$ of uniformly sampled rotation pairs
whose relative angle exceeds it (in each direction, so of order $9\%$ combined
--- matching the figure reported in this project's own development history).
SigmaFlow's \file{so3\_utils.py} narrows the clamp to
$[-1+10^{-7},\,1-10^{-7}]$, i.e.\ $\theta\in[0.026^\circ,179.97^\circ]$,
purely for \pyfunc{arccos}'s numerical domain, and documents the reason for
the change directly in the source (\file{so3\_flow\_matcher.py}, inline
comment). This document treats the fix as settled and uses the corrected
clamp implicitly wherever $\theta$ is discussed.
\end{caution}

\subsection{$SO(3)$ as a metric space, and Brownian motion's invariant measure}
\label{sec:so3-uniform}

We record two facts about $SO(3)$ that are used repeatedly from
\Cref{chap:diffusion} onward and are best derived once, here, in isolation
from any generative-modelling context.

\begin{proposition}[Bi-invariant Riemannian distance on $SO(3)$]
\label{prop:so3-distance}
With the metric of \Cref{prop:geodesic-oneparam}, the Riemannian distance
between $R_1,R_2\in SO(3)$ is $d(R_1,R_2) = \|\log(R_1^\top R_2)\|_g = \theta$,
the rotation angle of the relative rotation $R_1^\top R_2$.
\end{proposition}

\begin{proof}[Proof sketch]
By bi-invariance, $d(R_1,R_2)=d(I,R_1^\top R_2)$ (left-translate by
$R_1^{-1}=R_1^\top$), and $d(I,R)$ is the length of the geodesic from $I$ to
$R$, which by \Cref{prop:geodesic-oneparam} is $\exp(\tau\widehat\omega)$
for the $\omega$ with $\exp(\widehat\omega)=R$, $\|\omega\|$ minimal --- this
is exactly $\theta=\|\log(R)^\vee\|$ under the metric normalisation of
\Cref{prop:geodesic-oneparam} (chosen precisely so that the identification
$\so(3)\cong\RR^3$ of \Cref{def:hatvee-scratch} is a metric isometry, i.e.\
so that $\|\widehat\omega\|_g=\|\omega\|_2$ --- a direct check:
$g_I(\widehat\omega,\widehat\omega)=\tfrac12\tr(\widehat\omega^\top\widehat\omega)=\tfrac12\tr(-\widehat\omega^2)=\|\omega\|_2^2$
by a short computation using \Cref{prop:hat-cross}).
\end{proof}

\begin{proposition}[The Haar measure on $SO(3)$, in axis--angle coordinates]
\label{prop:haar-so3}
The measure on $SO(3)$ invariant under both left and right multiplication
(the \emph{Haar measure}, unique up to normalisation on a compact group;
existence/uniqueness is a classical fact we do not reprove --- see, e.g.,
Folland, \emph{A Course in Abstract Harmonic Analysis}, 1995, Ch.~2), when
$SO(3)$ is parameterised by axis $n\in S^2$ and angle $\theta\in[0,\pi]$ via
$R=\exp(\theta\widehat n)$, has density, with respect to
$d\theta\times(\text{uniform measure on }S^2)$, proportional to
$1-\cos\theta$.
\end{proposition}

\begin{proof}[Proof sketch]
This is the statement that the Riemannian volume form of
\Cref{prop:geodesic-oneparam}, expressed in axis--angle (geodesic normal)
coordinates, has Jacobian $\propto(1-\cos\theta)$ relative to the flat
product measure $d\theta\times d\Omega_{S^2}$; the computation is a standard
but slightly involved exercise in Riemannian normal coordinates on a compact
Lie group (the Jacobian of $\exp_I$ at $\theta\widehat n$, for a bi-invariant
metric, is $\prod_k \bigl(\tfrac{2\sin(\lambda_k\theta/2)}{\lambda_k\theta}\bigr)$
over the eigenvalues $\lambda_k$ of $\mathrm{ad}_{\widehat n}$; for $\so(3)$
these are $0,\pm i$, and evaluating the product reduces, after simplifying
the complex-conjugate pair, to $(1-\cos\theta)/(\theta^2/2)\cdot(\text{const})$
times the flat measure), which we do not carry out in full here and instead
take as a standard fact (see, e.g., Chirikjian, \emph{Stochastic Models,
Information Theory, and Lie Groups}, Vol.\,2, 2012, \S10.2). We record the
consequence that matters practically: to sample $R\sim\mathrm{Haar}(SO(3))$,
sample $n$ uniformly on $S^2$ and, independently, sample $\theta$ from the
density $\propto1-\cos\theta$ on $[0,\pi]$.
\end{proof}

\begin{codebox}
\textbf{Mathematical object:} $R\sim\mathrm{Haar}(SO(3))$ \\
\textbf{Implementation:} \file{so3\_utils.py}::\pyfunc{sample\_uniform} \\
\textbf{Tensor shape:} returns \shape{N, 3, 3} \\
\textbf{Interpretation:} builds a $1000$-point grid on $[0,\pi]$, numerically
integrates the density $\propto1-\cos\theta$ of \Cref{prop:haar-so3} into a
CDF, draws $\theta$ by inverse-CDF sampling via linear interpolation on
this grid (\pyfunc{np.interp}), draws an axis by normalising a standard
Gaussian vector in $\RR^3$ (a standard, exact way to sample uniformly on
$S^2$: the direction of an isotropic Gaussian is rotation-invariant in
distribution, since the Gaussian density $\propto e^{-\|x\|^2/2}$ depends
on $x$ only through $\|x\|$), and returns $\exp(\theta\widehat n)$. This
sampling routine runs on \texttt{numpy}'s global random state, \emph{not}
through \pyfunc{torch}'s seeded generator --- a genuine reproducibility
subtlety: seeding \pyfunc{torch.manual\_seed} alone does not fix the axes
or angles drawn by this function.
\end{codebox}

\section{Left and right trivialisation, resolved for this document}
\label{sec:so3-trivialisation}

We can now make \Cref{def:left-translation}'s warning fully concrete for
$SO(3)$, since this is a point at which, as flagged in the introduction to
this document, the existing thesis draft and the SigmaFlow source code do
not visibly agree, and it must be resolved by direct derivation rather than
by trusting either source.

\begin{definition}[Left- and right-trivialised angular velocity]
\label{def:trivialisation-scratch}
Let $R(\tau)$ be a smooth curve in $SO(3)$. Its \emph{left-trivialised} (or
\emph{world-frame}, or \emph{spatial}) angular velocity is
$\omega^{\mathrm L}(\tau)\in\RR^3$ with
$\dot R(\tau) = \widehat{\omega^{\mathrm L}(\tau)}\,R(\tau)$; its
\emph{right-trivialised} (or \emph{body-frame}) angular velocity is
$\omega^{\mathrm R}(\tau)\in\RR^3$ with
$\dot R(\tau) = R(\tau)\,\widehat{\omega^{\mathrm R}(\tau)}$.
\end{definition}

\begin{remark}[Terminology warning]
\label{rem:trivialisation-terminology}
Unfortunately, the words ``left'' and ``right'' in this context are used
inconsistently across the literature, because they can refer either to
\emph{which side of $R$ the hat matrix is multiplied on} (the convention
fixed in \Cref{def:trivialisation-scratch} above, which we adopt, since it
matches the standard Lie-theory usage in which $\omega^{\mathrm L}$ is
obtained by \emph{left}-translating a Lie-algebra element to $T_RG$, i.e.\
$dL_R|_e$) or to \emph{which of $\dot R R^{-1}$ vs.\ $R^{-1}\dot R$} is
meant, which some authors instead label by the opposite name. We fix
\Cref{def:trivialisation-scratch} once and use it consistently for the
remainder of this document; every other source cited or discussed is
translated into this convention explicitly before its formulas are used.
\end{remark}

By definition, $\widehat{\omega^{\mathrm L}}=\dot RR^\top=\dot RR^{-1}$ and
$\widehat{\omega^{\mathrm R}}=R^\top\dot R=R^{-1}\dot R$, and they are
related by $\omega^{\mathrm L}=R\omega^{\mathrm R}$ (apply
\Cref{prop:hat-cross}-style conjugation: $R\widehat{\omega^{\mathrm R}}R^\top=\widehat{R\omega^{\mathrm R}}$,
a direct consequence of \Cref{prop:so3-isometry}'s cross-product identity),
i.e.\ the two differ by \emph{rotating the vector itself by $R$}. They are
equal only when $R$ commutes with the relevant tangent direction, generally
not the case.

\begin{proposition}[The geodesic interpolant is naturally right-trivialised]
\label{prop:geodesic-right-triv}
Let $R(\tau)=R_0\exp(\tau\Xi)$ for fixed $R_0\in SO(3)$, $\Xi\in\so(3)$
(this is exactly the geodesic construction of \Cref{def:so3-interp-detail}
below). Then $\dot R(\tau) = R(\tau)\,\Xi$, i.e.\ the curve has \emph{constant}
right-trivialised (body-frame) angular velocity $\omega^{\mathrm R}=\Xi^\vee$,
in the sense of \Cref{def:trivialisation-scratch}.
\end{proposition}

\begin{proof}
Differentiate $R(\tau)=R_0\exp(\tau\Xi)$ directly:
$\dot R(\tau) = R_0\exp(\tau\Xi)\,\Xi = R(\tau)\,\Xi$, using that $\Xi$
commutes with $\exp(\tau\Xi)$ (both are polynomials/power series in $\Xi$).
This matches the second case of \Cref{def:trivialisation-scratch} exactly,
with $\omega^{\mathrm R}(\tau)\equiv\Xi^\vee$, constant in $\tau$.
\end{proof}

\begin{proposition}[The corresponding one-step update is a right multiplication]
\label{prop:euler-right-triv}
If a numerical scheme advances $R$ by one Euler step using a right-trivialised
velocity $v\in\RR^3$ over a time increment $dt$, the update that remains
exactly on $SO(3)$, consistent with \Cref{prop:geodesic-right-triv}, is
$R_{\mathrm{new}} = R\,\exp(dt\,\widehat v)$ (multiplying the exponential on
the \emph{right}); the analogous left-trivialised update would instead be
$R_{\mathrm{new}}=\exp(dt\,\widehat v)\,R$.
\end{proposition}

\Cref{prop:geodesic-right-triv,prop:euler-right-triv} settle the question
for this document: wherever SigmaFlow's geodesic interpolant
(\Cref{sec:fm-so3-scratch}) or its Euler integrator
(\Cref{sec:sigmaflow-sampling}) is discussed, the angular velocity supplied
by, and returned to, the network is the \textbf{right-trivialised (body-frame)}
angular velocity, consistent with the SigmaFlow source code's own comment
(\file{so3\_flow\_matcher.py}) and with the update rule
$R\mapsto R\exp(dt\,\widehat v)$ used by both \pyfunc{euler\_step} and
\pyfunc{expmap} in the actual implementation (both right-multiply the
exponential onto $R$, so --- once \Cref{def:trivialisation-scratch}'s
labelling convention is fixed and applied uniformly --- the two are in fact
mutually consistent, and the apparent conflict flagged by the code-research
pass resolves to a terminology mismatch between the source comment's word
choice and the existing thesis draft's word choice, not a mathematical
inconsistency in either). \Cref{chap:sigmadock} performs the same check,
independently, for SigmaDock's diffusion construction, since diffusion's
Brownian-motion increments (\Cref{sec:diff-so3}) are built from
\pyfunc{expmap}/\pyfunc{tangent\_gaussian}, the same right-multiplying
primitives, and the two chapters must therefore use the \emph{same}
trivialisation convention for their comparison in \Cref{chap:sigmaflow} to
be meaningful.

\begin{caution}
This resolution is a derivation carried out for this document; it should
not be taken to mean the original published SigmaDock description
necessarily uses matching language, only that the code that actually runs
is self-consistent under \Cref{def:trivialisation-scratch}. Where the
distinction between $\omega^{\mathrm L}$ and $\omega^{\mathrm R}$ affects a
formula materially (as opposed to being a bookkeeping choice with no
observable consequence, e.g.\ because a subsequent step immediately
re-expresses the result in world frame), this document states which one is
meant at the point of use, rather than relying on the reader to track it
globally.
\end{caution}

\section{The pose group $SE(3)$ and the product manifold $SE(3)^N$}
\label{sec:se3-scratch}

\begin{definition}[Special Euclidean group]
\label{def:se3-scratch}
$SE(3)$ is the group of rigid motions of $\RR^3$: as a set,
$SE(3) \cong SO(3)\times\RR^3$, with element $T=(R,x)$ acting on $y\in\RR^3$
by $T\cdot y = Ry+x$, and group operation
\[
(R_1,x_1)\circ(R_2,x_2) = (R_1R_2,\ R_1x_2+x_1)
\]
(composition of affine maps: apply $T_2$ then $T_1$).
\end{definition}

\begin{proposition}
$SE(3)$ is a $6$-dimensional Lie group, with identity $(I,0)$ and inverse
$(R,x)^{-1}=(R^\top,-R^\top x)$.
\end{proposition}

\begin{proof}
Immediate from \Cref{def:se3-scratch}: $\dim SE(3)=\dim SO(3)+\dim\RR^3=3+3=6$
as a manifold (a product of manifolds, \Cref{def:manifold-scratch}, has
dimension the sum of the factors' dimensions); associativity and the group
axioms follow from associativity of composition of affine maps; smoothness
of the group operations is immediate since \Cref{def:se3-scratch}'s formula
is polynomial in $(R_1,x_1,R_2,x_2)$.
\end{proof}

\paragraph{Metric choice.} $SE(3)$ admits more than one natural Riemannian
metric --- one may couple the rotational and translational blocks (e.g.\
scaling one relative to the other by a length parameter reflecting the
physical size of the body being moved) or keep them fully decoupled. We fix,
following Yim et al.\ (2023, cited precisely in \Cref{chap:sigmaflow} where
this choice is used), the \emph{product metric}
\begin{equation}
\label{eq:se3-product-metric-scratch}
g_{(R,x)}\bigl((\dot R_1,\dot x_1),(\dot R_2,\dot x_2)\bigr)
=
\langle\dot R_1,\dot R_2\rangle_{SO(3)} + \langle\dot x_1,\dot x_2\rangle_{\RR^3} .
\end{equation}

\begin{proposition}[$SE(3)$ with the product metric is a Riemannian product]
\label{prop:se3-product}
With \eqref{eq:se3-product-metric-scratch}, the exponential and logarithmic
maps, geodesics, and (once defined in \Cref{chap:diffusion}) Brownian
motion on $SE(3)$ all factorise componentwise:
\[
\exp_{(R,x)}(\dot R,\dot x) = \bigl(\exp_R(\dot R),\ x+\dot x\bigr),
\qquad
\log_{(R,x)}(R',x') = \bigl(\log_R(R'),\ x'-x\bigr) .
\]
\end{proposition}

\begin{proof}[Proof sketch]
A Riemannian product $(\mathcal M_1\times\mathcal M_2, g_1\oplus g_2)$ has
geodesics exactly the curves $(\gamma_1,\gamma_2)$ with $\gamma_i$ a
geodesic of $(\mathcal M_i,g_i)$ (since the geodesic equation, derived from
the metric via the Euler--Lagrange equations for the length functional,
decouples exactly when the metric does, as \eqref{eq:se3-product-metric-scratch}
does by construction), and $\RR^3$ with the flat Euclidean metric has
$\exp_x(v)=x+v$, $\log_x(y)=y-x$ trivially. Combining these two facts with
$SO(3)$'s own maps (\Cref{sec:rodrigues-scratch,sec:so3-log-scratch}) gives
the claim.
\end{proof}

\Cref{prop:se3-product} is, computationally, the single most consequential
fact in this chapter: it licenses treating translation and rotation
\emph{completely independently} in every later construction (the
conditional probability paths of \Cref{chap:diffusion,chap:flow-matching},
the network's two output heads of \Cref{chap:sigmadock,chap:sigmaflow}),
with no cross term anywhere, provided the metric
\eqref{eq:se3-product-metric-scratch} --- rather than some other, coupled
metric --- is the one actually in force. \Cref{chap:sigmadock} verifies,
directly against the code, that this decoupled treatment is indeed what is
implemented (no coupling term appears in either
\pyfunc{SE3Diffuser}/\pyfunc{SE3\_FlowMatcher}, both of which are pure
componentwise dispatchers, per the code-research findings recorded there).

\begin{definition}[The pose manifold of $N$ rigid fragments]
\label{def:pose-manifold-scratch}
The state space of a docking model that represents a ligand as $N$
independently posable rigid fragments (\Cref{sec:sigmadock-state-space})
is the product manifold
\[
\mathcal M_N = SE(3)^N,
\qquad
\mathbf T = \bigl(R^{(1)},x^{(1)},\dots,R^{(N)},x^{(N)}\bigr),
\]
again with the product metric, extending \eqref{eq:se3-product-metric-scratch}
fragmentwise.
\end{definition}

By the same reasoning as \Cref{prop:se3-product} applied $N$ times, every
construction on $\mathcal M_N$ factorises over fragments \emph{and} over
rotation/translation within each fragment: the entire generative model,
seen purely geometrically, is $2N$ independent copies of the two
constructions this document develops in full generality --- one Euclidean
($\RR^3$, for each fragment's translation) and one $SO(3)$ (for each
fragment's rotation) --- coupled \emph{only} through the neural network that
predicts the vector field/score, never through the geometry itself. This is
exactly the content of the ``inter-fragment correlations enter only through
the learned model'' design point discussed in
\Cref{sec:sigmadock-state-space}.

\section{The torsion space $\mathbb T^k$}
\label{sec:torus}

The remaining internal degree of freedom identified in
\Cref{sec:pose-dof}, the dihedral angle, lives on a space we have not yet
defined formally.

\begin{definition}[Circle group and torus]
\label{def:torus}
$\mathbb T := \RR/2\pi\ZZ$, the quotient of $\RR$ by the equivalence
$\phi\sim\phi+2\pi m$ ($m\in\ZZ$), with group operation addition mod $2\pi$;
equivalently, $\mathbb T\cong S^1=\{(\cos\phi,\sin\phi):\phi\in\RR\}\subset\RR^2$
via $\phi\mapsto(\cos\phi,\sin\phi)$, an isomorphism of Lie groups (the
group law on the right is complex multiplication of unit-modulus complex
numbers, matching addition of angles). $\mathbb T^k$ is the $k$-fold direct
product, an abelian Lie group of dimension $k$.
\end{definition}

Being abelian, $\mathbb T^k$'s Lie algebra is simply $\RR^k$ with the
identity map as exponential ($\exp(\phi)=\phi\bmod2\pi$ in each coordinate),
and the natural (bi-invariant, since the group is abelian and hence trivially
so) metric is flat Euclidean --- \emph{locally}. The one genuinely new
phenomenon $\mathbb T^k$ introduces, absent from $\RR^k$, is that
\emph{differences} of angles are only defined up to the $2\pi$-periodicity,
and there are two equally valid representatives differing by a full turn;
the convention used throughout is to take the \emph{shortest} representative.

\begin{definition}[Periodic (shortest-path) angular difference]
\label{def:angular-diff}
For $\phi_1,\phi_2\in\mathbb T$, the geodesic distance is
$d_{\mathbb T}(\phi_1,\phi_2) = \bigl|\bigl((\phi_2-\phi_1+\pi)\bmod2\pi\bigr)-\pi\bigr|\in[0,\pi]$,
and the logarithmic map at $\phi_1$ is the signed version,
$\log_{\phi_1}(\phi_2) = \bigl((\phi_2-\phi_1+\pi)\bmod2\pi\bigr)-\pi \in(-\pi,\pi]$,
the unique representative of $\phi_2-\phi_1\pmod{2\pi}$ lying in $(-\pi,\pi]$.
\end{definition}

\begin{remark}
\label{rem:torus-not-used}
SigmaDock and SigmaFlow, as reconstructed in
\Cref{chap:sigmadock,chap:sigmaflow} from the codebase, do \textbf{not}
generate torsion angles as an explicit degree of freedom at all: the
fragment-based parameterisation of \Cref{sec:sigmadock-state-space} treats
each rigid fragment's \emph{internal} torsions as fixed at their reference-conformer
values (\Cref{sec:conformational-manifold}), and only the $N$ copies of
$SE(3)$ --- never an explicit $\mathbb T^k$ factor --- are generated. This
chapter nonetheless develops $\mathbb T^k$ in full, for two reasons stated
precisely and not merely asserted: first, understanding \emph{why}
SigmaDock avoids torsional generation requires the volume-form argument of
\Cref{sec:sigmadock-state-space}, which itself requires comparing the torus
construction against the fragment construction on equal footing; second,
the torsion angle $\phi_{ijkl}$ of \Cref{sec:pose-dof} is exactly the
periodic dihedral variable this section formalises, and it reappears,
undiffused and ungenerated but still present as a fixed input, throughout
\Cref{chap:sigmadock,chap:sigmaflow}'s discussion of what a rigid fragment
template actually is.
\end{remark}

\begin{remark}[Parity and chirality]
\label{rem:parity-scratch}
We noted in \Cref{def:so3-def} that $O(3)\supsetneq SO(3)$ additionally
contains orientation-\emph{reversing} maps (reflections, $\det=-1$).
Molecules are, in general, chiral: a molecule and its mirror image can be
distinct chemical species (different smell, different biological activity,
etc.), so a docking model must not be invariant under reflection --- a
network that could not distinguish a ligand from its mirror image would
place both equally well into a pocket that in reality binds only one of
them. This is why \Cref{chap:equiformer} tracks a \emph{parity} label
alongside the degree $\ell$ once it introduces representations of $O(3)$
rather than only of $SO(3)$: a rotationally-equivariant-only architecture
that additionally happened to be reflection-invariant would be a strictly
worse inductive bias for this task, not a neutral simplification.
\end{remark}

\clearpage

% ================================================================
\chapter{Steerable Representations: Spherical Harmonics, Wigner Matrices, and Clebsch--Gordan Coefficients}
\label{chap:steerable}
% ================================================================

\Cref{chap:groups} equipped us with $SO(3)$ itself; this chapter equips us
with the machinery to build \emph{features that transform predictably under
$SO(3)$} --- scalars, ordinary $3$-vectors, and their higher-degree
analogues --- and to combine such features while preserving that
predictability. This is not optional scaffolding for the neural network of
\Cref{chap:equiformer}: it \emph{is} the architecture, in the precise sense
that once the objects of this chapter are defined, the network's layers
will turn out to be almost forced, as \Cref{chap:equiformer} shows in
detail. Everything here is representation theory of $SO(3)$, developed only
as far as needed and always tied back to $3$-dimensional physical intuition
(directions, angular momentum) rather than presented abstractly first.

\section{Representations of $SO(3)$}
\label{sec:representations}

\paragraph{Motivation.} A scalar (say, an atom's partial charge) does not
change when the whole molecule is rotated: it is the same number before and
after. A $3$-vector attached to an atom (say, the direction to its nearest
neighbour) \emph{does} change: rotating the molecule by $Q$ rotates the
vector by $Q$ too. Both are examples of the same general phenomenon --- a
feature associated with a rotation-dependent transformation rule --- and we
now formalise the full family of such transformation rules that $SO(3)$
admits.

\begin{definition}[Representation]
\label{def:representation-scratch}
A (finite-dimensional, real) \emph{representation} of $SO(3)$ is a smooth
map $D:SO(3)\to GL(V)$, for some real vector space $V$, satisfying
$D(Q_1Q_2)=D(Q_1)D(Q_2)$ for all $Q_1,Q_2\in SO(3)$ (a group homomorphism
into the invertible linear maps of $V$).
\end{definition}

The homomorphism property is exactly what makes $D$ a consistent
transformation rule: rotating first by $Q_2$ then by $Q_1$ has the same
effect on a feature as rotating once by the composite $Q_1Q_2$, matching
how physical rotations actually compose.

\begin{example}
$V=\RR$, $D(Q)=1$ for all $Q$ (the \emph{trivial} representation): this is
how scalars transform. $V=\RR^3$, $D(Q)=Q$ (the \emph{standard} or
\emph{defining} representation): this is how ordinary vectors transform.
Both are representations of $SO(3)$ in the sense of
\Cref{def:representation-scratch}; direct sums (block-diagonal
concatenations) and tensor products of representations are again
representations (immediate from the definition, checking the homomorphism
property component/factor-wise).
\end{example}

\begin{definition}[Irreducible representation]
\label{def:irrep-scratch}
A representation $D$ on $V$ is \emph{irreducible} (an \emph{irrep}) if $V$
has no proper, nonzero subspace $W\subsetneq V$ (other than $\{0\}$ and $V$
itself) with $D(Q)W\subseteq W$ for every $Q\in SO(3)$.
\end{definition}

Irreducibility is the representation-theoretic analogue of ``cannot be
decomposed further'': if $V$ \emph{did} have such an invariant subspace $W$,
the action of $SO(3)$ on $V$ would really just be two smaller, independent
actions (on $W$ and, in a suitable sense, on a complement), and nothing
would be lost by treating them separately. The classification of all
irreps of $SO(3)$ is the single structural fact this chapter is building
towards, and it is exactly what tells the network designer of
\Cref{chap:equiformer} which ``types'' of feature exist at all.

\begin{theorem}[Classification of the irreps of $SO(3)$]
\label{thm:so3-irrep-classification}
The (real, finite-dimensional) irreducible representations of $SO(3)$ are
indexed by a non-negative integer $\ell=0,1,2,\dots$, called the
\emph{degree}; the irrep of degree $\ell$ acts on a vector space
$V_\ell\cong\RR^{2\ell+1}$, and every finite-dimensional representation of
$SO(3)$ decomposes as a direct sum of copies of these.
\end{theorem}

We do not prove \Cref{thm:so3-irrep-classification} from scratch (it is a
classical result, most naturally proved via the representation theory of
the associated complex Lie algebra $\so(3,\CC)\cong\su(2)$ and its weight
theory, requiring substantially more Lie-algebra machinery than the rest of
this document uses; see, e.g., Hall, \emph{Lie Groups, Lie Algebras, and
Representations}, 2nd ed., 2015, Ch.~4). We instead \emph{construct} the
degree-$\ell$ irrep concretely and explicitly in
\Cref{sec:spherical-harmonics-derivation}, by exhibiting $V_\ell$ as a
specific, checkable space of functions on the sphere and verifying
irreducibility is unnecessary for our purposes once the construction is in
hand --- what we need is the transformation rule
\eqref{eq:sph-equivariance-scratch}, the dimension count $2\ell+1$, and the
tensor-product decomposition of \Cref{thm:cg-decomposition}, all of which
\Cref{sec:spherical-harmonics-derivation,sec:cg-derivation} derive directly.

We record the two cases familiar from \Cref{sec:generative-modelling-intro}
onward once and for all in the new language.

\begin{example}[$\ell=0$ and $\ell=1$]
\label{ex:ell01}
$V_0\cong\RR$ with $D^0(Q)=1$: type-$0$ features are scalars/invariants.
$V_1\cong\RR^3$ with $D^1(Q)=Q$ (the standard representation): type-$1$
features are ordinary $3$-vectors. Every feature encountered so far in this
document --- atomic charges, distances (type $0$); relative position
vectors, the translational and rotational velocity outputs of
\eqref{eq:network-signature-scratch} (type $1$) --- is one of these two
cases; $\ell\ge2$ is new content, introduced precisely because it lets the
network represent \emph{angular} patterns (e.g.\ ``the neighbours lie in a
ring around this axis'') that neither scalars nor single vectors can
express, as \Cref{chap:equiformer} discusses concretely.
\end{example}

\section{Spherical harmonics}
\label{sec:spherical-harmonics-derivation}

\paragraph{Motivation.} We want an explicit, computable realisation of each
degree-$\ell$ irrep --- not just the abstract knowledge that it exists and
has dimension $2\ell+1$, but an actual formula that can be evaluated on a
direction $\hat x\in S^2$ and fed into a neural network. Spherical
harmonics provide exactly this, and the cleanest way to see why they arise
is by analogy with a construction the reader may already know: Fourier
series on the circle.

\paragraph{Intuition: from the circle to the sphere.} A periodic function
on $\mathbb T$ (\Cref{def:torus}) decomposes into Fourier modes
$e^{im\phi}$, $m\in\ZZ$; the mode $e^{im\phi}$ transforms under rotation of
the circle by $\alpha$ (i.e.\ $\phi\mapsto\phi+\alpha$) by an overall phase,
$e^{im(\phi+\alpha)}=e^{im\alpha}e^{im\phi}$ --- a $1$-dimensional
representation of the rotation group of the circle, $SO(2)$, one for each
integer $m$. Spherical harmonics are the exact analogue for functions on
the sphere $S^2$ and the rotation group $SO(3)$: instead of a single integer
$m$ indexing $1$-dimensional pieces, we get a degree $\ell$ indexing
$(2\ell+1)$-dimensional pieces (the ``$m$'' index survives too, now ranging
over $-\ell,\dots,\ell$ \emph{within} a fixed $\ell$, and mixing among
themselves under a rotation, rather than each being separately invariant as
on the circle --- exactly the passage from the abelian $SO(2)$ to the
non-abelian $SO(3)$).

\begin{definition}[Spherical harmonics via harmonic homogeneous polynomials]
\label{def:sph-harm-def}
For $\ell\ge0$, let $\mathcal H_\ell$ be the space of \emph{harmonic}
(satisfying Laplace's equation $\Delta p=0$) \emph{homogeneous polynomials}
of degree $\ell$ in three real variables $(x_1,x_2,x_3)$. The (real)
\emph{spherical harmonics of degree $\ell$} are the restrictions of
$\mathcal H_\ell$ to the unit sphere $S^2$.
\end{definition}

\begin{proposition}[Dimension count]
\label{prop:sph-dim}
$\dim\mathcal H_\ell = 2\ell+1$.
\end{proposition}

\begin{proof}
The space of \emph{all} homogeneous degree-$\ell$ polynomials in $3$
variables has dimension $\binom{\ell+2}{2}$ (the number of monomials
$x_1^{a}x_2^{b}x_3^{c}$ with $a+b+c=\ell$). The Laplacian $\Delta$ maps this
space onto the space of homogeneous degree-$(\ell-2)$ polynomials
(surjectivity is a short, standard computation: given a target
degree-$(\ell-2)$ polynomial $q$, one directly exhibits a degree-$\ell$
polynomial $p$ with $\Delta p=q$, e.g.\ by induction on $\ell$ using that
$\Delta(\|x\|^2q)=2(2\ell+1)q+\|x\|^2\Delta q$ for homogeneous $q$ of degree
$\ell-2$, which for $\ell\ge2$ can be solved for $q$ in terms of a suitable
preimage), and its kernel is exactly $\mathcal H_\ell$. By rank--nullity,
\[
\dim\mathcal H_\ell
=
\binom{\ell+2}{2} - \binom{\ell}{2}
=
\frac{(\ell+2)(\ell+1)}2 - \frac{\ell(\ell-1)}2
=
2\ell+1 .
\]
\end{proof}

\begin{example}[$\ell=0,1$]
\label{ex:sph-low-degree}
$\mathcal H_0=\mathrm{span}\{1\}$ (constants; $\Delta1=0$ trivially),
restricting to the constant function on $S^2$, matching $\dim=1=2(0)+1$.
$\mathcal H_1=\mathrm{span}\{x_1,x_2,x_3\}$ (linear functions are
automatically harmonic, since $\Delta$ of a linear function is $0$),
restricting to the three coordinate functions on $S^2$, matching
$\dim=3=2(1)+1$; note that restricted to $S^2$, these are literally the
components of the unit vector $\hat x$ itself.
\end{example}

\begin{proposition}[Transformation rule of spherical harmonics]
\label{prop:sph-transform}
Let $Q\in SO(3)$ act on functions $f:S^2\to\RR$ by $(Q\cdot f)(\hat x) = f(Q^\top\hat x)$
(the natural action: the rotated function evaluated at $\hat x$ equals the
original function evaluated at the pre-rotation direction). Then this
action preserves $\mathcal H_\ell$ (restricted to $S^2$) for each $\ell$;
consequently, in any basis $Y^\ell_{-\ell},\dots,Y^\ell_\ell$ of degree-$\ell$
spherical harmonics, there is a matrix $D^\ell(Q)\in\RR^{(2\ell+1)\times(2\ell+1)}$
with
\begin{equation}
\label{eq:sph-equivariance-scratch}
Y^\ell_m(Q^\top\hat x) = \sum_{m'} D^\ell(Q)_{mm'}\,Y^\ell_{m'}(\hat x),
\qquad\text{equivalently}\qquad
Y^\ell(Q\hat x) = D^\ell(Q)\,Y^\ell(\hat x)
\end{equation}
(the second form obtained by relabelling $\hat x\mapsto Q\hat x$ and using
$Q^\top Q=I$), and $D^\ell$ is a representation of $SO(3)$ in the sense of
\Cref{def:representation-scratch}, of dimension $2\ell+1$, matching
\Cref{thm:so3-irrep-classification}'s degree-$\ell$ irrep.
\end{proposition}

\begin{proof}
If $p\in\mathcal H_\ell$ (a harmonic homogeneous degree-$\ell$ polynomial on
$\RR^3$), then $p\circ Q^\top$ (i.e.\ $x\mapsto p(Q^\top x)$) is again
homogeneous of degree $\ell$ (composing with a linear map preserves
homogeneity), and harmonic (the Laplacian is invariant under orthogonal
changes of coordinates: $\Delta(p\circ Q^\top)=(\Delta p)\circ Q^\top$ for
$Q$ orthogonal, a standard fact provable by the chain rule and
$Q^\top Q=I$), so $p\circ Q^\top\in\mathcal H_\ell$ too: the action preserves
$\mathcal H_\ell$. Since the action is linear on the finite-dimensional
space $\mathcal H_\ell$, expressing $Q\cdot Y^\ell_m$ in the chosen basis
gives exactly a matrix $D^\ell(Q)$ as claimed, and
$D^\ell(Q_1Q_2)=D^\ell(Q_1)D^\ell(Q_2)$ follows because the action of $Q$ on
functions is itself a group action (i.e.\ $(Q_1Q_2)\cdot f = Q_1\cdot(Q_2\cdot f)$,
a direct check from the definition), which is exactly the homomorphism
property required by \Cref{def:representation-scratch}.
\end{proof}

We do not additionally re-derive \emph{irreducibility} of $D^\ell$ here (it
follows from the classical theory cited after
\Cref{thm:so3-irrep-classification}, and the identification
$V_\ell\cong\mathcal H_\ell$ is the standard one); what \Cref{prop:sph-transform}
already gives us --- an explicit, computable, correctly-transforming basis
for every degree $\ell$ --- is everything the rest of this document needs.

\begin{codebox}
\textbf{Mathematical object:} $Y^\ell:S^2\to\RR^{2\ell+1}$, degree-$\ell$ real
spherical harmonics; the equivariant edge embedding
$\bigl(\psi(d_{ij}),Y^\ell(\hat x_{ij})\bigr)_{\ell\le L}$ \\
\textbf{Implementation:} \file{net/edge\_rot\_mat.py} (per-edge frame
alignment, see \Cref{chap:equiformer}) together with the \texttt{e3nn}
library's spherical-harmonics routines, called from within the
$SO(2)$-convolution machinery of \Cref{sec:equiformer-attention} rather than
as an explicit, separately-visible tensor-product step (a structural point
this document flags precisely, not glossed over, in
\Cref{sec:equiformer-attention}) \\
\textbf{Interpretation:} in EquiformerV2, $L=\texttt{lmax}=3$ by default
(\Cref{chap:equiformer}), so the code works with degrees $\ell\in\{0,1,2,3\}$,
i.e.\ steerable features of dimension $1+3+5+7=16$ per channel.
\end{codebox}

\subsection{Explicit low-degree formulas, and the connection to familiar operations}

We record the degree $0,1,2$ spherical harmonics explicitly (up to the
normalisation convention, which varies across sources and is immaterial to
every argument in this document, since it only rescales each $\ell$-block
by a constant that a subsequent learned linear layer can absorb):
\[
Y^0 = 1,
\qquad
Y^1 = (x_1,x_2,x_3)=\hat x,
\qquad
Y^2 \propto \bigl(x_1x_2,\ x_2x_3,\ 2x_3^2-x_1^2-x_2^2,\ x_3x_1,\ x_1^2-x_2^2\bigr) ,
\]
the last being (up to normalisation and a choice of basis) the five
independent components of the traceless symmetric matrix $\hat x\hat x^\top-\tfrac13I$,
i.e.\ exactly a (traceless) quadrupole moment of the direction $\hat x$.
This connects the abstract degree $\ell$ directly to physical/geometric
intuition: $\ell=0$ answers ``is there something here'' (an unsigned
scalar); $\ell=1$ answers ``in what single direction'' (a vector, capable of
representing one preferred direction); $\ell=2$ answers ``along what
\emph{axis}, ignoring sign'' (a nematic/quadrupolar pattern, e.g.\
distinguishing a planar ring of neighbours in the $xy$-plane from an
isotropic distribution, something no single vector can encode, since the
ring itself has no preferred in-plane direction and $Y^1$ averaged over the
ring would cancel to zero while $Y^2$ would not); higher $\ell$ resolve
correspondingly finer angular structure.

\section{Wigner-$D$ matrices}
\label{sec:wigner-d-derivation}

\Cref{prop:sph-transform} already \emph{is} the definition of the
Wigner-$D$ matrix: $D^\ell(Q)$ is, by construction, the matrix realising
$SO(3)$'s action on degree-$\ell$ spherical harmonics in a chosen basis. We
isolate the properties of $D^\ell$ that are used repeatedly.

\begin{proposition}[Basic properties of Wigner-$D$ matrices]
\label{prop:wigner-properties}
For every $\ell$: (i) $D^\ell(I)=I_{2\ell+1}$; (ii)
$D^\ell(Q_1Q_2)=D^\ell(Q_1)D^\ell(Q_2)$ (\Cref{prop:sph-transform}); (iii)
$D^\ell(Q)$ is orthogonal, $D^\ell(Q)^\top D^\ell(Q)=I$ (since the action of
\Cref{prop:sph-transform} preserves the natural $L^2(S^2)$ inner product,
under which a suitable choice of real spherical-harmonic basis is
orthonormal); (iv) $D^0(Q)=1$ for all $Q$ (the trivial representation, as
in \Cref{ex:ell01}); $D^1(Q)=Q$ itself, in the standard Cartesian basis
$Y^1=\hat x$ of \Cref{ex:sph-low-degree}.
\end{proposition}

\begin{proof}
(i) and (ii) restate \Cref{prop:sph-transform}. (iii): the action
$f\mapsto f\circ Q^\top$ preserves the rotation-invariant surface measure
on $S^2$, hence preserves the $L^2(S^2)$ inner product
$\langle f,g\rangle=\int_{S^2}fg\,d\Omega$; choosing $Y^\ell_m$ orthonormal
in this inner product (always possible by Gram--Schmidt within the
finite-dimensional space $\mathcal H_\ell|_{S^2}$) makes $D^\ell(Q)$ an
isometry of $\RR^{2\ell+1}$ with the standard inner product, i.e.\
orthogonal. (iv): $D^0(Q)=1$ since $Y^0\equiv1$ is fixed pointwise by every
rotation of its argument; $D^1(Q)=Q$ is immediate from
\Cref{eq:sph-equivariance-scratch} with $Y^1(\hat x)=\hat x$:
$Y^1(Q\hat x)=Q\hat x=Q\,Y^1(\hat x)$.
\end{proof}

\Cref{prop:wigner-properties}(iv) is worth dwelling on: it says that
$SO(3)$'s \emph{defining} representation (rotation matrices themselves,
which is how we have described $SO(3)$ throughout \Cref{chap:groups}) is
nothing but the degree-$1$ Wigner matrix. Every subsequent degree
$\ell\ge2$ is a genuinely new object, not reducible to ordinary matrix
multiplication by a rotation matrix, and this is exactly why
\Cref{chap:equiformer}'s architecture cannot simply reuse the familiar
$3\times3$ rotation-matrix machinery once $\ell\ge2$ features are involved
--- it needs the general $D^\ell(Q)$.

\begin{codebox}
\textbf{Mathematical object:} $D^\ell(Q)\in\RR^{(2\ell+1)\times(2\ell+1)}$, the
per-edge rotation applied to steerable features \\
\textbf{Implementation:} \file{net/wigner.py} (function assembling
$D^\ell$ block-diagonally from Euler angles, via \texttt{e3nn}), called from
\file{net/edge\_rot\_mat.py}::\pyfunc{RotationToWignerDMatrix} \\
\textbf{Tensor shape:} one $D^\ell$ per edge, per degree $\ell\le L$,
assembled into a single block-diagonal matrix of size
$\sum_{\ell=0}^L(2\ell+1)$ acting on the full multi-degree feature at once \\
\textbf{Interpretation:} every edge in the molecular graph gets its own
per-edge frame (\Cref{chap:equiformer}, \Cref{sec:equiformer-frame}), built
from the edge direction $\hat x_{ij}$; the corresponding Wigner-$D$ matrix
rotates a neighbour's steerable feature \emph{into} that edge-aligned
frame before any further processing, and the inverse rotates the result
back --- exactly the mechanism, detailed in \Cref{chap:equiformer}, by
which EquiformerV2 reduces full $SO(3)$-equivariant processing to a
computationally cheaper $SO(2)$-equivariant one.
\end{codebox}

\section{Tensor products and Clebsch--Gordan coefficients}
\label{sec:cg-derivation}

\paragraph{Motivation.} Message passing (\Cref{sec:gnn-scratch}) needs to
\emph{combine} two features --- e.g.\ a neighbour's current steerable
feature $h_j^{\ell_1}$ and the spherical-harmonic embedding $Y^{\ell_2}(\hat x_{ij})$
of the edge direction --- into a new feature. For ordinary (type-$0$,
scalar) features, ``combine'' just means multiply. For steerable features
of possibly different, positive degree, an arbitrary bilinear combination
will \emph{not} in general transform correctly under rotation; we need the
\emph{specific} bilinear combinations that do.

\begin{definition}[Tensor product representation]
\label{def:tensor-product-rep}
Given representations $D^{\ell_1}$ on $V_{\ell_1}$ and $D^{\ell_2}$ on
$V_{\ell_2}$, the \emph{tensor product representation} acts on
$V_{\ell_1}\otimes V_{\ell_2}$ by $(D^{\ell_1}\otimes D^{\ell_2})(Q)(u\otimes v) = D^{\ell_1}(Q)u\otimes D^{\ell_2}(Q)v$.
\end{definition}

This is again a representation of $SO(3)$ (immediate check of the
homomorphism property), but --- for $\ell_1,\ell_2\ge1$ --- it is generally
\emph{not} irreducible: $\dim(V_{\ell_1}\otimes V_{\ell_2})=(2\ell_1+1)(2\ell_2+1)$,
which for e.g.\ $\ell_1=\ell_2=1$ is $9$, not matching any single irrep
dimension $2\ell+1$ (odd numbers $1,3,5,7,\dots$, and $9=2(4)+1$ \emph{is}
odd, but we will see explicitly that the $9$-dimensional space splits into
\emph{three} pieces, of dimensions $1+3+5=9$, not one piece of dimension $9$).
By \Cref{thm:so3-irrep-classification}, every representation --- including
this one --- decomposes as a direct sum of irreps; the following theorem
says precisely which ones appear.

\begin{theorem}[Clebsch--Gordan decomposition]
\label{thm:cg-decomposition}
\[
V_{\ell_1}\otimes V_{\ell_2}
\ \cong\
\bigoplus_{\ell=|\ell_1-\ell_2|}^{\ell_1+\ell_2} V_\ell .
\]
\end{theorem}

\begin{proof}[Proof sketch, and a dimension check]
We do not derive \Cref{thm:cg-decomposition} from first principles (the
standard proof uses the weight/character theory of
$\so(3,\CC)\cong\su(2)$ cited after \Cref{thm:so3-irrep-classification}: the
weights of $V_{\ell_1}\otimes V_{\ell_2}$ are all sums $m_1+m_2$ with
$m_i\in\{-\ell_i,\dots,\ell_i\}$, and grouping these sums by which range of
consecutive integers they fill out recovers exactly the stated range of
$\ell$; see Hall (2015), Thm.~4.44 for $\su(2)$, transported to $SO(3)$ via
the double cover). We record instead the dimension identity that this
decomposition must, and does, satisfy, as a consistency check the reader can
verify unaided:
\[
\sum_{\ell=|\ell_1-\ell_2|}^{\ell_1+\ell_2}(2\ell+1)
=
(2\ell_1+1)(2\ell_2+1),
\]
provable by writing the left side as an arithmetic series (first term
$2|\ell_1-\ell_2|+1$, last term $2(\ell_1+\ell_2)+1$, number of terms
$\ell_1+\ell_2-|\ell_1-\ell_2|+1$, which for $\ell_1\ge\ell_2$ equals
$2\ell_2+1$) and simplifying; for $\ell_1=\ell_2=1$ this gives
$1+3+5=9=3\times3$, confirming the example above.
\end{proof}

\begin{definition}[Clebsch--Gordan coefficients and tensor product]
\label{def:cg-scratch}
The isomorphism of \Cref{thm:cg-decomposition}, written out in coordinates,
defines real numbers $C^{(\ell,m)}_{(\ell_1,m_1)(\ell_2,m_2)}$ (the
\emph{Clebsch--Gordan coefficients}, uniquely determined up to an overall
sign/phase convention per $(\ell_1,\ell_2,\ell)$ triple by orthonormality of
the bases involved) such that, for $u\in V_{\ell_1}$, $v\in V_{\ell_2}$, the
degree-$\ell$ component of their \emph{Clebsch--Gordan tensor product} is
\begin{equation}
\label{eq:cg-scratch}
\bigl(u\otimes_{\mathrm{cg}}v\bigr)^\ell_m
=
\sum_{m_1=-\ell_1}^{\ell_1}\sum_{m_2=-\ell_2}^{\ell_2}
C^{(\ell,m)}_{(\ell_1,m_1)(\ell_2,m_2)}\;u^{\ell_1}_{m_1}v^{\ell_2}_{m_2},
\qquad
|\ell_1-\ell_2|\le\ell\le\ell_1+\ell_2 .
\end{equation}
\end{definition}

\begin{proposition}[The Clebsch--Gordan product is equivariant]
\label{prop:cg-equivariant}
For $Q\in SO(3)$: $\bigl(D^{\ell_1}(Q)u\bigr)\otimes_{\mathrm{cg}}\bigl(D^{\ell_2}(Q)v\bigr) = D^\ell(Q)\bigl(u\otimes_{\mathrm{cg}}v\bigr)$
in each degree-$\ell$ component.
\end{proposition}

\begin{proof}
This is precisely the statement that
\eqref{eq:cg-scratch} realises an intertwiner (an equivariant linear map)
$V_{\ell_1}\otimes V_{\ell_2}\to V_\ell$: by construction (\Cref{def:cg-scratch}),
the Clebsch--Gordan coefficients are exactly the coordinates of the
isomorphism from \Cref{thm:cg-decomposition}, and that isomorphism, being
built from the representation-preserving decomposition itself, commutes
with the group action by definition of what it means for the direct sum
decomposition to consist of \emph{sub}representations.
\end{proof}

\Cref{prop:cg-equivariant} is the single fact that makes
\eqref{eq:cg-scratch} the ``$\times$'' of equivariant deep learning: it is
--- up to the choice of basis fixing the specific numbers
$C^{(\ell,m)}_{(\ell_1,m_1)(\ell_2,m_2)}$, a choice we do not need to make
explicit anywhere in this document, since every use of
\eqref{eq:cg-scratch} in \Cref{chap:equiformer} treats them as fixed,
precomputed constants --- the \emph{unique} (up to scale, by Schur's lemma,
\Cref{prop:schur-scratch} below) bilinear map $V_{\ell_1}\times V_{\ell_2}\to V_\ell$
respecting the rotation action, and it is therefore the only mechanism by
which two steerable features can be combined into a third without breaking
equivariance.

\begin{example}[Familiar special cases]
\label{ex:cg-special-cases}
$\ell_1=\ell_2=1$, $\ell=0$: the degree-$0$ (invariant) component of
$u\otimes_{\mathrm{cg}}v$ for $u,v\in\RR^3$ is (up to normalisation) the
ordinary dot product $\langle u,v\rangle$ --- the unique way to combine two
vectors into a rotation-invariant scalar, matching
\Cref{prop:invariance-obstruction-scratch}'s later discussion of what
distance-only networks can express. $\ell_1=\ell_2=1$, $\ell=1$: the
degree-$1$ component is (up to normalisation) the cross product $u\times v$
--- matching \Cref{prop:hat-cross}'s appearance of the cross product as an
$\so(3)$-valued (equivalently, since $\so(3)\cong\RR^3$, degree-$1$) object.
$\ell_1=0$: tensoring with a scalar is just scalar multiplication,
$(c\otimes_{\mathrm{cg}}v)^\ell = cv$ (only $\ell=\ell_2$ appears, since
$|\ell_1-\ell_2|=\ell_1+\ell_2=\ell_2$ when $\ell_1=0$), recovering ordinary
scaling of a steerable feature by an invariant weight --- the operation
that, per \Cref{prop:schur-scratch} below, is the \emph{only} linear
operation available on a single feature and hence the mechanism by which
every learned scalar weight in \Cref{chap:equiformer}'s network acts on
higher-degree channels.
\end{example}

\section{Consequences for equivariant linear maps and nonlinearities}
\label{sec:schur-consequences}

We record now, as consequences of the theory developed in this chapter, the
two structural constraints on any equivariant neural network layer, since
they are used repeatedly and verified directly against the code in
\Cref{chap:equiformer}.

\begin{proposition}[Schur's lemma, and the form of equivariant linear maps]
\label{prop:schur-scratch}
Let $W:V_{\ell_1}\to V_{\ell_2}$ be linear and satisfy
$W\,D^{\ell_1}(Q) = D^{\ell_2}(Q)\,W$ for all $Q\in SO(3)$ (an
\emph{intertwiner}). If $\ell_1\ne\ell_2$, then $W=0$. If $\ell_1=\ell_2=\ell$,
then $W=cI$ for some scalar $c\in\RR$.
\end{proposition}

\begin{proof}
This is Schur's lemma for real irreducible representations of a compact
group, specialised to $SO(3)$; we take the general statement as classical
(see, e.g., Fulton \& Harris, \emph{Representation Theory}, 1991, \S1.2,
noting that the real irreps of $SO(3)$ are of \emph{real type} --- their
complexifications remain irreducible, unlike, e.g., certain irreps of other
groups where the commutant would be $\CC$ rather than $\RR$ --- which is
exactly why the commutant here is $\RR I$ and not a larger algebra; this
real-type property is itself part of the classical classification cited
after \Cref{thm:so3-irrep-classification} and is not re-derived here).
\end{proof}

\begin{proposition}[Consequence for multi-channel equivariant linear layers]
\label{prop:multichannel-linear}
Let $h=\bigoplus_{\ell,c}h^{\ell,c}$ be a steerable feature with $C_\ell$
channels of each degree $\ell\le L$ (\Cref{def:steerable-scratch} below),
and let $W:\bigoplus_{\ell,c}V_\ell\to\bigoplus_{\ell,c'}V_\ell$ be linear
and equivariant (commuting with the diagonal action
$\bigoplus_{\ell,c}D^\ell(Q)$ on both sides). Then $W$ acts as
$h^{\ell,c}\mapsto\sum_{c'}W^\ell_{cc'}h^{\ell,c'}$ for scalar weights
$W^\ell_{cc'}\in\RR$: $W$ mixes channels only \emph{within} a fixed degree
$\ell$, never across degrees.
\end{proposition}

\begin{proof}
Write $W$ blockwise, indexed by (source degree, source channel) and
(target degree, target channel); each block
$W_{(\ell,c)\to(\ell',c')}:V_\ell\to V_{\ell'}$ is itself an intertwiner
(since $W$ commutes with the full diagonal action, and one may test
equivariance against elements supported in a single source/target block by
linearity), so \Cref{prop:schur-scratch} applies to each block separately:
it vanishes unless $\ell=\ell'$, and is a scalar multiple of the identity
when $\ell=\ell'$. Collecting the surviving (same-degree) blocks gives
exactly the stated form.
\end{proof}

\begin{definition}[Steerable feature; restated]
\label{def:steerable-scratch}
A \emph{steerable feature} (of a network node or edge) is a direct sum
$h=\bigoplus_{\ell=0}^L\bigoplus_{c=1}^{C_\ell}h^{\ell,c}$, $h^{\ell,c}\in V_\ell\cong\RR^{2\ell+1}$,
transforming under $Q\in SO(3)$ by $h\mapsto\bigl(\bigoplus_{\ell,c}D^\ell(Q)\bigr)h$
--- i.e.\ every one of the $C_\ell$ channels at degree $\ell$ transforms by
the \emph{same} $D^\ell(Q)$, independently of the others.
\end{definition}

\Cref{prop:multichannel-linear} is, in words, the entire content of ``an
equivariant linear layer has one learnable scalar per (source channel,
target channel, shared degree) triple, and cannot move information between
degrees'': every degree $\ell$ is processed by its \emph{own}, otherwise
completely ordinary, linear layer, exactly the pattern the
\pyfunc{SO3\_LinearV2} code correspondence in \Cref{chap:equiformer} will
confirm directly. Since linear maps can never mix degrees, and since
pointwise nonlinearities (e.g.\ $\mathrm{ReLU}$ applied entrywise) do
\emph{not}, in general, commute with $D^\ell(Q)$ for $\ell\ge1$ (a
direction-dependent nonlinearity applied to a rotated vector is generally
not the same as rotating the nonlinearly-transformed vector), the
Clebsch--Gordan product \eqref{eq:cg-scratch} of this chapter is the
\emph{only} remaining mechanism, besides invariant-gated rescaling
(\Cref{chap:equiformer}'s gate nonlinearity), by which an equivariant
network can create new degrees of information or genuinely nonlinearly
combine two existing steerable features --- and this is precisely why the
tensor product, rather than being one architectural ingredient among many,
is unavoidable, as promised at the end of \Cref{sec:gnn-scratch}.

\clearpage

% ================================================================
\part{Diffusion Models}
\label{part:diffusion}
% ================================================================

\chapter{Diffusion Models in $\RR^d$}
\label{chap:diffusion}

We now build the first of the two generative-modelling machineries
identified in \Cref{sec:generative-modelling-intro}: a way to construct a
probability path $(p_t)_{t\in[0,1]}$ implicitly, as the law of a stochastic
process, and to sample from $p_1=p_{\mathrm{data}}$ by simulating that
process. Everything in this chapter is developed in $\RR^d$ first; the
manifold (in particular $SO(3)$) generalisation is
\Cref{chap:diffusion-so3}, and SigmaDock's actual use of both is
reconstructed in full, with code correspondence, in \Cref{chap:sigmadock}.

\section{Denoising score matching}
\label{sec:score-matching-intro}

\Cref{sec:langevin} showed that knowing $\nabla_x\log p_{\mathrm{data}}$
would suffice for sampling, via Langevin dynamics, but left open how to
\emph{learn} this score from samples alone. The direct approach --- fit a
network $s_\theta(x)\approx\nabla_x\log p_{\mathrm{data}}(x)$ by minimising
$\EE_{x\sim p_{\mathrm{data}}}\bigl[\|s_\theta(x)-\nabla_x\log p_{\mathrm{data}}(x)\|^2\bigr]$
--- is circular: the target $\nabla_x\log p_{\mathrm{data}}$ is exactly the
unknown quantity. The resolution, exactly analogous in structure to
\Cref{prop:marginalisation}'s role in \Cref{sec:generative-modelling-intro},
is to introduce a \emph{noised} version of the data whose score \emph{is}
computable in closed form.

\begin{definition}[Gaussian perturbation kernel]
\label{def:gaussian-kernel}
For $\sigma>0$, let $p_\sigma(x\mid z) = \mathcal N(x\mid z,\sigma^2I_d)$,
and let $p_\sigma(x) = \int p_\sigma(x\mid z)\,p_{\mathrm{data}}(z)\,dz$ be
the corresponding \emph{noised marginal} (\Cref{prop:marginalisation}, with
$p_\sigma(\cdot\mid z)$ playing the role of the conditional density there).
\end{definition}

By \Cref{rem:gaussian-score}, the \emph{conditional} score is available in
closed form:
\begin{equation}
\label{eq:gaussian-conditional-score}
\nabla_x\log p_\sigma(x\mid z) = -\frac{x-z}{\sigma^2} .
\end{equation}
What we actually want, to run Langevin dynamics
targeting $p_\sigma$ rather than $p_{\mathrm{data}}$ directly (an
approximation we return to and justify via the noise ladder / SDE
construction of \Cref{sec:sde-perspective}), is the \emph{marginal} score
$\nabla_x\log p_\sigma(x)$, which is \emph{not} available in closed form,
since it requires the (intractable) integral defining $p_\sigma$. The
following theorem is the resolution, and its proof is the prototype for
every ``conditional loss has the same minimiser as marginal loss'' argument
in this document, including \Cref{thm:cfm-equivalence-scratch} in
\Cref{chap:flow-matching}.

\begin{theorem}[Denoising score matching]
\label{thm:dsm-equivalence}
Define
\[
\mathcal L_{\mathrm{ESM}}(\theta) = \EE_{x\sim p_\sigma}\Bigl[\bigl\|s_\theta(x)-\nabla_x\log p_\sigma(x)\bigr\|^2\Bigr]
\quad\text{(explicit/marginal score matching)},
\]
\begin{equation}
\label{eq:dsm}
\mathcal L_{\mathrm{DSM}}(\theta) = \EE_{z\sim p_{\mathrm{data}},\,x\sim p_\sigma(\cdot\mid z)}\Bigl[\bigl\|s_\theta(x)-\nabla_x\log p_\sigma(x\mid z)\bigr\|^2\Bigr]
\quad\text{(denoising score matching)} .
\end{equation}
Then $\mathcal L_{\mathrm{ESM}}(\theta) = \mathcal L_{\mathrm{DSM}}(\theta) + C$
with $C$ independent of $\theta$; in particular the two losses have
identical gradients in $\theta$ and the same minimisers.
\end{theorem}

\begin{proof}
Expand the square in $\mathcal L_{\mathrm{ESM}}$ under the marginal $p_\sigma(x)$:
\[
\mathcal L_{\mathrm{ESM}}(\theta)
=
\EE_{x\sim p_\sigma}\bigl[\|s_\theta(x)\|^2\bigr]
- 2\,\EE_{x\sim p_\sigma}\bigl[\langle s_\theta(x),\nabla_x\log p_\sigma(x)\rangle\bigr]
+ \EE_{x\sim p_\sigma}\bigl[\|\nabla_x\log p_\sigma(x)\|^2\bigr] .
\]
The third term does not depend on $\theta$; call it $C_1$. The first term
already matches $\mathcal L_{\mathrm{DSM}}$'s first term, since
$x\sim p_\sigma$ (marginally) is the same distribution as $x$ obtained by
first drawing $z\sim p_{\mathrm{data}}$ then $x\sim p_\sigma(\cdot\mid z)$
(\Cref{prop:marginalisation}), so
$\EE_{x\sim p_\sigma}[\|s_\theta(x)\|^2] = \EE_{z,x}[\|s_\theta(x)\|^2]$.
It remains to handle the cross term. Using
$p_\sigma(x)\nabla_x\log p_\sigma(x) = \nabla_x p_\sigma(x) = \int \nabla_x p_\sigma(x\mid z)\,p_{\mathrm{data}}(z)\,dz
= \int p_\sigma(x\mid z)\nabla_x\log p_\sigma(x\mid z)\,p_{\mathrm{data}}(z)\,dz$
(differentiating under the integral sign, and using
$\nabla_x p_\sigma(x\mid z)=p_\sigma(x\mid z)\nabla_x\log p_\sigma(x\mid z)$),
\begin{align*}
\EE_{x\sim p_\sigma}\bigl[\langle s_\theta(x),\nabla_x\log p_\sigma(x)\rangle\bigr]
&=
\int \langle s_\theta(x),\nabla_x\log p_\sigma(x)\rangle\,p_\sigma(x)\,dx\\
&=
\int\!\!\int \langle s_\theta(x),\nabla_x\log p_\sigma(x\mid z)\rangle\,p_\sigma(x\mid z)\,p_{\mathrm{data}}(z)\,dz\,dx\\
&=
\EE_{z\sim p_{\mathrm{data}},\,x\sim p_\sigma(\cdot\mid z)}\bigl[\langle s_\theta(x),\nabla_x\log p_\sigma(x\mid z)\rangle\bigr],
\end{align*}
which is exactly $\mathcal L_{\mathrm{DSM}}$'s cross term. Collecting terms,
$\mathcal L_{\mathrm{ESM}}(\theta) = \mathcal L_{\mathrm{DSM}}(\theta) + C_1 - C_2$,
where $C_2=\EE_{z,x}[\|\nabla_x\log p_\sigma(x\mid z)\|^2]$ is also
$\theta$-independent, giving $C=C_1-C_2$.
\end{proof}

The practical content of \Cref{thm:dsm-equivalence}: to train a score
network, draw a clean data point $z\sim p_{\mathrm{data}}$ (a sample from
the training set), draw a noised version
$x = z + \sigma\varepsilon$, $\varepsilon\sim\mathcal N(0,I_d)$
(sampling from $p_\sigma(\cdot\mid z)$ directly, rather than needing the
marginal), and regress $s_\theta(x)$ against the \emph{known}, closed-form
target $-\varepsilon/\sigma$ (substituting $x-z=\sigma\varepsilon$ into
\eqref{eq:gaussian-conditional-score}). No integral over
$p_{\mathrm{data}}$ or $p_\sigma$ is ever evaluated.

\section{The SDE perspective, and why a single noise level is not enough}
\label{sec:sde-perspective}

A single $\sigma$ is not sufficient for high-quality sampling: for large
$\sigma$, $p_\sigma\approx p_{\mathrm{init}}$ is easy to sample from via
Langevin dynamics (the perturbation dominates, and the resulting density is
close to an isotropic Gaussian, which mixes quickly under Langevin
dynamics) but far from $p_{\mathrm{data}}$, and vice versa for small
$\sigma$. Song and Ermon's original resolution
(Song \& Ermon, \emph{Generative Modeling by Estimating Gradients of the
Data Distribution}, NeurIPS 2019) trains a single, noise-conditional network
$s_\theta(x,\sigma)$ across a decreasing ladder of noise levels
$\sigma_1>\dots>\sigma_N$ and runs annealed Langevin dynamics, decreasing
$\sigma$ over the course of sampling. Song et al.\ (\emph{Score-Based
Generative Modeling through Stochastic Differential Equations}, ICLR 2021)
subsequently observed that this discrete ladder is the time-discretisation
of a single, continuous-time stochastic process, and it is this continuous
formulation, not the discrete ladder, that both SigmaDock's actual
implementation and this document use throughout.

Write the \emph{noising} process in the reversed time variable $s=1-t$
(\Cref{rem:time-convention-global}), so it runs from data ($s=0$) to noise
($s=1$):
\begin{equation}
\label{eq:forward-sde-scratch}
dY_s = f(Y_s,s)\,ds + g(s)\,dW_s,
\qquad
Y_0\sim p_{\mathrm{data}} ,
\end{equation}
with drift $f:\RR^d\times[0,1]\to\RR^d$, diffusion coefficient
$g:[0,1]\to\RR_{\ge0}$, and $W$ a standard Wiener process
(\Cref{def:wiener}). Write $p_t$ for the law of $Y_{1-t}$ (i.e.\ time is
relabelled to our global generation-time convention, so $p_0$ is the law of
$Y_1$, the fully-noised end, and $p_1$ is the law of $Y_0=p_{\mathrm{data}}$).

\begin{example}[Variance-exploding (VE) SDE]
\label{ex:ve-sde}
$f\equiv0$, $g(s)^2 = \tfrac{d}{ds}\sigma(s)^2$ for an increasing schedule
$\sigma(s)$. Then \eqref{eq:forward-sde-scratch} integrates in closed form
to $Y_s = Y_0 + \sigma(s)\varepsilon$, $\varepsilon\sim\mathcal N(0,I_d)$
(a direct check: this is a driftless SDE whose solution is Brownian motion
time-changed by $\sigma(s)^2$, and its marginal variance at time $s$ is
exactly $\sigma(s)^2$ by It\^o isometry, which we do not derive here but
which matches the perturbation kernel of \Cref{def:gaussian-kernel} with
$\sigma=\sigma(s)$) --- exactly recovering the noise ladder of
\Cref{sec:score-matching-intro} as the time slices $\sigma(s_i)=\sigma_i$
of a single continuous process. This is the SDE family used, as
\Cref{chap:sigmadock} confirms directly against the code, by SigmaDock's
\emph{rotational} diffusion process.
\end{example}

\begin{example}[Variance-preserving (VP) SDE]
\label{ex:vp-sde}
$f(x,s)=-\tfrac12\beta(s)x$, $g(s)=\sqrt{\beta(s)}$, for a noise-rate
schedule $\beta(s)\ge0$. Solving the resulting (linear, hence exactly
solvable) SDE gives
\[
Y_s = \alpha(s)\,Y_0 + \sqrt{1-\alpha(s)^2}\,\varepsilon,
\qquad
\alpha(s) = \exp\Bigl(-\tfrac12\int_0^s\beta(r)\,dr\Bigr),
\]
i.e.\ $p_s(\cdot\mid Y_0) = \mathcal N\bigl(\alpha(s)Y_0,\ (1-\alpha(s)^2)I_d\bigr)$
--- unlike the VE case, the \emph{signal} $Y_0$ is itself attenuated, and
the total variance $\alpha(s)^2\cdot\mathrm{Var}(Y_0) + (1-\alpha(s)^2)$
stays bounded (indeed, if $\mathrm{Var}(Y_0)=1$, stays exactly at $1$) as
$s\to1$, rather than exploding --- hence the name. This is the SDE family
used, as \Cref{chap:sigmadock} again confirms directly, by SigmaDock's
\emph{translational} diffusion process; the exact solution just derived is
this document's own re-derivation of the closed-form perturbation kernel
that the code computes via \pyfunc{R3Diffuser}'s
\pyfunc{integrated\_beta\_t}, not a claim quoted from elsewhere.
\end{example}

We record, once, the general fact that both examples instantiate, since it
is exactly \Cref{def:gaussian-kernel}'s Gaussian-perturbation-kernel
construction generalised to an SDE with (possibly time- and
state-dependent) linear drift.

\begin{proposition}[Perturbation kernel of a linear SDE]
\label{prop:linear-sde-kernel}
If $f(x,s)=-\tfrac12\beta(s)x$ (including $\beta\equiv0$, the VE case) and
$g(s)^2=\tfrac{d}{ds}\sigma(s)^2$ is otherwise free, the solution of
\eqref{eq:forward-sde-scratch} has Gaussian conditional law
$p_s(\cdot\mid Y_0)=\mathcal N\bigl(\alpha(s)Y_0,\beta(s)^2 I_d\bigr)$
for schedules $\alpha(s),\beta(s)$ determined by $f,g$ (linear SDEs are
exactly solvable; we do not re-derive the general variation-of-constants
formula here, having exhibited it explicitly for the two cases used by this
document in \Cref{ex:ve-sde,ex:vp-sde}).
\end{proposition}

\subsection{Score, weighting, and the training objective in continuous time}

Given a schedule as in \Cref{prop:linear-sde-kernel}, the conditional score
is again available in closed form by \Cref{rem:gaussian-score}, and the
continuous-time denoising score matching objective
(obtained from the single-noise-level objective \eqref{eq:dsm} by letting
the fixed $\sigma$ become the time-dependent $\beta(s)$ of
\Cref{prop:linear-sde-kernel} and averaging over $s\sim\mathcal U(0,1)$ with
weight $\lambda(s)$ --- the literal continuous-time analogue of
\Cref{thm:dsm-equivalence} applied at every noise level simultaneously) is
\begin{equation}
\label{eq:dsm-continuous}
\mathcal L(\theta) = \EE_{s\sim\mathcal U(0,1),\,Y_0\sim p_{\mathrm{data}},\,Y_s\sim p_s(\cdot\mid Y_0)}
\Bigl[\lambda(s)\bigl\|s_\theta(Y_s,s) - \nabla_{Y_s}\log p_s(Y_s\mid Y_0)\bigr\|^2\Bigr] ,
\end{equation}
for a positive weighting function $\lambda(s)$. The role of $\lambda$ is
purely practical, not conceptual: by \Cref{thm:dsm-equivalence} (applied
noise-level by noise-level, which is valid since the theorem's proof never
uses anything special about a fixed $\sigma$, only that $x$ is drawn from
\emph{some} fixed conditional given $z$), \emph{any} positive $\lambda(s)$
gives a loss with the same minimiser, $s_\theta(x,s)=\nabla_x\log p_s(x)$
--- but the target's own scale, $\|\nabla_x\log p_s(x\mid z)\|=\|\varepsilon\|/\beta(s)$
in \Cref{prop:linear-sde-kernel}'s notation, varies enormously across $s$
(it blows up as $\beta(s)\to0$, i.e.\ as $s\to0$, towards clean data), so an
unweighted loss lets the large-noise (easy, large-target) terms dominate
the gradient at the expense of the small-noise (hard, small-target) terms
that determine fine sample detail. The standard choice,
$\lambda(s)=\beta(s)^2$ (i.e.\ $\lambda(s)=1/\EE\|\nabla_x\log p_s(x\mid z)\|^2$
up to the constant $d$, since $\EE\|\varepsilon/\beta(s)\|^2=d/\beta(s)^2$
for $\varepsilon\sim\mathcal N(0,I_d)$), exactly cancels this scale
variation, and reappears, generalised to $SO(3)$, as the
\pyfunc{score\_scaling} construction reconstructed precisely in
\Cref{sec:diff-so3-numerics}.

\subsection{The reverse-time SDE}

\Cref{eq:forward-sde-scratch} destroys structure (runs from data to noise);
generation needs the reverse. The following classical theorem is what
licenses training a score network and then \emph{simulating backwards} to
generate samples.

\begin{theorem}[Anderson's reverse-time SDE]
\label{thm:anderson}
Let $(Y_s)_{s\in[0,1]}$ solve \eqref{eq:forward-sde-scratch} with marginal
laws $q_s=\mathrm{Law}(Y_s)$. Under mild regularity conditions (existence
and smoothness of $q_s$ for all $s$, and suitable growth conditions on
$f,g$; see Anderson, \emph{Reverse-Time Diffusion Equation Models},
Stochastic Processes and their Applications, 1982), the time-reversed
process $\widetilde Y_s := Y_{1-s}$ solves the SDE
\[
d\widetilde Y_s
=
\Bigl[-f(\widetilde Y_s,1-s) + g(1-s)^2\,\nabla_x\log q_{1-s}(\widetilde Y_s)\Bigr]\,ds
+
g(1-s)\,d\overline W_s ,
\]
where $\overline W$ is a Wiener process for the reversed filtration (a
technical point about which $\sigma$-algebra the increments are independent
of, which does not affect how the SDE is simulated and which we do not
elaborate on further).
\end{theorem}

We do not reprove \Cref{thm:anderson} (it is a statement about how the
Fokker--Planck equation governing $q_s$'s time evolution transforms under
time reversal, requiring It\^o calculus we have not developed; see the
cited reference, or Haussmann \& Pardoux, \emph{Time Reversal of Diffusions},
Annals of Probability, 1986, for the technically complete statement). We
translate it into our generation-time convention, $t=1-s$: writing
$X_t:=\widetilde Y_{1-t}=Y_t$ (so $p_t=q_t$, matching our earlier
notation), $X_0=Y_1\sim p_{\mathrm{init}}$ (approximately, the fully-noised
end), $X_1=Y_0\sim p_{\mathrm{data}}$, and, substituting $s=1-t$,
$ds=-dt$, into \Cref{thm:anderson}'s SDE and flipping the direction of
integration accordingly,
\begin{equation}
\label{eq:reverse-sde-scratch}
dX_t
=
\Bigl[-f(X_t,1-t) + g(1-t)^2\nabla_x\log p_t(X_t)\Bigr]\,dt
+
g(1-t)\,d\overline W_t,
\qquad
X_0\sim p_{\mathrm{init}} ,
\end{equation}
integrated \emph{forward} in $t$ from $0$ to $1$. This is the SDE that
generation simulates, and it requires exactly one previously-undetermined
quantity: the score $\nabla_x\log p_t(x)$, precisely what
\eqref{eq:dsm-continuous} trains a network to approximate. Substituting
$s_\theta\approx\nabla_x\log p_t$ into \eqref{eq:reverse-sde-scratch} and
discretising by Euler--Maruyama (the direct generalisation of
\eqref{eq:ula}'s update rule to a state- and time-dependent drift/diffusion)
gives the sampling algorithm; \Cref{chap:sigmadock} exhibits the exact
discretisation SigmaDock's code uses, which matches this recipe precisely
for the translational (VP) component.

\section{The probability flow ODE, and the bridge to flow matching}
\label{sec:pf-ode-scratch}

\begin{theorem}[Probability flow ODE]
\label{thm:pf-ode-scratch}
Under the same conditions as \Cref{thm:anderson}, the deterministic ODE
\begin{equation}
\label{eq:pf-ode-scratch}
\frac{dX_t}{dt} = u_t(X_t) := -f(X_t,1-t) + \tfrac12 g(1-t)^2\,\nabla_x\log p_t(X_t)
\end{equation}
has the \emph{same} marginal laws $(p_t)_{t\in[0,1]}$ as
\eqref{eq:reverse-sde-scratch} (and hence as the original forward
process \eqref{eq:forward-sde-scratch}, relabelled).
\end{theorem}

\begin{proof}[Proof sketch]
Both \eqref{eq:reverse-sde-scratch} and \eqref{eq:pf-ode-scratch} are
constructed so that their associated Fokker--Planck / continuity equations
(the PDE governing how a density evolves under the dynamics; for an SDE
$dX=b(X,t)dt+g(t)dW$, this is
$\partial_tp_t = -\nabla\cdot(b\,p_t) + \tfrac12g(t)^2\Delta p_t$, and for
an ODE $dX/dt=u_t(X)$ it is the driftless-diffusion special case
$\partial_tp_t=-\nabla\cdot(u_tp_t)$, which is exactly
\Cref{thm:continuity} of \Cref{chap:flow-matching}, proved there in full)
coincide: substituting $b=-f(\cdot,1-t)+g(1-t)^2\nabla\log p_t$ into the
SDE's Fokker--Planck equation and using
$\nabla\cdot(p_t\nabla\log p_t)=\nabla\cdot(\nabla p_t)=\Delta p_t$ shows
the diffusion term $\tfrac12g(1-t)^2\Delta p_t$ exactly cancels against
half of the drift's contribution, leaving the ODE's continuity equation
with $u_t$ as in \eqref{eq:pf-ode-scratch}. We do not carry out this
cancellation in full generality here (it requires the Fokker--Planck
equation, which \Cref{chap:flow-matching} develops in the driftless case
needed there), but note it is a standard, mechanical computation once
\Cref{thm:continuity} is available --- see, e.g., Song et al.\ (2021,
Appendix D.1) for the complete derivation in this SDE setting.
\end{proof}

\Cref{thm:pf-ode-scratch} is exactly the content of
\Cref{sec:pf-ode-scratch}'s claim, made precise: \emph{every} score-based diffusion
model secretly defines a deterministic flow with the same marginals, whose
velocity field \eqref{eq:pf-ode-scratch} is an explicit, invertible
function of the score. \Cref{chap:flow-matching} develops the
flow-matching framework as though starting from scratch, and
\Cref{sec:pf-ode-scratch} (there) proves the general form of this correspondence;
what \Cref{thm:pf-ode-scratch} establishes here is only the specific fact,
needed to state \Cref{chap:sigmadock}'s Remark on ``items requiring
verification'', that SigmaDock's own sampler \emph{could} in principle be
run as a deterministic ODE (setting the SDE's injected noise to zero) with
unchanged marginals --- and \Cref{chap:sigmadock} confirms, directly from
the code, that this is in fact the \emph{default} configuration SigmaDock
ships with (\file{conf/sampling/base.yaml}'s default \texttt{noise\_scale:0.0}),
narrowing the ``stochastic vs.\ deterministic'' distinction between the two
methods considerably.

\clearpage

% ================================================================
\chapter{Diffusion Models on $SO(3)$}
\label{chap:diffusion-so3}

\Cref{chap:diffusion} developed diffusion entirely in $\RR^d$; by
\Cref{prop:se3-product}, this already suffices for the \emph{translational}
component of a pose. This chapter develops the rotational analogue: what
does ``add Gaussian noise'' mean on $SO(3)$, where addition is undefined,
and what replaces the closed-form Gaussian kernel of
\Cref{def:gaussian-kernel}?

\section{Brownian motion on a Riemannian manifold}
\label{sec:manifold-brownian}

\paragraph{Motivation.} In $\RR^d$, Brownian motion can be built either as
the limit of a random walk with i.i.d.\ Gaussian increments, or as the
unique (in law) continuous process with independent Gaussian increments of
variance proportional to elapsed time. On a curved manifold, ``increment''
must be reinterpreted: an increment is a step in a \emph{tangent space},
which must then be mapped back onto the manifold.

\begin{definition}[Manifold Brownian motion, via a random-walk limit]
\label{def:manifold-brownian}
Fix $x_0\in\mathcal M$ and a step size $\Delta\tau>0$. Define a discrete
process $x_0,x_{\Delta\tau},x_{2\Delta\tau},\dots$ by
\[
x_{(k+1)\Delta\tau} = \exp_{x_{k\Delta\tau}}\bigl(\sqrt{\Delta\tau}\,\xi_k\bigr),
\qquad
\xi_k\overset{\mathrm{iid}}\sim\mathcal N\bigl(0,I_{\dim\mathcal M}\bigr)\ \text{in } T_{x_{k\Delta\tau}}\mathcal M
\]
(a Gaussian tangent vector sampled in the tangent space at the current
point, using an orthonormal-basis identification $T_x\mathcal M\cong\RR^{\dim\mathcal M}$
via the metric $g_x$, then pushed back onto $\mathcal M$ by the exponential
map). As $\Delta\tau\to0$, this process converges (in an appropriate sense;
we do not make this precise, taking it as the operational definition
throughout this document, following the treatment standard in the
geometric-diffusion-model literature, e.g.\ Leach et al., \emph{Denoising
Diffusion Probabilistic Models on SO(3) for Rotational Alignment}, ICLR
Workshop 2022, and De Bortoli et al., \emph{Riemannian Score-Based
Generative Modelling}, NeurIPS 2022) to a continuous-time process
$B^{\mathcal M}_\tau$, \emph{Brownian motion on $\mathcal M$}.
\end{definition}

\begin{codebox}
\textbf{Mathematical object:} one step of $B^{SO(3)}_\tau$: sample a
Gaussian tangent vector and push it onto the manifold \\
\textbf{Implementation:} \file{so3\_utils.py}::\pyfunc{tangent\_gaussian}
(samples $\widehat\varepsilon$, $\varepsilon\sim\mathcal N(0,I_3)$, and
left-multiplies by $R$: returns $R\widehat\varepsilon$, a
right-trivialised --- \Cref{sec:so3-trivialisation}, note the tangent
vector's own components are drawn in the \emph{same} $\RR^3$ regardless of
$R$, then transported by right-multiplication, i.e.\ this is exactly
\Cref{prop:euler-right-triv}'s convention --- Gaussian tangent vector at
$R$) composed with \file{so3\_utils.py}::\pyfunc{expmap} (the map
$\exp_R$) \\
\textbf{Tensor shape:} \shape{..., 3, 3} $\to$ \shape{..., 3, 3} \\
\textbf{Interpretation:} \pyfunc{expmap}, per \Cref{sec:rodrigues-scratch}'s
code correspondence, computes $\exp_R(v)$ via \texttt{torch.matrix\_exp}
after re-symmetrising the skew part for numerical stability, rather than
Rodrigues' closed form; this is the code-level realisation of one step of
\Cref{def:manifold-brownian}'s random walk, used directly inside
SigmaDock's reverse-time sampler (\Cref{sec:diff-so3-numerics}).
\end{codebox}

\begin{proposition}[Stationary distribution of $SO(3)$ Brownian motion]
\label{prop:so3-brownian-stationary}
$B^{SO(3)}_\tau$'s law converges, as $\tau\to\infty$, to the Haar (uniform)
measure of \Cref{prop:haar-so3}, independently of the starting point $x_0$.
\end{proposition}

\begin{proof}[Proof sketch]
Brownian motion on a compact Riemannian manifold is reversible with respect
to, and converges to, the Riemannian volume measure --- the manifold
analogue of \Cref{thm:langevin}'s statement that $\RR^d$ Langevin dynamics
converges to $p\propto e^{-U}$, specialised to the driftless case $U\equiv0$
(constant potential, i.e.\ pure diffusion with no preferred point), whose
stationary measure is simply the volume measure itself; on $SO(3)$ this
volume measure is (proportional to, hence equal to after normalisation)
the Haar measure, since Haar measure \emph{is} the (unique, up to
normalisation) bi-invariant, and in particular the Riemannian-volume-induced,
measure. We do not reprove the general convergence statement (it again
requires the manifold Fokker--Planck / heat equation, whose generator is
the Laplace--Beltrami operator, \Cref{def:laplace-beltrami} below), but note
it is exactly why $p_{\mathrm{init}}=\mathcal U(SO(3))$, not some other
distribution, is the correct rotational prior for a diffusion model whose
forward process runs Brownian motion for a long enough time
--- \Cref{sec:sigmadock-forward} confirms this is precisely what the code
uses.
\end{proof}

This is the promised structural point from \Cref{sec:manifolds-scratch}:
unlike $\RR^d$ Brownian motion, whose stationary behaviour is
``variance $\to\infty$'', a \emph{compact} manifold's Brownian motion has a
genuine equilibrium distribution determined entirely by the manifold's
geometry, and $SO(3)$'s compactness (it is a closed, bounded subset of
$\RR^{3\times3}$) is what makes the variance-exploding schedule of
\Cref{ex:ve-sde} converge to something useful rather than literally
exploding without bound, as it would in $\RR^d$.

\begin{definition}[Laplace--Beltrami operator, informal]
\label{def:laplace-beltrami}
The \emph{Laplace--Beltrami operator} $\Delta_{\mathcal M}$ is the natural
generalisation of the Euclidean Laplacian $\Delta=\sum_i\partial_i^2$ to a
Riemannian manifold, defined so that it is the infinitesimal generator of
Brownian motion on $\mathcal M$ (\Cref{def:manifold-brownian}) in exactly
the sense that $\Delta$ is the generator of ordinary $\RR^d$ Brownian
motion: the transition density $k_\tau(x,y)$ of $B^{\mathcal M}_\tau$
(started at $x$) solves the \emph{heat equation}
$\partial_\tau k_\tau = \Delta_{\mathcal M}k_\tau$, exactly as the ordinary
Gaussian kernel solves $\partial_\tau k_\tau=\Delta k_\tau$ in $\RR^d$. We
do not need $\Delta_{\mathcal M}$'s coordinate formula anywhere in this
document; only the existence and role of the heat kernel it generates,
constructed explicitly for $SO(3)$ in \Cref{sec:igso3-derivation} below.
\end{definition}

\section{The heat kernel on $SO(3)$: $\mathrm{IGSO}(3)$}
\label{sec:igso3-derivation}

\paragraph{Motivation.} \Cref{def:gaussian-kernel}'s Gaussian perturbation
kernel $p_\sigma(x\mid z)=\mathcal N(x\mid z,\sigma^2I)$ is, in
$\RR^d$-specific language, exactly the \emph{heat kernel}: the transition
density of Brownian motion after time $\sigma^2$. We now want its
$SO(3)$-analogue: the transition density $k_{\sigma^2}(R_0,R)$ of
$SO(3)$-Brownian motion (\Cref{def:manifold-brownian}) run for
``time'' $\sigma^2$ starting at $R_0$.

\begin{theorem}[The isotropic Gaussian distribution on $SO(3)$]
\label{thm:igso3-scratch}
The heat kernel on $SO(3)$, called $\mathrm{IGSO}(3)(R_0,\sigma^2)$, depends
on $R_0,R$ only through the relative-rotation angle
$\omega=\Omega(R_0^\top R)$ of \Cref{thm:so3-log-scratch} (bi-invariance of
the metric, hence of the Laplace--Beltrami operator and its heat kernel,
under simultaneous left/right translation), and its density with respect
to the Haar measure admits the spectral (Peter--Weyl) expansion
\begin{equation}
\label{eq:igso3-scratch}
f_\sigma(\omega)
=
\sum_{\ell=0}^{\infty}(2\ell+1)\,\chi_\ell(\omega)\,\exp\!\Bigl(-\tfrac12\ell(\ell+1)\sigma^2\Bigr),
\qquad
\chi_\ell(\omega) = \frac{\sin\bigl((\ell+\tfrac12)\omega\bigr)}{\sin(\omega/2)} ,
\end{equation}
where $\chi_\ell$ is the character of the degree-$\ell$ irrep of $SO(3)$
(\Cref{thm:so3-irrep-classification}).
\end{theorem}

\begin{proof}[Proof sketch]
This is the Peter--Weyl expansion of the heat kernel of a compact Lie
group: the heat kernel, viewed as a function on the group (using
bi-invariance to reduce to a function of one variable via
$k_{\sigma^2}(R_0,R)=k_{\sigma^2}(I,R_0^\top R)=:f_\sigma(\omega)$), decomposes
in the orthogonal basis of \emph{matrix coefficients} of the irreducible
representations $D^\ell$ (the Peter--Weyl theorem: $L^2(G)$ for a compact
group $G$ decomposes as a direct sum, over irreps, of matrix-coefficient
spaces), each an eigenfunction of the Laplace--Beltrami operator with
eigenvalue $-\ell(\ell+1)$ (the Casimir eigenvalue of the degree-$\ell$
irrep of $SO(3)$, a standard fact from the representation theory cited
after \Cref{thm:so3-irrep-classification}), and the coefficient
$e^{-\ell(\ell+1)\sigma^2/2}$ is exactly the corresponding solution
$e^{\lambda\tau}$ of the heat equation's eigenvalue ODE with
$\lambda=-\ell(\ell+1)/2$, $\tau=\sigma^2$; summing the trace of $D^\ell$
over the class-function structure of a bi-invariant kernel (i.e.\ summing
each matrix-coefficient block's diagonal, which is exactly the character
$\chi_\ell$) and using the classical closed form for $SO(3)$'s degree-$\ell$
character, $\chi_\ell(\omega)=\sin((\ell+\tfrac12)\omega)/\sin(\omega/2)$
(itself derivable from the eigenvalues of $D^\ell(\exp(\theta\widehat n))$,
which are $\{e^{im\theta}:m=-\ell,\dots,\ell\}$ in a suitable complex basis,
summing as a geometric series), gives \eqref{eq:igso3-scratch}. We take the
Peter--Weyl decomposition and the Casimir eigenvalue computation as
classical facts and do not reprove them (see Chirikjian, 2012, cited
earlier, \S10.3, for the complete derivation).
\end{proof}

\begin{remark}[The series must be truncated in practice]
\label{rem:igso3-truncation}
\Cref{eq:igso3-scratch} is an infinite sum; \Cref{sec:diff-so3-numerics}
below records exactly how the actual SigmaDock implementation truncates,
tabulates, and interpolates it, since this is a genuine engineering
component of the method, not an afterthought --- and, as
\Cref{sec:fm-so3-scratch} in \Cref{chap:flow-matching} shows, it is
precisely this machinery that flow matching eliminates entirely.
\end{remark}

\begin{proposition}[The angle marginal density]
\label{prop:igso3-angle-marginal}
The density of $\omega=\Omega(R_0^\top R)$ alone (as opposed to the full
density on $SO(3)$) is $\tfrac{1-\cos\omega}\pi\,f_\sigma(\omega)$: the
IGSO(3) density $f_\sigma$ times the same Haar-measure Jacobian
$1-\cos\omega$ that appeared in \Cref{prop:haar-so3} (up to the same
normalisation constant), since $\mathrm{IGSO}(3)(R_0,\sigma^2)$, in the
limit $\sigma\to\infty$, must reduce to the Haar measure of
\Cref{prop:haar-so3} by \Cref{prop:so3-brownian-stationary}, and this is
consistent with \eqref{eq:igso3-scratch} since every term with $\ell\ge1$
decays to $0$ as $\sigma\to\infty$, leaving only the $\ell=0$ term
$f_\infty(\omega)=1\cdot\chi_0(\omega)\cdot1 = 1$ (using
$\chi_0(\omega)=\sin(\omega/2)/\sin(\omega/2)=1$), recovering exactly
\Cref{prop:haar-so3}'s density $\propto1-\cos\omega$.
\end{proposition}

\subsection{The score of the $\mathrm{IGSO}(3)$ kernel}
\label{sec:igso3-score-derivation}

\begin{proposition}[Riemannian score of $\mathrm{IGSO}(3)$]
\label{prop:igso3-score-scratch}
Under the left-trivialised (world-frame) identification of
\Cref{def:trivialisation-scratch} at $R$, the score of
$p_{\sigma}(R\mid R_0):=\mathrm{IGSO}(3)(R_0,\sigma^2)$-density at $R$, as a
tangent vector in $T_RSO(3)\cong\RR^3$, is
\begin{equation}
\label{eq:igso3-score-scratch}
\nabla_R\log p_\sigma(R\mid R_0)
=
\frac{v}{\|v\|}\cdot\frac{\partial}{\partial\omega}\log f_\sigma(\omega)\Big|_{\omega=\|v\|},
\qquad
v = \log(R_0^\top R)^\vee \in\RR^3 .
\end{equation}
\end{proposition}

\begin{proof}[Proof sketch]
Since (\Cref{thm:igso3-scratch}) $p_\sigma(R\mid R_0)=f_\sigma(\omega)$
depends on $R$ only through the scalar $\omega=\Omega(R_0^\top R)$, the
chain rule gives $\nabla_R\log p_\sigma = f_\sigma'(\omega)/f_\sigma(\omega)\cdot\nabla_R\omega$
(the Riemannian gradient of the composite scalar function). The gradient of
the angle function $\omega(R)=\Omega(R_0^\top R)$, as a function on the
manifold, points in the direction of the tangent vector that most rapidly
increases $\omega$ --- exactly the unit tangent vector along the geodesic
away from $R_0$, which by \Cref{prop:so3-distance} is
$v/\|v\|=\log(R_0^\top R)^\vee/\omega$ --- with magnitude $1$ (since $\omega$
itself is arc length along that geodesic, by
\Cref{prop:so3-distance} again, and the gradient of arc length along its
own geodesic has unit norm, the direct manifold analogue of
$\nabla_x\|x\|=x/\|x\|$ in $\RR^d$). Combining gives
\eqref{eq:igso3-score-scratch}.
\end{proof}

\begin{codebox}
\textbf{Mathematical object:} $\nabla_R\log p_\sigma(R\mid R_0)$ \\
\textbf{Implementation:} \file{so3\_diffuser.py}::\pyfunc{SO3Diffuser.score} \\
\textbf{Tensor shape:} \shape{..., 3, 3} (world-frame element of $\so(3)$,
i.e.\ the code returns $R\cdot\widehat v\cdot(\partial_\omega\log f_\sigma/\omega)$
rather than the bare $\RR^3$ vector, matching the world-frame
identification of \Cref{def:trivialisation-scratch}) \\
\textbf{Interpretation:} $\partial_\omega\log f_\sigma$ is \emph{not}
recomputed by differentiating the tabulated density table at score-evaluation
time; it is evaluated \emph{live} from the analytic term-by-term derivative
of the truncated series \eqref{eq:igso3-scratch} (function
\pyfunc{d\_log\_f\_d\_omega}, \Cref{sec:diff-so3-numerics}), for the exact,
continuous $\sigma(t)$ of the query, not a value snapped to the
precomputed grid.
\end{codebox}

\subsection{Numerical implementation: truncation, tabulation, and caching}
\label{sec:diff-so3-numerics}

We now record, precisely and only from the code (not reconstructed
analytically), how SigmaDock evaluates the objects of the preceding two
subsections in practice, since \Cref{rem:igso3-truncation} promised this
level of detail and it is exactly the machinery \Cref{sec:fm-so3-scratch}
will show flow matching does not need.

\begin{codebox}
\textbf{Mathematical object:} the truncated series
$f_\sigma(\omega)=\sum_{\ell=0}^{L-1}(2\ell+1)\chi_\ell(\omega)e^{-\ell(\ell+1)\sigma^2/2}$
of \Cref{thm:igso3-scratch}, and its analytic derivative
$\partial_\omega f_\sigma$, $\partial_\omega\log f_\sigma$ \\
\textbf{Implementation:} \file{so3\_diffuser.py}, module-level functions
\pyfunc{igso3\_expansion} (the truncated sum, default $L=1000$ terms, i.e.\
degrees $\ell=0,\dots,999$), \pyfunc{igso3\_density} (multiplies in the
angle-marginal Jacobian $1-\cos\omega$ of
\Cref{prop:igso3-angle-marginal}), \pyfunc{d\_f\_igso3\_d\_omega} (the
term-by-term derivative, via the quotient rule applied to each
$\chi_\ell(\omega)=\sin((\ell+\tfrac12)\omega)/\sin(\omega/2)$),
\pyfunc{d\_log\_f\_d\_omega} \\
\textbf{Interpretation:} \pyfunc{d\_log\_f\_d\_omega} additionally
hard-branches (via boolean masking, not the soft \pyfunc{isclose}-blend
pattern of \Cref{sec:so3-log-scratch}) between the literal ratio
$\partial_\omega f/f$ and the small-noise approximation $-\omega/\sigma^2$
whenever $\sigma^2<0.3$, since the literal ratio becomes numerically
ill-conditioned at small $\sigma$ (both numerator and denominator of the
truncated series become extremely small simultaneously). We have not
independently re-derived that $-\omega/\sigma^2$ is the correct
$\sigma\to0$ limit of $\partial_\omega\log f_\sigma$ from
\eqref{eq:igso3-scratch} (it is plausible --- as $\sigma\to0$,
$\mathrm{IGSO}(3)(R_0,\sigma^2)$ should concentrate at $R_0$ and locally
resemble a flat Gaussian, whose log-density derivative in the radial
direction is exactly $-\omega/\sigma^2$ by \Cref{rem:gaussian-score}'s
formula specialised to one radial coordinate --- but we record this as the
code's own documented claim, flagged per this document's stated policy on
uncertainty, not as a result independently verified here).
\end{codebox}

\begin{codebox}
\textbf{Mathematical object:} the full apparatus of caching, tabulation,
and inverse-CDF sampling used at both training and inference time \\
\textbf{Implementation:} \file{so3\_diffuser.py}::\pyfunc{SO3Diffuser.\_\_init\_\_},
\pyfunc{sample\_igso3\_angle}, \pyfunc{sample}, \pyfunc{sample\_ref} \\
\textbf{Interpretation:} at construction time, the diffuser builds a
$1000\times2000$ grid: $1000$ noise levels $\sigma_i$ (linearly spaced in
$t\in[0,1]$, then mapped through the schedule
$\sigma(t)=\log\bigl(te^{\sigma_{\max}}+(1-t)e^{\sigma_{\min}}\bigr)$, the
\emph{logarithmic} schedule) $\times$ $2000$ angles $\omega_j$ (linearly
spaced in $(0,\pi]$, deliberately excluding $\omega=0$ to avoid the
$\sin(\omega/2)=0$ singularity visible in \eqref{eq:igso3-scratch}), and
evaluates $f_{\sigma_i}(\omega_j)$ (via \pyfunc{igso3\_density}, with
\texttt{marginal=True}, i.e.\ including the Jacobian of
\Cref{prop:igso3-angle-marginal}) and
$\partial_{\omega}\log f_{\sigma_i}(\omega_j)$ (via
\pyfunc{d\_log\_f\_d\_omega}) at every grid point, once, and caches both
tables to disk, together with a row-wise cumulative-sum CDF table (a
Riemann-sum numerical integral of the density table, safe-normalised by
its own final value clamped away from zero). The cache directory name
encodes every hyperparameter of the grid (matching the
\texttt{eps\_1000\_omega\_2000\_min\_sigma\_...\_schedule\_logarithmic}-style
paths visible in this project's own training logs). \emph{Sampling} an
angle at a queried $t$ (\pyfunc{sample\_igso3\_angle}) looks up the nearest
tabulated $\sigma_i$ row via \pyfunc{torch.bucketize}, draws
$u\sim\mathrm{Unif}(0,1)$, locates $u$'s bracketing interval in that row's
cached CDF, and \emph{linearly interpolates} between the two bracketing
grid angles --- piecewise-linear inverse-CDF sampling on a fixed
$2000$-point grid, not exact analytic inversion. \pyfunc{sample\_ref} (the
prior sampler, i.e.\ $R\sim p_{\mathrm{init}}$) calls \pyfunc{sample} at
$t=1$, the schedule's maximum-noise end --- an \emph{approximation} to the
exact Haar-uniform prior of \Cref{prop:so3-brownian-stationary}, exact only
in the limit $\sigma_{\max}\to\infty$, not literally the closed-form Haar
sampler of \Cref{prop:haar-so3} (that routine, \pyfunc{sample\_uniform}, is
a separate function, used instead by \emph{SigmaFlow}'s prior, as recorded
in \Cref{sec:fm-so3-scratch}). By contrast, the \emph{score itself}
(\Cref{prop:igso3-score-scratch}'s code correspondence, immediately above)
is \textbf{not} read from this cached table at evaluation time --- it is
recomputed live via \pyfunc{d\_log\_f\_d\_omega} at the query's exact,
continuous $(\omega,\sigma)$; the cached \texttt{score\_norms} table exists
only to build the training loss-weight $\lambda(t)$
(\pyfunc{score\_scaling}, matching \eqref{eq:dsm-continuous}'s $\lambda(s)$
in this document's notation), computed once at construction time as a
weighted second moment,
$\lambda_i^{-1/2}=\sqrt{\EE_\omega[(\partial_\omega\log f_{\sigma_i})^2]/3}$
(the division by $3$ normalising a full-vector squared-norm expectation
down to a per-coordinate one), over the angle marginal at noise level
$\sigma_i$.
\end{codebox}

\section{Riemannian denoising score matching and the reverse-time SDE on $SO(3)$}
\label{sec:diff-so3}

With \Cref{prop:igso3-score-scratch} in hand, the entire machinery of
\Cref{chap:diffusion} transfers to $SO(3)$ by the single substitution
``Euclidean gradient $\to$ Riemannian gradient'', exactly as
\Cref{def:metric-scratch}'s generalisation of $\RR^d$ calculus to
manifolds promised at the start of \Cref{chap:groups}.

\begin{theorem}[Riemannian denoising score matching]
\label{thm:rdsm-scratch}
Define, analogously to \Cref{thm:dsm-equivalence},
\[
\mathcal L_{\mathrm{RDSM}}(\theta)
=
\EE_{t,\,R_1\sim p_{\mathrm{data}},\,R_t\sim p_t(\cdot\mid R_1)}
\Bigl[\lambda_t\bigl\|s_\theta(R_t,t) - \nabla_{R_t}\log p_t(R_t\mid R_1)\bigr\|_g^2\Bigr],
\]
with $p_t(\cdot\mid R_1) = \mathrm{IGSO}(3)(R_1,\sigma(1-t)^2)$
(\Cref{rem:time-convention-global}'s convention: $R_1$ is data, so the
conditioning is on the $t=1$ endpoint, and the noise level increases as $t$
decreases, i.e.\ as we move away from data towards $t=0$). Then, exactly as
in \Cref{thm:dsm-equivalence}'s proof (which uses nothing about the ambient
space beyond the existence of the gradient/score and the marginalisation
identity of \Cref{prop:marginalisation}, both of which hold verbatim on a
Riemannian manifold with $\nabla_x$ replaced by the Riemannian gradient and
$dx$ by the Riemannian volume element), minimising
$\mathcal L_{\mathrm{RDSM}}$ has the same minimiser,
$s_\theta(R,t)\approx\nabla_R\log p_t(R)$, as the (intractable) marginal
Riemannian score matching loss.
\end{theorem}

We do not repeat the proof (it is, verbatim, the computation of
\Cref{thm:dsm-equivalence}, with every Euclidean integral replaced by an
integral against the Riemannian volume measure, valid since none of the
manipulations used --- integration by parts under the gradient, Fubini,
linearity of expectation --- depend on the ambient space being flat).

Sampling then simulates the manifold reverse-time SDE, the direct
$SO(3)$-analogue of \eqref{eq:reverse-sde-scratch}:
\begin{equation}
\label{eq:reverse-so3-scratch}
dR_t = g(1-t)^2\,\nabla_R\log p_t(R_t)\,dt + g(1-t)\,d\overline B_t^{SO(3)},
\qquad R_0\sim p_{\mathrm{init}} ,
\end{equation}
(no drift term $f$, since \Cref{ex:ve-sde}'s VE-type schedule, which
SigmaDock's rotational component uses per the code correspondence of
\Cref{sec:diff-so3-numerics}, has $f\equiv0$), discretised as a
\emph{geodesic random walk}: at each step, take one
\Cref{def:manifold-brownian}-style Brownian increment, scaled and drifted
by the (learned, substituted) score, and map it onto $SO(3)$ via
$\exp_{R_t}$ --- guaranteeing, by construction, that the discretised
trajectory never leaves $SO(3)$, for any step size, unlike a naive
Euclidean discretisation of a constrained system.

\begin{codebox}
\textbf{Mathematical object:} one Euler--Maruyama step of
\eqref{eq:reverse-so3-scratch} \\
\textbf{Implementation:} \file{so3\_diffuser.py}::\pyfunc{SO3Diffuser.reverse} \\
\textbf{Tensor shape:} \shape{..., 3, 3} $\to$ \shape{..., 3, 3} \\
\textbf{Interpretation:} computes
\texttt{perturb}$= g_t^2\cdot\text{score}_t\cdot dt + \text{noise\_scale}\cdot g_t\cdot\sqrt{dt}\cdot z$
with $z=$\pyfunc{tangent\_gaussian}$(R_t)$ (\Cref{sec:manifold-brownian}'s
code correspondence) and $g_t=$\pyfunc{diffusion\_coef}$(t)$, then returns
$R_{t-dt}=$\pyfunc{expmap}$(R_t,\text{perturb})$ --- exactly
\eqref{eq:reverse-so3-scratch}'s Euler--Maruyama discretisation, with the
manifold-respecting update $\exp_{R_t}(\cdot)$ in place of the naive
Euclidean $R_t+(\cdot)$. An optional \texttt{noise\_scale} parameter
(default, per \file{conf/sampling/base.yaml}, exactly $0.0$ --- see
\Cref{chap:sigmadock}) rescales the injected noise term $z$ independently
of the score-drift term, allowing the sampler to be run anywhere between
fully stochastic (\texttt{noise\_scale=1}) and fully deterministic
(\texttt{noise\_scale=0}, in which case \eqref{eq:reverse-so3-scratch}'s
SDE degenerates to \Cref{thm:pf-ode-scratch}'s probability flow ODE,
specialised to $SO(3)$ by the same cancellation argument, restated but not
re-derived here).
\end{codebox}

This completes the diffusion side of the theory. \Cref{chap:sigmadock}
assembles \Cref{chap:diffusion,chap:diffusion-so3} into the full SigmaDock
model --- combining the VP translational process of \Cref{ex:vp-sde} and
the VE rotational process of this chapter via the product-manifold
structure of \Cref{prop:se3-product}, and connecting every formula derived
here to the exact network architecture (\Cref{chap:equiformer}) that
predicts the required scores.

\clearpage

% ================================================================
\part{Architecture}
\label{part:architecture}
% ================================================================

\chapter{From Sets to Geometry: Attention, Message Passing, and the Equivariance Obstruction}
\label{chap:pre-equiformer}

\Cref{chap:diffusion,chap:diffusion-so3} reduced ``generate a pose'' to
``learn a score/vector-field network'', but said nothing about what that
network's \emph{input} looks like or how it must be built. This chapter
develops that architecture from the ground up, assuming only that the
reader knows what a matrix multiplication and a feed-forward layer are; by
its end we will have derived, not merely stated, why an architecture built
from scalars and distances alone \emph{cannot} do the job, which is exactly
the gap \Cref{chap:equiformer} closes.

\section{The network's signature}
\label{sec:network-signature-scratch}

Consolidating \Cref{chap:diffusion,chap:diffusion-so3}: the object to be
learned, whether called a score ($s_\theta$) or (in
\Cref{chap:flow-matching}) a velocity ($u_\theta$), is a function
\begin{equation}
\label{eq:network-signature-scratch}
u_\theta:\ \mathcal M_N\times[0,1]\ \longrightarrow\ T\mathcal M_N,
\qquad
u_\theta(\mathbf T,t) = \bigl(\omega^{(i)}_\theta,\,v^{(i)}_\theta\bigr)_{i=1}^N
\in(\RR^3\times\RR^3)^N ,
\end{equation}
where, by the left/right-trivialisation identification fixed in
\Cref{sec:so3-trivialisation}, $\omega^{(i)}_\theta$ is fragment $i$'s
rotational output and $v^{(i)}_\theta$ its translational output, each an
element of $\RR^3$ per fragment (never a full $3\times3$ matrix; the
tangent-space identification of \Cref{chap:groups} is precisely what
collapses ``an infinitesimal rotation'' to ``three numbers''). The input
$\mathbf T\in\mathcal M_N$ is not, however, presented to the network in
this abstract group-element form; it is presented as the concrete Cartesian
coordinates $\mathbf x=\varphi(\mathbf T)\in\RR^{n\times3}$ of every atom
(protein and ligand) obtained by applying each fragment's current rigid
motion to its reference-conformer template
(\Cref{sec:fragment-reconstruction} makes this map $\varphi$ precise), plus
each atom's discrete chemical identity, together with $t$. The network,
therefore, is a map from an \emph{unordered, geometric} collection of
labelled points in $\RR^3$ to a per-fragment vector-pair output, and the
rest of this chapter and the next derive what such a map is allowed to look
like.

\section{Two structural requirements}
\label{sec:two-requirements}

\begin{description}
\item[Permutation equivariance.] Atoms in $\mathbf x$ carry no canonical
order; they are a \emph{set} (indexed arbitrarily for storage), not a
sequence. Reindexing the input must permute the corresponding output
identically and change nothing else.
\item[$SE(3)$-equivariance.] The physical process being modelled (a ligand
binding a pocket) does not depend on the pose of the laboratory reference
frame. Applying a rigid motion to the entire input (protein and ligand
together) must transform the predicted vector field exactly as
\Cref{def:equivariance-full} below makes precise --- \emph{not} leave it
unchanged, and \emph{not} change it arbitrarily.
\end{description}

\begin{definition}[Group action, invariance, equivariance]
\label{def:equivariance-full}
Let a group $G$ act on a set (or manifold) $\mathcal X$ via maps
$\Phi_Q:\mathcal X\to\mathcal X$, $Q\in G$, satisfying
$\Phi_{Q_1Q_2}=\Phi_{Q_1}\circ\Phi_{Q_2}$ (a group action, the natural
generalisation of a linear representation, \Cref{def:representation-scratch},
to a possibly nonlinear space). A function $\rho:\mathcal X\to\RR$ is
\emph{$G$-invariant} if $\rho(\Phi_Q(x))=\rho(x)$ for all $Q,x$. A function
$u:\mathcal X\to T\mathcal X$ (with $G$ also acting on the tangent bundle by
the differential $d\Phi_Q$) is \emph{$G$-equivariant} if
$u(\Phi_Q(x)) = d\Phi_Q|_x\,u(x)$ for all $Q,x$: transform-then-evaluate
equals evaluate-then-transform.
\end{definition}

For the permutation group $G=S_n$ acting on $\mathbf x\in\RR^{n\times3}$ by
reindexing, and for the rigid-motion group $SE(3)$ acting by
$\Phi_{(Q,b)}(\mathbf x)_i = Q\mathbf x_i+b$, \Cref{def:equivariance-full}
specialises to exactly the two requirements stated informally above; for
the rotational part of \eqref{eq:network-signature-scratch}, using the
world/body-frame bookkeeping fixed in \Cref{sec:so3-trivialisation}
(equivariance of the \emph{body-frame} rotational output under a
\emph{world-frame} rotation of the whole scene means the body-frame output
is left \emph{unchanged}, since it is expressed relative to the fragment's
own, co-rotating frame --- an important, easily-misstated point, discussed
again with the network's actual output convention in
\Cref{sec:equiformer-readout}).

\begin{proposition}[Why equivariance, not merely invariance, is the right constraint]
\label{prop:equivariant-invariant-scratch}
Let $G$ act on $\mathcal M$ by isometries (maps preserving the Riemannian
metric $g$; both the permutation action on a graph's node set and the
rigid-motion action on $\mathcal M_N$ satisfy this), let $p_0$ be
$G$-invariant, and let $u_\theta$ be $G$-equivariant. Then every marginal
$p_t^\theta$ of the flow $\tfrac{d}{dt}X_t=u_\theta(X_t,t)$, $X_0\sim p_0$
(\Cref{chap:flow-matching}; the analogous statement for the diffusion SDEs
of \Cref{chap:diffusion,chap:diffusion-so3} holds by the same argument,
stated precisely in \Cref{thm:sigmadock-symmetry}), is $G$-invariant; in
particular the generated data distribution $p_1^\theta$ is.
\end{proposition}

\begin{proof}
Let $Y_t := \Phi_Q(X_t)$. Since $\Phi_Q$ is a fixed diffeomorphism,
$\dot Y_t = d\Phi_Q|_{X_t}\dot X_t = d\Phi_Q|_{X_t}\,u_\theta(X_t,t)
= u_\theta(\Phi_Q(X_t),t) = u_\theta(Y_t,t)$, using equivariance of
$u_\theta$ in the last step: $Y_t$ solves the \emph{same} ODE as $X_t$.
Since $Y_0=\Phi_Q(X_0)$ has law $\Phi_{Q\#}p_0 = p_0$ ($G$-invariance of
$p_0$), and solutions of an ODE are uniquely determined by their initial
condition (under standard regularity, e.g.\ Lipschitz $u_\theta$; a
classical existence/uniqueness fact for ODEs we take as given), $Y_t$ and
$X_t$ have the \emph{same} law for every $t$. But $Y_t=\Phi_Q(X_t)$ by
definition, so $\mathrm{Law}(X_t) = \mathrm{Law}(\Phi_Q(X_t)) = \Phi_{Q\#}\mathrm{Law}(X_t)$:
$p_t^\theta:=\mathrm{Law}(X_t)$ is $G$-invariant.
\end{proof}

This is the reason equivariance is not a stylistic preference: it is a
\emph{hard, exact} guarantee --- for \emph{every} value of the learned
parameters $\theta$, at \emph{every} point during training, not only at
convergence --- that the generative model respects the symmetries of the
physical problem, obtained purely from an architectural constraint rather
than needing to be learned from data (data augmentation could
\emph{approximately} teach a non-equivariant network the same symmetry, at
the cost of a strictly larger effective hypothesis space to search and no
exact guarantee at any finite amount of training).

\section{Attention, and the permutation-equivariance half of the problem}
\label{sec:transformer-scratch}

\begin{definition}[Scaled dot-product self-attention; restated]
\label{def:attention-scratch}
For tokens $X=(x_1,\dots,x_n)^\top\in\RR^{n\times d}$ and projections
$W_Q,W_K\in\RR^{d\times d_k}$, $W_V\in\RR^{d\times d_v}$, set
$Q=XW_Q,K=XW_K,V=XW_V$ and
\[
\Attn(X) = AV,
\qquad
A = \softmax\!\Bigl(\frac{QK^\top}{\sqrt{d_k}}\Bigr)\in\RR^{n\times n}
\]
(row-wise softmax, so $A_{ij}\ge0$, $\sum_jA_{ij}=1$).
\end{definition}

Row $i$ of $\Attn(X)$ is a convex combination $\sum_jA_{ij}v_j$ of value
vectors, weighted by the learned compatibility $\langle q_i,k_j\rangle$
between tokens $i$ and $j$; the $1/\sqrt{d_k}$ scaling keeps logits at unit
scale for i.i.d.\ unit-variance $q,k$ entries, preventing early softmax
saturation. Multiple parallel \emph{heads} (independent
$W_Q,W_K,W_V$ triples, concatenated and mixed by a further $W_O$) and a
pre-normalised residual block,
\begin{equation}
\label{eq:transformer-block-scratch}
X' = X + \Attn(\LN(X)),
\qquad
X'' = X' + \MLP(\LN(X')) ,
\end{equation}
complete the ordinary Transformer layer.

\begin{proposition}[Permutation equivariance of attention]
\label{prop:perm-equivariance-scratch}
For a permutation matrix $P\in\RR^{n\times n}$: $\Attn(PX)=P\Attn(X)$, and
\eqref{eq:transformer-block-scratch} inherits this property.
\end{proposition}

\begin{proof}
$Q,K,V\mapsto PQ,PK,PV$ (the projections act on the feature axis, so
permuting rows of $X$ permutes rows of $Q,K,V$ identically). Hence
$QK^\top\mapsto PQK^\top P^\top$; since row-wise softmax acts identically on
each row regardless of row order, $A\mapsto PAP^\top$. Then
$AV\mapsto PAP^\top PV = PAV$ using $P^\top P=I$. $\LN$ and $\MLP$ act
identically, tokenwise, on every row, hence commute with $P$.
\end{proof}

Attention thus solves permutation equivariance exactly \emph{because} it
treats its input as an (index-labelled but otherwise symmetric) set: the
formula \eqref{def:attention-scratch} never references the token index $i$
except through $x_i$ itself. Two further observations connect attention to
the constructions of the next section: first, \eqref{def:attention-scratch}
\emph{is} message passing (\Cref{def:mpnn-scratch} below) on the
\emph{complete} graph, with learned edge weight $A_{ij}$; second, its
$\mathcal O(n^2d)$ cost, from computing $QK^\top$ densely, is why in
practice attention (and the equivariant version developed below) is
restricted to a sparse graph --- a radius or $k$-nearest-neighbour graph ---
rather than the complete graph, recovering tractability while keeping the
learned, data-dependent weighting.

What \eqref{def:attention-scratch} does \emph{not} provide is any use of
geometry: $\langle q_i,k_j\rangle$ depends on the tokens' \emph{features},
never on where in $\RR^3$ the corresponding atoms sit. Concatenating raw
coordinates $x_i\in\RR^3$ into the feature vector would inject geometry but
destroy $SE(3)$-equivariance immediately (the resulting logits, and hence
$A$, would change under a global rotation, since nothing in
\eqref{def:attention-scratch} constrains how a rotated coordinate enters
the dot product). Geometry must instead enter through
\emph{relationships between atoms} --- the subject of the next section.

\section{Geometric graphs and message passing}
\label{sec:gnn-scratch}

\begin{definition}[Geometric graph; restated]
\label{def:geometric-graph-scratch}
A \emph{geometric graph} is $\mathcal G=(\mathcal V,\mathcal E,\{h_i\},\{x_i\})$:
nodes $i\in\mathcal V$ with $SE(3)$-invariant features $h_i\in\RR^d$ (atom
type, charge, fragment/entity type, time embedding of $t$) and positions
$x_i\in\RR^3$; edges $\mathcal E\subseteq\mathcal V\times\mathcal V$. Write
$x_{ij}=x_j-x_i$, $d_{ij}=\|x_{ij}\|$; we take $\mathcal E$ to be a radius
graph, $\{(i,j):d_{ij}<r_{\mathrm{cut}}\}$, for some fixed or per-edge-type
cutoff $r_{\mathrm{cut}}$.
\end{definition}

\begin{definition}[Message passing layer; restated]
\label{def:mpnn-scratch}
A message passing layer computes
\begin{equation}
\label{eq:mpnn-scratch}
m_{ij} = \phi_m(h_i,h_j,e_{ij}),
\qquad
m_i = \bigoplus_{j\in\mathcal N(i)}m_{ij},
\qquad
h_i' = \phi_u(h_i,m_i),
\end{equation}
for learned maps $\phi_m,\phi_u$, edge attributes $e_{ij}$, and a
permutation-invariant aggregator $\bigoplus$ (sum or mean over $j$).
\end{definition}

Since $\bigoplus$ depends only on the \emph{multiset} $\{m_{ij}\}$, not on
any ordering of $\mathcal N(i)$, \eqref{eq:mpnn-scratch} is permutation
equivariant by the same argument as \Cref{prop:perm-equivariance-scratch}
(indeed \Cref{def:attention-scratch} is the special case $\mathcal N(i)=\mathcal V$
for all $i$, $\bigoplus$ a learned convex combination). Geometry enters
through $e_{ij}$; the invariant (and hence, per
\Cref{sec:invariance-obstruction} below, safe-but-limited) choice expands
the distance in a radial basis,
\begin{equation}
\label{eq:rbf-scratch}
e_{ij} = \bigl(\psi_1(d_{ij}),\dots,\psi_K(d_{ij})\bigr),
\qquad
\psi_k(d) = \exp\bigl(-\gamma(d-\mu_k)^2\bigr)\,\chi(d),
\end{equation}
with basis centres $\mu_k$ and a smooth cutoff envelope $\chi$ vanishing
(along with enough derivatives to avoid a kink) at $d=r_{\mathrm{cut}}$: an
edge that is about to leave the radius graph as atoms move contributes a
message that continuously shrinks to zero \emph{before} it is dropped from
$\mathcal E$ entirely, rather than discontinuously disappearing --- without
this, the learned vector field would be discontinuous in $\mathbf x$
exactly at the cutoff radius, with predictable consequences for the
stability of the ODE/SDE integrators of
\Cref{chap:diffusion,chap:flow-matching}.

\begin{codebox}
\textbf{Mathematical object:} the radial basis $\psi_k(d)\chi(d)$ of
\eqref{eq:rbf-scratch} \\
\textbf{Implementation:} \file{net/encoder.py}::\pyfunc{BesselBasis},
\pyfunc{PolynomialCutoff} \\
\textbf{Interpretation:} the actual basis used is the (MACE-style) Bessel
form $\psi_k(d)=\sqrt{2/r_{\max}}\,\sin(k\pi d_{\mathrm{scaled}}/r_{\max})/(d_{\mathrm{scaled}}+\eps)$
(trainable frequencies $k\pi/r_{\max}$, not fixed Gaussian centres as in
\eqref{eq:rbf-scratch} --- this document presents the Gaussian form first
purely because it is the more elementary starting point, and records the
Bessel form actually used as a documented substitution, mathematically
serving an identical role: an orthogonal-ish family of radial basis
functions), multiplied by a degree-$8$ polynomial envelope
\[
\chi(d) = \Bigl[1-\tfrac{(p+1)(p+2)}2\Bigl(\tfrac d{r_{\max}}\Bigr)^p
+p(p+2)\Bigl(\tfrac d{r_{\max}}\Bigr)^{p+1}-\tfrac{p(p+1)}2\Bigl(\tfrac d{r_{\max}}\Bigr)^{p+2}\Bigr]\ind[d<r_{\max}],
\quad p=8,
\]
constructed exactly so that $\chi$ and its first two derivatives vanish at
$d=r_{\max}$ (a standard DimeNet/MACE-style envelope), confirming the
smooth-cutoff argument above at the level of the actual, exact formula
used. \textbf{Exact cutoff radii used}, per \file{oracle.py}: protein--ligand
interaction edges (\texttt{inter\_complex}) and ligand--ligand
cross-fragment edges (\texttt{inter\_fragments}), both $4.0\,$\r A;
C$\alpha$--C$\alpha$ protein-virtual edges, $18\,$\r A; ligand-virtual--protein-virtual
edges, $15\,$\r A --- these are distinct from, and should not be confused
with, the separate pocket-definition cutoffs
(\file{data.py}: \texttt{pocket\_distance\_cutoff}$=6.0$,
\texttt{pocket\_virtual\_cutoff}$=15.0$) governing which protein atoms are
extracted into the graph \emph{at all}, upstream of any of the above
edge-construction cutoffs.
\end{codebox}

\section{The obstruction: invariant features cannot output a vector}
\label{sec:invariance-obstruction}

We can now state, precisely, why \eqref{eq:mpnn-scratch} as given is
insufficient for \eqref{eq:network-signature-scratch}'s task.

\begin{proposition}[Invariant features cannot produce an equivariant vector output]
\label{prop:invariance-obstruction-scratch}
Suppose every node feature $h_i$ (initial and, by induction through
\eqref{eq:mpnn-scratch}, at every subsequent layer, since $\phi_m,\phi_u$
are ordinary functions of $SE(3)$-invariant inputs $h_i,h_j,e_{ij}$) is
$SE(3)$-invariant, and a readout $v_i=\rho(h_i)\in\RR^3$ is defined as a
function of $h_i$ alone. If $v_i$ is required to be $SE(3)$-equivariant in
the sense of \eqref{eq:equivariance-concrete} below, then $v_i\equiv0$.
\end{proposition}

\begin{equation}
\label{eq:equivariance-concrete}
v_i\bigl(\Phi_{(Q,b)}(\mathbf x)\bigr) = Q\,v_i(\mathbf x)
\qquad\forall (Q,b)\in SE(3) .
\end{equation}

\begin{proof}
By hypothesis, $h_i(\Phi_{(Q,b)}(\mathbf x)) = h_i(\mathbf x)$ for every
$(Q,b)$ (invariance), so
$v_i(\Phi_{(Q,b)}(\mathbf x)) = \rho(h_i(\Phi_{(Q,b)}(\mathbf x))) = \rho(h_i(\mathbf x)) = v_i(\mathbf x)$:
the left side of \eqref{eq:equivariance-concrete} equals $v_i(\mathbf x)$
identically, so \eqref{eq:equivariance-concrete} forces
$Qv_i(\mathbf x)=v_i(\mathbf x)$ for \emph{every} $Q\in SO(3)$. The only
vector in $\RR^3$ fixed by every rotation is $0$ (if $v\ne0$, take $Q$ a
rotation by $180^\circ$ about an axis perpendicular to $v$: $Qv=-v\ne v$).
\end{proof}

This is the obstruction promised at the end of \Cref{chap:groups}'s
motivation: a purely distance-based (invariant-feature) message passing
network can predict \emph{any scalar} (an energy, an affinity, a
confidence score) but literally cannot, for any choice of $\phi_m,\phi_u,\rho$,
output the translational or rotational velocity
\eqref{eq:network-signature-scratch} requires. Invariance is too strong a
symmetry constraint for a vector-valued task; equivariance (a strictly
weaker requirement, since invariance is the special case ``equivariant with
respect to the trivial representation'', \Cref{ex:ell01}) is needed
instead, and providing it is the entire content of
\Cref{chap:equiformer}.

\section{A minimal repair, and why it is not enough}
\label{sec:egnn-scratch}

Before building the full steerable machinery, it is instructive to see the
cheapest possible fix, since its specific limitations motivate exactly the
representation-theoretic apparatus of \Cref{chap:steerable}.

\begin{definition}[EGNN-style equivariant readout]
\label{def:egnn-scratch}
\[
v_i = \sum_{j\in\mathcal N(i)}\phi_v(h_i,h_j,d_{ij})\,x_{ij},
\qquad
\phi_v:\ \text{scalar-valued, built from invariant inputs only.}
\]
\end{definition}

\begin{proposition}
\Cref{def:egnn-scratch}'s $v_i$ satisfies \eqref{eq:equivariance-concrete}.
\end{proposition}

\begin{proof}
Under $\mathbf x\mapsto Q\mathbf x+b$: $x_{ij}\mapsto Qx_{ij}$ (translations
cancel in a difference of two points; the term ``$+b$'' subtracted from
itself), and $d_{ij}=\|x_{ij}\|$ is unchanged (\Cref{prop:so3-isometry}), so
$\phi_v(h_i,h_j,d_{ij})$, depending only on invariants, is unchanged too.
Hence $v_i\mapsto\sum_j\phi_v(\dots)\,Qx_{ij} = Q\sum_j\phi_v(\dots)x_{ij}=Qv_i$.
\end{proof}

This construction --- known as EGNN (E($n$)-Equivariant Graph Neural
Networks; Satorras et al., ICML 2021) --- is genuinely equivariant and
genuinely cheap, but its expressive limitation is structural, not merely
empirical, and worth stating precisely rather than only qualitatively.
\Cref{def:egnn-scratch}'s $v_i$ is a \emph{scalar-weighted sum of degree-$1$
vectors}: in the language of \Cref{chap:steerable}, it uses only the
$\ell_1=0$ (scalar weight $\phi_v$) $\otimes_{\mathrm{cg}}$ $\ell_2=1$
(vector $x_{ij}$) $\to\ell=1$ case of \Cref{ex:cg-special-cases}, and
\emph{no other Clebsch--Gordan path at all}. Two concrete
consequences follow directly from \Cref{prop:multichannel-linear}'s
degree-mixing constraint and \Cref{thm:cg-decomposition}'s decomposition
rule, not from vague ``expressivity'' intuition:

\begin{enumerate}
\item \textbf{Angular patterns among three or more neighbours are invisible.}
A ring of neighbours arranged symmetrically around node $i$, all at the
same distance and evenly spaced in angle, produces $\sum_jx_{ij}\approx0$
by symmetry (the vectors cancel), so \Cref{def:egnn-scratch}'s readout is
insensitive to the ring's orientation, even though that orientation is
physically meaningful (e.g.\ the plane of an aromatic ring relative to a
binding pocket wall); \Cref{ex:sph-low-degree}'s explicit degree-$2$
harmonic $\propto2x_3^2-x_1^2-x_2^2$ (a quadrupole) does \emph{not} cancel
under the same symmetric averaging, and this is precisely the sense in
which $\ell\ge2$ features carry information that $\ell\le1$ features
cannot.
\item \textbf{Axial (rotation-vector) outputs get the wrong reflection
parity unless tracked explicitly.} $\omega^{(i)}_\theta\in\so(3)\cong\RR^3$
is an \emph{axial} vector (\Cref{rem:parity-scratch}): under a reflection,
$\widehat\omega\mapsto -\det(Q)\,Q\widehat\omega Q^\top$ for $Q\in O(3)$
picks up an extra sign relative to an ordinary (\emph{polar}) vector like
$x_{ij}$ under the same reflection, a fact provable directly from
\Cref{prop:hat-cross}'s cross-product identity together with the general
transformation rule $(Ax)\times(Ay)=\det(A)A^{-\top}(x\times y)$ quoted in
the proof of \Cref{prop:so3-isometry}, now with $A=Q\in O(3)\setminus SO(3)$
so $\det Q=-1$. A construction of exactly the form of
\Cref{def:egnn-scratch}, built only from polar vectors $x_{ij}$, assigns
$\omega_\theta^{(i)}$ the \emph{polar} parity by default, which is simply
the wrong object unless the reflection sign is tracked and corrected
elsewhere --- an easy silent error, and the reason \Cref{chap:equiformer}
(and, per \Cref{rem:parity-scratch}, the discussion of molecular chirality)
carries a parity label alongside $\ell$.
\end{enumerate}

What is needed, and what \Cref{chap:steerable} already built, is: features
of every angular degree $\ell$, not only $0$ and $1$; and a way to combine
two such features, of \emph{any} degrees, into a third, without leaving the
class of equivariant objects --- exactly the Clebsch--Gordan tensor product
\eqref{eq:cg-scratch}. \Cref{chap:equiformer} now assembles this into a
complete, code-verified architecture.

\clearpage

% ================================================================
\chapter{Equiformer and EquiformerV2}
\label{chap:equiformer}
% ================================================================

This chapter is deliberately the slowest and most granular in the
document, per the stated goal of enabling a reader to open the actual
network source (\file{net/} in \file{SigmaFlow\_Development/src/sigmadock/})
and follow it line by line afterward. We assume from the reader only
\Cref{chap:groups,chap:steerable} (manifolds, $SO(3)$, spherical harmonics,
Wigner-$D$, Clebsch--Gordan) and \Cref{chap:pre-equiformer} (attention,
message passing, the invariance obstruction); nothing else from earlier
chapters is needed to follow the architecture itself, though its role as
the score/vector-field network of \Cref{chap:diffusion,chap:diffusion-so3}
is what it is \emph{for}.

The architecture implemented is EquiformerV2 (Liao et al., \emph{EquiformerV2:
Improved Equivariant Transformer for Scaling to Higher-Degree
Representations}, ICLR 2024), building on the original Equiformer (Liao \&
Smidt, ICLR 2023) and the eSCN reduction (Passaro \& Zitnick, \emph{Reducing
SO(3) Convolutions to SO(2) for Efficient Equivariant GNNs}, ICML 2023). We
derive the ideas in the order that makes them motivated, not necessarily
the historical order, and flag explicitly, in
\Cref{sec:v1-vs-v2-scratch}, the one place where the actual shipped
configuration is a genuine \emph{hybrid} of the two papers' recipes rather
than a pure implementation of either.

\section{From scalar messages to steerable messages}
\label{sec:steerable-messages}

\paragraph{Motivation.} \Cref{def:mpnn-scratch}'s message
$m_{ij}=\phi_m(h_i,h_j,e_{ij})$ was, implicitly, scalar-valued (or, at best,
per \Cref{def:egnn-scratch}, a scalar times a fixed direction vector). We
now let \emph{every} object in sight --- node features, messages, edge
embeddings --- be a full steerable feature in the sense of
\Cref{def:steerable-scratch}: a direct sum over degrees
$\ell=0,\dots,L$ (EquiformerV2's default, confirmed directly from the
constructor defaults in \file{net/model.py}, is $L=\texttt{lmax}=3$), each
degree carrying some number $C_\ell$ of channels.

\begin{definition}[Steerable node/edge feature, restated for the network]
\label{def:steerable-network-feature}
A node $i$'s feature is
$h_i = \bigoplus_{\ell=0}^L\bigoplus_{c=1}^{C_\ell}h_i^{\ell,c}\in\bigoplus_\ell(\RR^{2\ell+1})^{C_\ell}$,
transforming, under a rotation $Q$ of the whole scene, by
$h_i\mapsto\bigl(\bigoplus_{\ell,c}D^\ell(Q)\bigr)h_i$
(\Cref{def:steerable-scratch}). EquiformerV2's default channel counts are
uniform across degree, $C_\ell\equiv C=\texttt{sphere\_channels}=256$
(constructor default, \file{net/model.py}).
\end{definition}

The embedding step (\Cref{sec:node-embedding} below) initialises every
node's $\ell=0$ channels from ordinary categorical/continuous input
features (\Cref{sec:invariance-obstruction}'s invariant quantities are
exactly the correct content for a type-$0$ channel) and leaves every
$\ell\ge1$ channel at exactly zero; the network's job, across its layers,
is to populate the higher-degree channels with genuinely angular
information extracted from the geometry, and finally to read a $3$-vector
out of the $\ell=1$ channels as the network's geometric output
(\Cref{sec:equiformer-readout}).

\section{Embedding geometry: spherical harmonics on edges}
\label{sec:edge-embedding-scratch}

\paragraph{Motivation.} \Cref{chap:steerable} constructed, for each
degree $\ell$, an explicit map $Y^\ell:S^2\to\RR^{2\ell+1}$ satisfying the
equivariance law \eqref{eq:sph-equivariance-scratch}. This is exactly the
tool needed to lift a \emph{geometric} quantity --- the direction
$\hat x_{ij}=x_{ij}/d_{ij}$ of an edge --- into a \emph{steerable feature}
that can be combined, via the Clebsch--Gordan product, with a node's
existing steerable feature.

\begin{definition}[Equivariant edge embedding]
\label{def:edge-embedding-scratch}
Each edge $(i,j)$ is embedded as the pair
$\bigl(\psi(d_{ij})\chi(d_{ij}),\ Y^\ell(\hat x_{ij})_{\ell\le L}\bigr)$: an
invariant radial part (\Cref{eq:rbf-scratch}, code-verified form recorded in
\Cref{sec:gnn-scratch}'s code correspondence) and an equivariant angular
part, one $(2\ell+1)$-dimensional block per degree up to $L$.
\end{definition}

Recall from \Cref{ex:sph-low-degree} that $Y^0\equiv\mathrm{const}$ and
(up to normalisation) $Y^1(\hat x)=\hat x$; substituting $L=0$ or $L=1$ into
\Cref{def:edge-embedding-scratch} exactly recovers \Cref{sec:gnn-scratch}'s
distance-only GNN or \Cref{def:egnn-scratch}'s EGNN readout respectively:
the steerable-feature construction is a strict generalisation, not a
different framework, of everything built so far.

\section{The eSCN reduction: aligning every edge to its own frame}
\label{sec:equiformer-frame}

\paragraph{Motivation, and the problem eSCN solves.} A naive
implementation of steerable message passing would, at every edge, form a
full Clebsch--Gordan tensor product \eqref{eq:cg-scratch} between the
neighbour's feature $h_j$ (up to degree $L$) and the edge's spherical
harmonic $Y^{\ell'}(\hat x_{ij})$ (up to degree $L$), for every pair of
degrees $(\ell,\ell')$ --- a computation whose cost grows steeply with $L$,
since the number of $(\ell_1,\ell_2)\to\ell$ paths in
\Cref{thm:cg-decomposition} grows roughly cubically in $L$ (each pair
$(\ell_1,\ell_2)$ contributes up to $\min(2\ell_1,2\ell_2)+1$ output
degrees). EquiformerV2's central algorithmic contribution, inherited from
eSCN, is an exact reformulation that avoids ever explicitly instantiating a
Clebsch--Gordan coefficient table.

\begin{proposition}[Rotating an edge onto a fixed axis reduces $SO(3)$ to $SO(2)$]
\label{prop:esc-reduction}
Let $\hat x_{ij}\in S^2$ be an edge direction, and let $Q_{ij}\in SO(3)$ be
\emph{any} rotation with $Q_{ij}\hat x_{ij}=e_1$ (the first standard basis
vector; such $Q_{ij}$ always exists, and is unique up to an arbitrary
further rotation about $e_1$ itself, since $\{Q:Qe_1=e_1\}\cong SO(2)$).
Working in the rotated frame (applying $D^\ell(Q_{ij})$ to every steerable
feature, per \Cref{eq:sph-equivariance-scratch}), the edge direction itself
becomes the \emph{fixed} vector $e_1$, and \emph{every} subsequent
operation that only needs to be equivariant with respect to rotations
fixing the edge direction --- rather than all of $SO(3)$ --- may
therefore use the much smaller symmetry group $SO(2)$ (rotations about
$e_1$) instead of the full $SO(3)$, since $\hat x_{ij}$ itself no longer
transforms (it is a fixed constant, $e_1$, in this frame) and only the
\emph{other} steerable quantities (the rotated node features) need to
respect the residual $SO(2)$ symmetry about the now-fixed axis.
\end{proposition}

\begin{proof}[Justification, not a formal theorem]
This is a change-of-frame argument, not a new mathematical fact requiring
proof beyond \Cref{eq:sph-equivariance-scratch} itself: given any
$SO(3)$-equivariant operation $F$ acting on (steerable feature, edge
direction) pairs, define $\widetilde F(h,e_1) := D^\ell(Q_{ij})F\bigl(D^\ell(Q_{ij})^{-1}h,\hat x_{ij}\bigr)$;
by construction $\widetilde F$ agrees with $F$ after transporting through
the fixed rotation $Q_{ij}$, and since its second argument is now always
the fixed vector $e_1$, $\widetilde F$ need only respect rotations fixing
$e_1$ (the residual $SO(2)$) to guarantee $F$ itself was $SO(3)$-equivariant
--- the constraint on the \emph{first} argument's transformation law has
been reduced from a full $SO(3)$ representation to an $SO(2)$
representation, at the cost of applying $D^\ell(Q_{ij})$ and its inverse
once per edge.
\end{proof}

This is the eSCN insight stated precisely: \emph{per-edge}, rotate into a
frame where the edge direction is a fixed axis, do all the expensive
feature-mixing work under the much cheaper $SO(2)$-equivariance constraint
(an abelian group, whose irreps are $1$-dimensional --- ordinary complex,
equivalently $2\times2$-real, numbers indexed by a single integer $m$,
rather than $SO(3)$'s $(2\ell+1)$-dimensional irreps requiring the full
Clebsch--Gordan machinery), then rotate back.

\begin{codebox}
\textbf{Mathematical object:} the per-edge rotation $Q_{ij}$ of
\Cref{prop:esc-reduction} and its Wigner-$D$ realisation \\
\textbf{Implementation:} \file{net/edge\_rot\_mat.py}::\pyfunc{init\_edge\_rot\_mat}
(builds the orthonormal $3\times3$ frame: first axis $=\hat x_{ij}$; second
axis chosen, per-edge, as the least-aligned of several candidate reference
vectors, to avoid a near-degenerate cross product when constructing the
third axis; near-zero-length edges, $\|x_{ij}\|<10^{-5}$, fall back to a
fixed default direction, logging a warning) composed with
\file{net/wigner.py} (assembling the full block-diagonal Wigner-$D$ matrix,
via \texttt{e3nn}'s \texttt{o3.wigner\_D} applied to the Euler angles of
$Q_{ij}$, extracted via \texttt{o3.xyz\_to\_angles}) \\
\textbf{Tensor shape:} one $Q_{ij}\in\RR^{3\times3}$, and one block-diagonal
Wigner-$D$ matrix of size $\sum_{\ell=0}^L(2\ell+1)=\sum_{\ell=0}^3(2\ell+1)=16$
(for the default $L=3$), \emph{per edge} \\
\textbf{Interpretation:} every neighbour's steerable feature is rotated
into its (edge-specific!) aligned frame \emph{before} the message
computation of \Cref{sec:so2-conv} below, and the resulting message is
rotated back \emph{afterward}, via
\pyfunc{SO3\_Embedding.\_rotate}/\pyfunc{\_rotate\_inv}
(\file{transformer\_block.py}) --- i.e.\ this per-edge alignment happens
\emph{inside} every attention block, not once globally, since a different
edge has a different direction and hence a different aligned frame.
\end{codebox}

\section{$SO(2)$-equivariant convolution: the actual EquiformerV2 message}
\label{sec:so2-conv}

\paragraph{What an $SO(2)$-equivariant linear map looks like.} By the same
Schur's-lemma argument as \Cref{prop:schur-scratch}, but for the
\emph{abelian} group $SO(2)$ (whose irreps, unlike $SO(3)$'s, are complex
$1$-dimensional, indexed by an integer $m\in\ZZ$, with $Q_\alpha$ (rotation
by angle $\alpha$) acting on the degree-$m$ irrep by multiplication by
$e^{im\alpha}$): an $SO(2)$-equivariant linear map on the $m$-indexed
components of a steerable feature, once expressed in the aligned frame of
\Cref{prop:esc-reduction}, may mix channels \emph{within} a fixed order
$|m|$ (matching real and imaginary/paired components together) but not
across different values of $|m|$, and (since a real degree-$\ell$
spherical-harmonic component of order $m$ and one of order $-m$ transform
as a genuine, non-collapsible complex-conjugate pair under $SO(2)$ for
$m\ne0$) the map on this pair is exactly a $2\times2$ real matrix realising
multiplication by a complex number, while the single $m=0$ component per
degree is an ordinary real scalar with no complex structure.

\begin{definition}[$SO(2)$ convolution, informal]
\label{def:so2-conv}
Given a steerable feature $h$ (in the aligned frame), an $SO(2)$-equivariant
linear layer computes, for each order $m=0,1,\dots,\texttt{mmax}$: an
ordinary real-linear map on the (concatenated, across degrees $\ell\ge|m|$)
$m=0$ components; and, for $m\ge1$, a complex-linear map
$(x_r,x_i)\mapsto(W_1x_r-W_2x_i,\ W_1x_i+W_2x_r)$ on the (real, imaginary)
parts of the $m$-order components, with $W_1,W_2$ learned real matrices
mixing channels within that order. All weights ($W_1,W_2$, and the $m=0$
map's weights) are generated from an ordinary MLP applied to
\emph{invariant} scalars only (the radial+chemistry edge embedding), per
edge, per edge type.
\end{definition}

\begin{codebox}
\textbf{Mathematical object:} the $SO(2)$-equivariant message layer of
\Cref{def:so2-conv} \\
\textbf{Implementation:} \file{net/so2\_ops.py}::\pyfunc{SO2\_Convolution},
\pyfunc{SO2\_m\_Convolution} \\
\textbf{Interpretation:} this is, precisely, the departure from a naive
Clebsch--Gordan-tensor-product message
$m_{ij}=W_{ij}\bigl(h_j\otimes_{\mathrm{cg}}Y(\hat x_{ij})\bigr)$ that a
first-principles-only reading of \Cref{chap:steerable} might suggest as the
message formula, and it is worth being explicit that the \emph{executed
code does not form this Clebsch--Gordan product at all}: after the
edge-frame rotation of \Cref{sec:equiformer-frame}, the edge direction is
the fixed vector $e_1$, so mixing a node's steerable feature with the
(now-fixed, hence not needing to be explicitly multiplied in via CG
coefficients) edge direction reduces, by
\Cref{prop:esc-reduction}, to an $SO(2)$-equivariant linear operation on
the node feature alone, and it is \emph{this} operation, not a tensor
product, that \pyfunc{SO2\_Convolution} implements. Both are, by
\Cref{prop:esc-reduction}'s justification, mathematically legitimate
realisations of an equivariant message; the eSCN reduction is a genuine
\emph{algorithmic} improvement (removing the CG-path combinatorics
entirely, replacing them with ordinary/complex linear layers), not merely
a reparameterisation of the same computation, and this is the precise,
verifiable sense in which EquiformerV2 ``improves on'' the original
Equiformer's direct tensor-product message construction. The per-edge
radial MLPs generating $W_1,W_2$ have their \textbf{bias terms removed}
specifically for the transient (protein--ligand interaction) edge types
--- confirmed directly from the architecture description this monograph
is auditing, and consistent with \Cref{sec:gnn-scratch}'s smooth-cutoff
argument: with a bias, a message would not vanish even when the radial
features $\psi(d)\chi(d)$ do, at $d\to r_{\mathrm{cut}}$, breaking the
continuity property the envelope $\chi$ was introduced to guarantee.
\end{codebox}

\section{Equivariant nonlinearity: gating and the separable $S^2$ activation}
\label{sec:equiformer-nonlinearity}

\paragraph{Why an ordinary nonlinearity breaks equivariance.} By
\Cref{prop:schur-scratch}, the \emph{only} linear equivariant operations
are per-degree scalar rescalings; a network built from linear layers alone
is itself linear (up to the invariant-scalar-weight structure) and
severely limited in expressive power, exactly as an ordinary neural network
with no nonlinearity would be. But an entrywise nonlinearity
$\sigma$ applied directly to a degree-$\ell\ge1$ feature, $h\mapsto\sigma(h)$,
does not commute with $D^\ell(Q)$ in general (a direction-dependent
nonlinear reshaping of a vector is generally not the same before and after
rotating the vector), so it is not equivariant, and cannot be used
directly.

\begin{definition}[Gate nonlinearity]
\label{def:gate-scratch}
For a steerable feature channel $h^{\ell,c}$ ($\ell\ge1$) and an associated,
separately-computed \emph{invariant} (type-$0$) scalar $g^{\ell,c}$,
\[
h^{\ell,c}\ \longmapsto\ \sigma(g^{\ell,c})\cdot h^{\ell,c},
\]
for an ordinary scalar nonlinearity $\sigma$ (e.g.\ a sigmoid or SiLU).
\end{definition}

\begin{proposition}
\Cref{def:gate-scratch}'s map is equivariant.
\end{proposition}

\begin{proof}
Under a rotation, $g^{\ell,c}\mapsto g^{\ell,c}$ (invariant, by
construction --- it is produced by a separate, invariant-only computation
path) and $h^{\ell,c}\mapsto D^\ell(Q)h^{\ell,c}$, so
$\sigma(g^{\ell,c})h^{\ell,c}\mapsto\sigma(g^{\ell,c})D^\ell(Q)h^{\ell,c} = D^\ell(Q)\bigl(\sigma(g^{\ell,c})h^{\ell,c}\bigr)$
using linearity of $D^\ell(Q)$ and that $\sigma(g^{\ell,c})$, an unchanged
scalar, commutes past the matrix.
\end{proof}

Geometrically: the gate rescales a steerable feature's \emph{magnitude}
nonlinearly (via the invariant scale $\sigma(g)$) while leaving its
\emph{direction} (its transformation behaviour) completely untouched ---
exactly the one degree of freedom (an overall scalar multiplier) that
\Cref{prop:schur-scratch} says is available on a single channel without
violating equivariance, now made nonlinear in the \emph{scalar} channel
that controls it.

\begin{codebox}
\textbf{Mathematical object:} the equivariant nonlinearity actually used by
default \\
\textbf{Implementation:} \file{net/activation.py}::\pyfunc{SeparableS2Activation}
(config default: \texttt{use\_gate\_act=False}, \texttt{use\_sep\_s2\_act=True}) \\
\textbf{Interpretation:} EquiformerV2's default nonlinearity is \emph{not}
literally \Cref{def:gate-scratch} but a closely related, more expressive
construction: the steerable feature is projected onto a discrete grid of
points on $S^2$ (via \texttt{e3nn}'s \pyfunc{ToS2Grid}, resolution set by
$L$/\texttt{mmax}), an \emph{ordinary}, pointwise nonlinearity is applied at
each grid point (legitimate precisely because evaluating a steerable
function at a rotated grid point equals evaluating the un-rotated function
at the inverse-rotated point --- the grid-sampling analogue of
\Cref{prop:sph-transform}'s transformation law, applied to function values
rather than coefficients), and the result is projected back
(\pyfunc{FromS2Grid}) into spherical-harmonic coefficients. An additional,
separate scalar ``gating'' path (extra invariant channels produced
alongside the main $SO(2)$-convolution output) modulates this activation
further, combining both mechanisms. We record this as the precise
implemented nonlinearity without re-deriving the grid-sampling equivariance
argument in full technical detail (it is a standard construction in the
equivariant-networks literature, e.g.\ Cohen et al., \emph{Spherical CNNs},
ICLR 2018).
\end{codebox}

\section{Equivariant linear layers and normalisation}
\label{sec:equiformer-linear-norm}

\begin{codebox}
\textbf{Mathematical object:} the general equivariant linear map of
\Cref{prop:multichannel-linear} \\
\textbf{Implementation:} \file{net/so3.py}::\pyfunc{SO3\_LinearV2} \\
\textbf{Tensor shape:} weight tensor \shape{lmax+1, out\_features, in\_features} \\
\textbf{Interpretation:} implements exactly
$h^{\ell,c}\mapsto\sum_{c'}W^\ell_{cc'}h^{\ell,c'}$ via
\texttt{einsum("bmi,moi->bmo", ...)} with the weight broadcast identically
across the $2\ell+1$ components $m$ of a given degree $\ell$ (one weight
matrix \emph{per degree}, shared across all $m$ within it, exactly matching
\Cref{prop:schur-scratch}'s statement that an intertwiner is a scalar
multiple of the identity \emph{on} $V_\ell$, generalised here to multiple
channels per degree via \Cref{prop:multichannel-linear}); and a bias term
is added \textbf{only to the $\ell=0$ slice} of the output
(\texttt{out[:,0:1,:] += bias}) --- the code's own structural enforcement
of the fact, provable directly from \Cref{prop:schur-scratch} (an additive
constant is a map from the trivial representation, $\ell=0$, into the
target; by Schur's lemma such a map can only land in the $\ell=0$
component of the target without violating equivariance, since a nonzero
constant vector in a degree-$\ell\ge1$ channel would not itself be fixed
by every rotation, contradicting $D^\ell(Q)\cdot\mathrm{bias}=\mathrm{bias}$
required for a translation-equivariant \emph{additive} bias), that
higher-degree channels cannot carry an additive bias at all.
\end{codebox}

\begin{codebox}
\textbf{Mathematical object:} equivariant normalisation \\
\textbf{Implementation:} \file{net/layer\_norm.py} (\texttt{norm\_type}
default \texttt{"rms\_norm\_sh"}) \\
\textbf{Interpretation:} normalises each degree-$\ell$, channel-$c$ block
$h^{\ell,c}$ by (an RMS-style estimate of) its own \emph{norm}
$\|h^{\ell,c}\|_2$ --- an invariant quantity by \Cref{prop:so3-isometry}'s
orthogonality of $D^\ell(Q)$, since $\|D^\ell(Q)v\|=\|v\|$ for orthogonal
$D^\ell(Q)$ (\Cref{prop:wigner-properties}(iii)) --- rather than by any
coordinate-wise statistic, which is what makes the normalisation
equivariant (dividing every component of $h^{\ell,c}$ by the same
invariant scalar $\|h^{\ell,c}\|$ commutes with $D^\ell(Q)$ by the same
scalar-multiplication argument as \Cref{def:gate-scratch}'s gate). The
exact formula implemented (beyond ``RMS of the norm, per spherical-harmonic
degree'') was not read in the code-research pass underlying this document
and is flagged, per this document's stated policy, as not independently
confirmed beyond its name and role.
\end{codebox}

\section{Equivariant attention weights}
\label{sec:equiformer-attention}

\paragraph{The constraint attention weights must satisfy.} We now assemble
messages and weights into the full attention block, and the first thing to
determine is what kind of object the attention weight $a_{ij}$
(\Cref{def:attention-scratch}'s analogue) is allowed to be.

\begin{proposition}[Attention weights must be invariant scalars]
\label{prop:attention-weights-scalar}
Let $h_i' = h_i + \sum_{j\in\mathcal N(i)}a_{ij}\,m_{ij}$, with each message
$m_{ij}$ steerable of a fixed degree $\ell$. If the $a_{ij}$ are
$SE(3)$-invariant scalars, this update is equivariant. If $a_{ij}$ were
instead permitted to be an arbitrary equivariant (not merely invariant)
quantity of matching degree, the sum $\sum_ja_{ij}m_{ij}$ would generally
fail to commute with $D^\ell(Q)$, since $a_{ij}$ itself would then also
transform under $Q$, and the product of two degree-$\ell$-transforming
quantities is not, in general, again degree-$\ell$ (\Cref{thm:cg-decomposition}).
\end{proposition}

\begin{proof}
If $a_{ij}$ is invariant and $m_{ij}\mapsto D^\ell(Q)m_{ij}$:
$\sum_ja_{ij}m_{ij}\mapsto\sum_ja_{ij}D^\ell(Q)m_{ij}=D^\ell(Q)\sum_ja_{ij}m_{ij}$,
using linearity of $D^\ell(Q)$ and that scalar $a_{ij}$ commutes past it.
Adding $h_i\mapsto D^\ell(Q)h_i$ (assuming $h_i$ is also degree-$\ell$;
otherwise, apply the same argument degree-by-degree in a multi-degree sum)
preserves equivariance of $h_i'$. If $a_{ij}$ instead transformed under $Q$
(e.g.\ if it were itself a component of a degree-$\ell'>0$ feature), the
product $a_{ij}m_{ij}$ would transform by a Clebsch--Gordan combination of
degrees $\ell'$ and $\ell$, generally landing (partially) outside degree
$\ell$, breaking the claimed degree-$\ell$ equivariance of the sum.
\end{proof}

This forces, structurally, exactly the design choice \Cref{chap:pre-equiformer}
already used implicitly for ordinary attention (\Cref{def:attention-scratch}'s
softmax output is a scalar) and is why an equivariant network cannot simply
reuse a dot-product attention score computed \emph{directly} between two
full multi-degree steerable features (which would not, in general, be
invariant, unless the specific $\ell\to0$ Clebsch--Gordan path of
\Cref{ex:cg-special-cases} is used deliberately) without first projecting
onto an invariant quantity.

\begin{codebox}
\textbf{Mathematical object:} the attention weight $a_{ij}$ of
\Cref{prop:attention-weights-scalar} \\
\textbf{Implementation:} \file{net/transformer\_block.py}::\pyfunc{SO2EquivariantGraphAttention.forward} \\
\textbf{Interpretation:} $a_{ij}$ is computed \textbf{only} from the extra
$\ell=0,m=0$ scalar output channels of the first $SO(2)$-convolution
(\Cref{sec:so2-conv}) applied to the edge message --- i.e.\ the network
allocates a few dedicated invariant ``gating'' channels alongside its main
steerable output, specifically for this purpose --- via a
$\LN\to\mathrm{SmoothLeakyReLU}\to$ learned per-head dot product, then a
standard \texttt{torch\_geometric.utils.softmax} keyed on the
\emph{target} node index (\texttt{edge\_index[1]}), i.e.\ a per-target-node
softmax over its incoming edges, exactly \Cref{def:attention-scratch}'s
row-wise softmax specialised to a sparse graph. This is, precisely, the
architectural enforcement of \Cref{prop:attention-weights-scalar}: there is
structurally no other path in the code by which a non-scalar, degree-$\ge1$
quantity could enter the attention weight.
\end{codebox}

\subsection{Assembling the full attention block}

\begin{proposition}[The full block is equivariant]
\label{prop:full-block-equivariant}
Let messages $m_{ij}$ be constructed (via \Cref{sec:equiformer-frame,sec:so2-conv})
to be equivariant, and $a_{ij}$ invariant scalars
(\Cref{prop:attention-weights-scalar}). Then
$h_i' = h_i + \sum_{j\in\mathcal N(i)}a_{ij}\,m_{ij}$ is equivariant,
degree-by-degree.
\end{proposition}

This is exactly \Cref{prop:attention-weights-scalar}'s proof, restated as
the conclusion for the assembled block. In code, the full per-layer
pipeline (\pyfunc{TransBlockV2}, \file{net/transformer\_block.py}) is: (1)
concatenate source and target node features; (2) rotate into the per-edge
frame (\Cref{sec:equiformer-frame}); (3) first $SO(2)$-convolution
(\Cref{sec:so2-conv}), also producing the invariant gating channels used
for attention weights; (4) equivariant nonlinearity
(\Cref{sec:equiformer-nonlinearity}); (5) second $SO(2)$-convolution; (6)
compute $a_{ij}$ from the gating channels (\Cref{sec:equiformer-attention});
(7) rotate messages back out of the edge frame; (8) scatter-sum
$a_{ij}m_{ij}$ into target nodes (an \texttt{index\_add\_} along
\texttt{edge\_index[1]}, the permutation-equivariant aggregator of
\Cref{def:mpnn-scratch}); (9) final equivariant linear projection
(\pyfunc{SO3\_LinearV2}, \Cref{sec:equiformer-linear-norm}); embedded, as in
\Cref{eq:transformer-block-scratch}, in a pre-normalised residual structure
alongside a per-node equivariant feed-forward block.

\section{Node embedding: turning chemistry into a steerable feature}
\label{sec:node-embedding}

Before any message passing, every node needs an initial steerable feature
(\Cref{sec:steerable-messages}); by \Cref{sec:invariance-obstruction}'s
argument in reverse, only the $\ell=0$ (invariant) channels can be
populated directly from ordinary categorical/continuous chemistry, since
nothing in the raw per-atom input (an atom type, a formal charge, a time
value) carries any directional information to seed a higher-degree channel
with.

\begin{codebox}
\textbf{Mathematical object:} the initial (type-$0$-only) node embedding \\
\textbf{Implementation:} \file{net/encoder.py}::\pyfunc{AtomDiffusionEncoder},
called once per node-entity type from \file{net/model.py}::\pyfunc{compute\_node\_embedding} \\
\textbf{Interpretation:} for each of the (default) \textbf{ten} categorical
atom features (\texttt{atom\_feature\_dims}$=(54,5,3,4,4,7,2,2,4,3)$,
\file{config.py} --- cardinalities only; the exact chemical identity of
each of the ten slots, e.g.\ which one is atomic number versus
hybridisation, was not independently confirmed against the preprocessing
code in the pass underlying this document and is left as an open item, not
asserted), a separate \pyfunc{nn.Embedding} table (with one extra row for
an explicit ``unknown/out-of-range'' category) produces a vector; these ten
vectors are combined not by concatenation-then-MLP but --- in the
configuration actually used, \texttt{linear\_aggregate=False} at every call
site --- by \textbf{summing them and dividing by $\sqrt{10}$}, a
variance-preserving sum (if the ten embeddings were independent with unit
variance per coordinate, the sum's variance would be $10\times$ larger, and
dividing by $\sqrt{10}$ restores unit variance --- the same normalisation
principle as the $1/\sqrt{d_k}$ scaling of \Cref{def:attention-scratch}'s
attention logits). Continuous auxiliary features (e.g.\ protein residue/ESM
embeddings, where used) are instead passed through a small MLP and
\emph{added}. The time embedding $t_{\mathrm{emb}}$ --- a
\texttt{sinusoidal} embedding (default \texttt{t\_emb\_dim=32},
\texttt{t\_emb\_scale=10000.0}; the exact functional form was not
independently re-derived from the code in this pass, but the name and
parameters strongly suggest the standard Transformer construction
$\bigl(\sin(t/10000^{2k/d}),\cos(t/10000^{2k/d})\bigr)_{k}$, flagged here as
plausible but not code-confirmed) --- is concatenated (not summed) onto the
aggregated categorical embedding and passed through one further linear
layer (\pyfunc{time\_projector}). \textbf{Ligand-fragment virtual nodes}
(one per fragment, per \Cref{sec:sigmadock-graph} below) do
\emph{not} get an independent embedding at all: their $\ell=0$ feature is
the \emph{mean} (a \texttt{scatter}, \texttt{reduce="mean"}) of the
already-computed embeddings of every real atom belonging to that fragment,
propagated along dedicated virtual-to-atom (\texttt{ligand\_v2a}) edges ---
i.e.\ a fragment's virtual node is, by construction, exactly
$SE(3)$-invariant-feature-wise a summary of its members, before any message
passing layer has even run.
\end{codebox}

\section{Reading out the vector field}
\label{sec:equiformer-readout}

After $\texttt{num\_layers}=6$ (default) blocks, each atom carries a full
steerable feature $h_i=\bigoplus_{\ell=0}^3\bigoplus_{c=1}^{256}h_i^{\ell,c}$.
The network's geometric output, per
\eqref{eq:network-signature-scratch}, is a single $3$-vector per
\emph{atom} (not yet per fragment; the reduction to per-fragment
force/torque via Newton--Euler mechanics is a separate step, reconstructed
precisely in \Cref{chap:sigmadock}), read off from the $\ell=1$ channels by
one further equivariant linear layer.

\begin{codebox}
\textbf{Mathematical object:} the per-atom pseudo-force readout \\
\textbf{Implementation:} \file{net/model.py}::\pyfunc{EquiformerV2.forward},
final \pyfunc{force\_block} (an \pyfunc{SO2EquivariantGraphAttention} with a
single output channel), then
\texttt{forces = force\_block\_output.embedding.narrow(1, 1, 3)} \\
\textbf{Tensor shape:} \shape{N\_atoms, 3} \\
\textbf{Interpretation:} \pyfunc{narrow(1, 1, 3)} slices the flattened
$(L+1)^2=16$-length coefficient axis starting at index $1$, length $3$ ---
i.e.\ it selects exactly the three $\ell=1$ ($m=-1,0,1$) coefficients,
discarding the $\ell=0$ coefficient (index $0$) and every $\ell\ge2$
coefficient (indices $4$ and up) entirely. By
\Cref{prop:wigner-properties}(iv), this $\RR^3$-valued output transforms
under a global rotation $Q$ exactly as $D^1(Q)=Q$, i.e.\ as an ordinary
polar $3$-vector in the \emph{world frame} in which the network's input
coordinates were expressed --- this is precisely the ``atom-wise pseudo-force
in the world frame'' object that \Cref{chap:sigmadock}'s Newton--Euler
reduction (net force, torque, per fragment) consumes, and it is only at
that later, non-network step that a fragment's rotational output is
converted into the \emph{body-frame} (right-trivialised,
\Cref{sec:so3-trivialisation}) angular quantity
$\omega^{(i)}_\theta\in\so(3)$ that
\eqref{eq:network-signature-scratch} ultimately requires --- the network
itself, note carefully, never outputs a body-frame quantity directly; the
conversion is a geometric post-processing step external to
\pyfunc{EquiformerV2} proper, and \Cref{chap:sigmadock} traces it precisely.
\end{codebox}

\section{EquiformerV1 vs.\ EquiformerV2: what actually changed, and what this codebase actually ships}
\label{sec:v1-vs-v2-scratch}

We can now state the V1-versus-V2 comparison precisely, having derived
every ingredient involved, rather than merely asserting a version number.

\begin{description}
\item[Original Equiformer (V1).] Messages are formed by an explicit
Clebsch--Gordan tensor product \eqref{eq:cg-scratch} between a neighbour's
steerable feature and the edge's spherical harmonic embedding
(\Cref{def:edge-embedding-scratch}), with per-path learned weights; no
per-edge frame alignment is used. Attention weights are computed from
invariant edge/node scalars via an ordinary MLP $\to$ dot product $\to$
softmax pipeline, structurally identical to
\Cref{def:attention-scratch}'s original Transformer recipe, merely applied
to steerable rather than plain features.
\item[EquiformerV2 (eSCN-based).] Messages are formed via the per-edge
frame alignment of \Cref{sec:equiformer-frame} and the $SO(2)$-convolution
of \Cref{sec:so2-conv}, \emph{replacing} the explicit Clebsch--Gordan
tensor product with an algorithmically cheaper, exactly equivalent (by
\Cref{prop:esc-reduction}'s justification) construction. V2's own paper
additionally proposes an alternative attention-weight mechanism,
\texttt{use\_s2\_act\_attn}, computing attention directly from grid-sampled
($S^2$-activation-based) features rather than V1's separate scalar-dot-product
path.
\end{description}

\begin{caution}[This codebase implements a hybrid, not pure V2]
\label{caution:hybrid-architecture}
The constructor of \pyfunc{EquiformerV2} in this codebase sets
\texttt{use\_s2\_act\_attn: bool = False} \emph{and enforces this with an
explicit assertion} (\texttt{assert not self.use\_s2\_act\_attn}, with an
inline comment recording that this path is unused). That is: the message
mechanism is EquiformerV2's eSCN-based $SO(2)$-convolution
(\Cref{sec:so2-conv}) throughout, but the \emph{attention-weight}
computation follows the \textbf{original Equiformer(V1) recipe} ---
separate invariant scalar channels, MLP, per-head dot product, softmax
(\Cref{sec:equiformer-attention}) --- rather than V2's own alternative
\texttt{use\_s2\_act\_attn=True} path. This is a genuine, verifiable,
non-cosmetic architectural fact about the shipped model, not a matter of
degree or approximation, and this document states it explicitly rather
than describing the backbone as ``EquiformerV2'' without qualification, as
a less careful account might.
\end{caution}

\section{Summary: architecture, hyperparameters, and parameter count}
\label{sec:equiformer-summary}

\begin{center}
\begin{tabular}{@{}ll@{}}
\toprule
Hyperparameter & Default value (constructor, \file{net/model.py}) \\
\midrule
Maximum degree $L$ (\texttt{lmax}) & $3$ \\
Maximum order $M$ (\texttt{mmax}, $SO(2)$-conv truncation) & $2$ \\
Number of resolutions & $1$ (single $(L,M)$ pair) \\
Channels per degree (\texttt{sphere\_channels}) & $256$ \\
Number of attention layers (\texttt{num\_layers}) & $6$ \\
Attention heads (\texttt{num\_heads}) & $4$ \\
Edge/radial embedding dim & $64$ / $32$ basis functions \\
Polynomial cutoff order & $8$ \\
Time embedding & sinusoidal, dim $32$, scale $10^4$ \\
Normalisation & \texttt{rms\_norm\_sh} (per-degree RMS norm) \\
Radial cutoff function & Bessel \\
\bottomrule
\end{tabular}
\end{center}
Total learnable parameters, as reported in every training log inspected
for this document: $\mathbf{14{,}979{,}377}$.

\begin{remark}[What this chapter has, and has not, independently verified]
Several fine-grained implementation details --- the exact RMS-norm formula,
the exact sinusoidal time-embedding formula, the internals of the
\pyfunc{S2Activation}/\pyfunc{GateActivation} classes as distinguished from
\pyfunc{SeparableS2Activation}, and the precise semantics of each of the
ten categorical atom features --- were, per this document's stated policy,
flagged rather than guessed at in the code-research pass underlying this
chapter, because the corresponding source files (\file{net/layer\_norm.py},
\file{net/timestep\_embedder.py}, \file{net/activation.py} in full,
\file{chem/processing.py}) were not read to completion. Everything else in
this chapter --- the edge-frame mechanism, the $SO(2)$-convolution message
construction, the attention-weight computation and its V1/V2 hybrid status,
the equivariant linear layer's degree-block structure, the node embedding's
aggregation formula, and the final $\ell=1$ readout --- was confirmed
directly against the source.
\end{remark}

\clearpage

% ================================================================
\part{SigmaDock}
\label{part:sigmadock}
% ================================================================

\chapter{SigmaDock: Full Reconstruction}
\label{chap:sigmadock}

Every mathematical ingredient is now available:
\Cref{chap:groups,chap:steerable} give the geometry,
\Cref{chap:diffusion,chap:diffusion-so3} give the generative-modelling
machinery, and \Cref{chap:pre-equiformer,chap:equiformer} give the network.
This chapter assembles them into SigmaDock exactly as implemented, tracing
every design decision back to either a theorem already proved or a specific
line of code, and marking clearly the few points that remain genuinely
open (not independently re-derivable from a static reading of the code
alone).

\section{Choosing the state space}
\label{sec:sigmadock-state-space}

\Cref{sec:pose-dof} identified three coordinate systems in which a pose
could be represented: all $3n$ Cartesian coordinates, the
$(k{+}6)$-dimensional torsional parameterisation
$\bigl(SE(3)\times\mathbb T^k\bigr)$, or a fragment-based parameterisation
$\bigl(SE(3)^m\bigr)$. SigmaDock's central architectural decision is to use
the third, and the argument for doing so is a statement about the induced
\emph{measure}, not merely about dimensionality --- worth deriving
carefully, since it is what licenses the product-manifold construction
(\Cref{prop:se3-product}) underlying every diffusion/flow-matching formula
in this document.

\begin{definition}[Pullback volume form]
\label{def:pullback-volume}
Let $\Omega:\mathcal Y\to\RR^{n\times3}$ be a smooth parameterisation map
from a parameter space $\mathcal Y$ into Cartesian coordinates, with
Jacobian $J_\Omega=\partial\Omega/\partial y$ and Gram matrix
$G_\Omega=J_\Omega^\top J_\Omega$. The \emph{pullback volume form} induced
on $\mathcal Y$ is $d\mu_\Omega(y) = \sqrt{\det G_\Omega(y)}\,dy$ ---
the Riemannian volume form of the metric on $\mathcal Y$ obtained by
declaring $\Omega$ a local isometry onto its image with the ambient
Euclidean metric.
\end{definition}

\begin{theorem}[Torsional entanglement]
\label{thm:entanglement-scratch}
Let $\psi:SE(3)\times\mathbb T^k\to\RR^{n\times3}$ be the torsional
parameterisation of a ligand with $k$ rotatable bonds (global rigid motion
composed with dihedral-angle-driven internal reconstruction), and let
$\varphi:SE(3)^m\to\RR^{n\times3}$ be the fragment parameterisation with
$m$ pairwise-disjoint rigid fragments. Then, generically (for a molecular
topology in which some rotatable bond's rotation displaces atoms that a
different rotatable bond also displaces --- true for essentially every
branched or chain topology with more than one rotatable bond),
$G_\psi$ is non-diagonal and $d\mu_\psi$ is not a product measure over the
individual torsion coordinates. By contrast, $J_\varphi$ is block-diagonal
(one block per fragment, since moving fragment $i$'s rigid motion changes
no coordinate belonging to fragment $j\ne i$), so
$\det G_\varphi = \prod_{i=1}^m\det(J_i^\top J_i)$ and $d\mu_\varphi$
\emph{is} a product measure over fragments.
\end{theorem}

\begin{proof}
\emph{Fragment case.} By definition of $\varphi$, atom $a$'s position
depends only on the rigid motion of the fragment containing $a$:
$\partial\varphi_a/\partial(\text{fragment }j\text{'s coordinates})=0$
whenever $a\notin\text{fragment }j$. Hence $J_\varphi$, with rows grouped
by atom and columns grouped by fragment, is block-diagonal (a given row has
nonzero entries only in the column-block of its own fragment), and for a
block-diagonal matrix $J=\bigoplus_iJ_i$, $J^\top J=\bigoplus_iJ_i^\top J_i$
is also block-diagonal, with $\det$ of a block-diagonal matrix the product
of the blocks' determinants.

\emph{Torsional case.} Let $\phi_i,\phi_j$ be two rotatable-bond angles
whose associated atom-displacement fields overlap (an atom $a$ moved by
rotating about bond $i$ is also moved by rotating about bond $j$ --- true
whenever the two bonds lie on a common path through the molecular graph
with atom $a$ downstream of both, which occurs for any chain or branched
topology with $\ge2$ non-independent rotatable bonds). Then
$(G_\psi)_{ij}=\sum_a(\partial_{\phi_i}\psi_a)\cdot(\partial_{\phi_j}\psi_a)$
is generically nonzero, since it is a sum, over the atoms moved by
\emph{both} torsions, of dot products of two generically-non-orthogonal
displacement vectors --- there is no structural reason for this sum to
vanish identically across all molecular geometries, and it generically does
not. Since $(G_\psi)_{ij}\ne0$ for some $i\ne j$, $G_\psi$ is not diagonal,
and $\det G_\psi$ does not factor as a product over the individual
coordinates $\phi_i$ (a non-diagonal Gram matrix's determinant is not, in
general, the product of its diagonal entries), so $d\mu_\psi(\bm\phi)$
cannot be written as $\prod_id\mu_i(\phi_i)$ for single-coordinate
measures $d\mu_i$.
\end{proof}

The mechanism, made concrete: the Jacobian column for torsion $\phi_i$ is
the Cartesian velocity field generated by rotating everything
\emph{downstream} of bond $i$ about that bond's axis, and in a chain or
branched molecule this field moves \emph{many} atoms, with the sets of
atoms moved by different torsions typically overlapping. Three
\emph{practical} consequences follow directly from
\Cref{thm:entanglement-scratch}, each of which would separately motivate
avoiding the torsional parameterisation for a diffusion or flow-matching
model specifically (as opposed to, e.g., a search-based method, for which
the measure's factorisation is less directly load-bearing):

\begin{enumerate}
\item \textbf{Independent conditional noise does not stay independent.} A
forward noising kernel that factorises as a product over torsion
coordinates --- exactly the kind of kernel \Cref{chap:diffusion} builds,
one independent Gaussian/heat-kernel perturbation per coordinate --- maps,
under $\psi$, to a correlated, anisotropic perturbation of the Cartesian
positions the network actually observes (since $\psi$'s Jacobian mixes the
torsion coordinates together non-trivially, per
\Cref{thm:entanglement-scratch}). The network must then implicitly learn
to invert this metric-induced coupling, an unnecessarily hard learning
problem compared to \Cref{prop:se3-product}'s clean, coupling-free product
structure on $SE(3)^m$.
\item \textbf{Global and internal motion do not decouple.} A torsional
increment generically shifts the fragment's centre of mass
($\sum_am_a\partial_{\phi_i}\psi_a\ne0$ in general, since the atoms
downstream of bond $i$ do not, in general, have their displacements sum to
zero) and induces a net torque, so the $\mathbb T^k$ and $SE(3)$ factors of
the torsional parameterisation are not, geometrically, independent either
--- any diffusion or flow constructed to treat them as independent (again,
the natural, simple construction) is therefore only approximately correct,
exactly (an infinitesimal-step) in the limit, not exactly at any finite
step size.
\item \textbf{The dihedral definition itself has a gauge ambiguity.} A
dihedral angle $\phi_{ABCD}$ does not, by itself, specify \emph{which side}
of the bond $BC$ is held fixed while the other rotates; a change
$\Delta\phi$ can be realised by rotating the $A$-side, the $D$-side, or any
split between the two, and these give physically different (though
dihedral-angle-equal) intermediate trajectories. Any concrete
implementation must fix a convention, and a learned score/vector-field
network must be trained consistently with whichever convention sampling
uses --- a consistency requirement with no natural enforcement mechanism,
unlike the fragment parameterisation, where no such choice ever arises
(each fragment simply \emph{is} a specified rigid body).
\end{enumerate}

Fragmentation dissolves all three simultaneously, \emph{because} of
\Cref{thm:entanglement-scratch}'s block-diagonal Jacobian: moving one
fragment's rigid motion changes no other fragment's coordinates at all, so
independent per-fragment noise \emph{stays} independent under $\varphi$,
global and (now nonexistent, at the fragment level) internal motion trivially
coincide, and there is no gauge ambiguity in specifying a rigid motion. The
price is dimensional: $6m$ degrees of freedom against the torsional
parameterisation's $k+6$, and $m>k+1-m$ generically (fragmenting at every
rotatable bond, without merging, gives $\hat m=k+1$ fragments, so $6\hat m=6k+6$
against $k+6$ --- \emph{more} parameters than the torsional route, not
fewer), which is exactly the deficit \Cref{sec:sigmadock-fr3d} below
addresses.

\begin{caution}[Status of this argument]
\Cref{thm:entanglement-scratch} and its practical consequences are this
document's own reconstruction and elaboration of the reasoning given in
the thesis draft this monograph audits (\file{Texte/main.tex}); we have
verified the mathematical content --- the block-diagonal-Jacobian argument
for fragments, and the generic-non-diagonality argument for torsions ---
independently, by direct derivation, and found it sound, but we have
\emph{not} located a specific published SigmaDock source stating the
theorem in this exact form, and it is presented here as a rigorous,
self-contained argument for a design choice the codebase clearly makes
(fragment-based, not torsional, generation), rather than as a verbatim
quotation from a paper.
\end{caution}

\section{The conformational manifold: what rigidity costs}
\label{sec:conformational-manifold}

Treating each fragment as perfectly rigid is itself an approximation, and
its cost should be quantified, not merely assumed acceptable. In vacuum,
the accessible conformations of a ligand (fixed molecular topology)
concentrate, under a Boltzmann distribution at physiological temperature,
near a manifold
$\mathcal M_c=\{x_c:g(x_c)\approx0\}$ where $g$'s components softly encode
bond lengths ($\|r_{AB}\|-d_0$, $d_0$ the equilibrium bond length) and bond
angles ($\hat u\cdot\hat v-\cos\tau_0$), \emph{excluding} dihedrals (which
are, by contrast, the ligand's soft/free coordinates, exactly the ones
\Cref{sec:pose-dof} identified as free rotatable-bond degrees of freedom).
$\mathcal M_c$ is sampled, in practice, by conformer-generation software
(RDKit's ETKDGv3 algorithm) rather than by literally simulating the
Boltzmann distribution.

The question relevant to a fragment-based (rigid-template) docking model is
whether a bound (experimentally observed) pose is well approximated by
rigidly repositioning \emph{some} sample from $\mathcal M_c$: writing
$\psi(\bm\phi)$ for the dihedral-to-Cartesian reconstruction map and
$\Psi(p,R,\bm\phi)=(p,R)\cdot\psi(\bm\phi)$ for the composite (rigid motion
applied to an internally-reconstructed conformer), the relevant quantity is
\[
\min_{p\in\RR^3,\,R\in SO(3),\,\bm\phi\in\mathbb T^k}
\mathrm{RMSD}\bigl(x_b,\,\Psi(p,R,\bm\phi)\bigr),
\qquad
x_b\sim\pi_{\mathcal M_b}
\]
(the bound-pose distribution). This document does not re-derive this
empirical figure from scratch (it is a dataset-dependent, measured
quantity, not a theorem); we record it, as reported in the thesis draft
this monograph audits, and flag it as an empirical claim requiring the
cited dataset (Astex Diverse Set) to independently reproduce: the best of
five ETKDGv3 conformers aligns to the bound pose with mean RMSD
$0.21\,$\r A (and $0.39\,$\r A for a single conformer), well under the
$2\,$\r A success threshold of \Cref{def:rmsd}. The consequence for this
document: the roughly $0.2$\,\r A gap between a rigid-fragment
reconstruction and the true bound pose is an \textbf{error floor}
inherited by \emph{any} model built on this fragmentation, including
SigmaFlow, and is not attributable to the choice between diffusion and
flow matching --- it bounds what either method can, even in principle,
achieve, and should not be mistaken for a deficiency of one method
specifically when comparing the two (\Cref{chap:sigmaflow}).

\section{FR3D: reducing fragment count}
\label{sec:sigmadock-fr3d}

Cutting \emph{every} rotatable bond gives $\hat m=k+1$ fragments and
$6\hat m=6k+6$ degrees of freedom, exceeding the torsional
parameterisation's $k+6$ --- the naive fragment construction has bought
\Cref{thm:entanglement-scratch}'s decoupled measure at the cost of
\emph{more}, not fewer, dimensions than the alternative it replaces. FR3D
(Fragment REduction) recovers some of this cost back by merging adjacent
fragments: starting from the maximal $\hat m$-fragment decomposition, a
stochastic recursive procedure repeatedly merges two fragments across a
shared torsional bond, terminating at an irreducible fragment set of size
$m<\hat m$; empirically (per the audited thesis draft, not independently
re-derived here) $m\approx\tfrac23(k+1)$. Because which bonds get merged is
itself stochastic and (per the source description) equiprobable across
valid cut-sets, re-running the fragmentation on the same molecule at a
different training epoch can yield a genuinely different $m$-fragment
decomposition, functioning simultaneously as a data augmentation scheme.

\begin{caution}
The \emph{exact} merge heuristic and stopping criterion implemented by
FR3D were not confirmed line-by-line against
\file{fragmentation.py} in the code-research pass underlying this
document (the file is over $1500$ lines; only the rotatable-bond-detection
and triangulation-mapping functions, described precisely below and in
\Cref{sec:triangulation-precise}, were read to completion). What
\emph{was} confirmed directly is (i) rotatable-bond detection itself
(\Cref{sec:sigmadock-torsion-detection} below), (ii) that the file's own
top-of-module documentation describes a ``recursive merge'' procedure,
consistent with the description above, and (iii) the resulting dummy-atom
bookkeeping (\Cref{sec:sigmadock-fr3d-dummies} below). The precise
stochastic merge algorithm itself is therefore reported here as the
audited thesis draft's own description, not as independently re-derived
or line-verified fact, and is flagged accordingly.
\end{caution}

\subsection{Rotatable-bond detection, as actually implemented}
\label{sec:sigmadock-torsion-detection}

\begin{codebox}
\textbf{Mathematical object:} the set of rotatable bonds defining the cut
points of \Cref{sec:pose-dof} \\
\textbf{Implementation:} \file{fragmentation.py}::\pyfunc{detect\_torsional\_bonds} \\
\textbf{Interpretation:} uses RDKit's \pyfunc{TorsionFingerprints.CalculateTorsionLists}
to enumerate torsion quadruples, then filters to bonds that are: not part
of any ring, of bond type \texttt{SINGLE}, and (governed by an
\texttt{ignore\_conjugated} flag, default \texttt{False} per this
project's own training logs) optionally excluding conjugated bonds not
touching a ring. This is a specific, code-verified criterion, more
concrete than a generic ``rotatable bond'' description: not merely ``any
single, non-ring bond'', but specifically those bonds RDKit's own torsion
enumeration recognises as defining a well-formed dihedral quadruple.
\end{codebox}

\subsection{Dummy atoms and fragment boundaries}
\label{sec:sigmadock-fr3d-dummies}

Cutting a bond leaves each resulting fragment with a ``dangling'' bond at
the cut point; SigmaDock represents this by introducing a \emph{dummy
atom} on either side, so that a fragment's atom count can exceed the
corresponding subset of the original ligand's atom count
($|\mathcal G_{\mathrm{ligand}}|\le\sum_i|\mathcal G_{F_i}|$, one extra
dummy per retained cut). A dummy atom is retained only if it remains
\emph{free} (marks a torsional bond that FR3D's merge procedure did
\emph{not} collapse into a single fragment); if FR3D merges the two sides
of a torsional bond back together, the corresponding dummy is
\emph{pruned}, since retaining it would spuriously freeze a dihedral that
--- per \Cref{sec:conformational-manifold}'s reachability argument --- must
remain free for the rigid-fragment reconstruction to stay close to the true
bound-pose distribution.

\section{Triangulation: an invariant, soft geometric prior at fragment boundaries}
\label{sec:triangulation-precise}

Fragmentation is lossy in one specific, quantifiable sense: once two
fragments move independently, nothing in $SE(3)^m$ itself prevents the
bond angles \emph{across} the cut bond from taking chemically absurd
values (\Cref{chap:sigmaflow}'s empirical analysis of exactly this failure
mode, in the flow-matching setting, is one of the central quantitative
results of this project's broader investigation, referenced there).
SigmaDock's answer is not to hard-constrain the manifold (which would
reintroduce \Cref{thm:entanglement-scratch}'s coupling problems) but to
\emph{condition the network} on an elementary, exact piece of Euclidean
geometry.

\begin{theorem}[Triangulation determines the boundary bond angles]
\label{thm:triangulation-precise}
Let $\overline{BC}$ be a cut torsional bond, with $B$ in fragment
$\mathcal A\ni A$ and $C$ in fragment $\mathcal D\ni D$ ($A,D$ each bonded
to $B,C$ respectively, within their own fragment). Fixing the four
intra-fragment/bond lengths $\|A-B\|$, $\|B-C\|$, $\|C-D\|$ and the two
\emph{cross-fragment} distances $\|A-C\|$, $\|B-D\|$ determines the two
boundary bond angles $\angle(A,B,C)$, $\angle(B,C,D)$ uniquely, while
leaving the dihedral $\phi_{ABCD}$ completely unconstrained.
\end{theorem}

\begin{proof}
The triple $\{\|A-B\|,\|B-C\|,\|A-C\|\}$ determines the triangle $(A,B,C)$
up to rigid motion (the classical side-side-side congruence criterion), and
the law of cosines gives
$\cos\angle(A,B,C) = \dfrac{\|A-B\|^2+\|B-C\|^2-\|A-C\|^2}{2\|A-B\|\|B-C\|}$,
a function of the three fixed lengths alone; identically,
$\cos\angle(B,C,D)$ is determined by $\{\|B-C\|,\|C-D\|,\|B-D\|\}$. Neither
computation constrains the relative rotation of the two triangles' planes
about the shared edge $\overline{BC}$, which is exactly the dihedral
$\phi_{ABCD}$ by \Cref{sec:pose-dof}'s definition.
\end{proof}

The two distances the theorem needs, $\|A-C\|$ and $\|B-D\|$, are
manifestly $SE(3)$-invariant (ordinary Euclidean distances,
\Cref{prop:so3-isometry}), and hence admissible, per
\Cref{sec:invariance-obstruction}'s discussion, as ordinary type-$0$
(scalar) edge features consumable by an equivariant network without any
special handling.

\begin{codebox}
\textbf{Mathematical object:} the triangulation conditioning feature \\
\textbf{Implementation:} the deviation formula
$d_{ij}(t) = \bigl|\,\|x_i(t)-x_j(t)\| - \|x_i^{\mathrm{ref}}-x_j^{\mathrm{ref}}\|\,\bigr|$,
computed at \file{net/model.py}::\pyfunc{compute\_edge\_embedding}'s
``boundary'' edge branch, using \texttt{pos\_0} as the reference geometry
and \texttt{pos\_t} as the current (noised/interpolated) geometry; the
\emph{atom pairs} entering this formula are enumerated by
\file{fragmentation.py}::\pyfunc{get\_triangle\_equality\_mapping} \\
\textbf{Interpretation:} the network is given, at every step $t$, the
absolute deviation of each triangulation distance from its equilibrium
(reference-conformer) value, so $d_{ij}(t)\to0$ as $t\to1$ (approaching
data, in this document's generation-time convention) is exactly the signal
that ``this boundary is currently geometrically wrong, correct it''; no
projection or hard equality is ever imposed, only this soft conditioning
signal. \textbf{More general than \Cref{thm:triangulation-precise}'s
minimal statement:} the actual code maps every anchor atom to \emph{all}
atoms directly bonded to each of its neighbouring anchors (not merely one
representative atom $A$ or $D$ per side), so more than the minimal two
cross-fragment distances of the Theorem are typically supplied per cut
bond in practice --- the Theorem establishes that two such distances
\emph{suffice} in principle, not that exactly two are used. An
\texttt{ignore\_triangulation} ablation flag exists
(\file{data.py}), consistent with the audited thesis draft's report of an
ablation study in which removing this conditioning costs a documented
$12.8$ percentage points of PoseBusters validity (Def.~\ref{def:posebusters})
--- the single largest effect size reported in that study, though we have
not independently reproduced this specific number and report it as the
audited source's own claim.
\end{codebox}

\begin{remark}[Degrees of freedom of the soft-constrained manifold]
\label{rem:soft-constraint-dof}
Were the triangulation distances imposed as \emph{hard} equality
constraints (which they are not --- \Cref{thm:triangulation-precise}'s
feature is a soft conditioning signal, never a projection), the
constrained submanifold $\mathcal M_{\mathrm{frag}}=\{z\in SE(3)^m:c(z)=0\}$
would, for a tree-structured fragment graph ($e=m-1$ cut bonds), have two
non-degenerate distance constraints per joint removing five of the six
relative degrees of freedom between adjacent fragments (leaving exactly the
one free dihedral, as \Cref{thm:triangulation-precise}'s proof shows), so
$\dim\mathcal M_{\mathrm{frag}} = 6m-5(m-1) = m+5$ --- interpolating between
the torsional parameterisation's $k+6$ and the unconstrained $SE(3)^m$'s
$6m$. Since the actual constraints are soft, sampled trajectories only
\emph{concentrate} near $\mathcal M_{\mathrm{frag}}$ as $t\to1$; this
dimension count is a useful sanity check on the geometry, not a claim
about the trained network's exact behaviour.
\end{remark}

\section{Symmetry: gauge invariance and stochastic $SE(3)$-equivariance}
\label{sec:sigmadock-symmetry-precise}

Two distinct symmetry statements matter for SigmaDock, and conflating them
(as the broader literature sometimes does) obscures which architectural
property is responsible for which guarantee.

\begin{description}
\item[Gauge invariance.] A rigid fragment's \emph{local} coordinate frame
(the coordinate system its reference-conformer template is expressed in,
before any rigid motion is applied) has no canonical choice: for any fixed
$R_s\in SO(3)$, replacing the local template $\tilde x_F$ by
$R_s\tilde x_F$ and simultaneously replacing the fragment's rotation $R_F$
by $R_FR_s^\top$ describes the \emph{same} global pose,
$R_F\tilde x_F+p_F = (R_FR_s^\top)(R_s\tilde x_F)+p_F$. A correctly
constructed loss or sampler must not depend on this arbitrary choice.
\item[Stochastic $SE(3)$-equivariance.] The \emph{generated distribution}
should satisfy $p_\theta(dz\mid R_0\cdot\mathcal G_{\mathrm{dock}}) = R_0\cdot p_\theta(dz\mid\mathcal G_{\mathrm{dock}})$
for every $R_0\in SO(3)$: rotating the input pocket rotates the entire
predicted pose \emph{distribution} and nothing else about it.
\end{description}

\begin{theorem}[SigmaDock's symmetry]
\label{thm:sigmadock-symmetry}
SigmaDock's training objective and sampling procedure are invariant to the
(arbitrary) choice of fragment local-frame orientation, and the induced
kernel $p_\theta(z\mid\mathcal G_{\mathrm{dock}})$ is stochastically
$SO(3)$-equivariant.
\end{theorem}

This is the stochastic-process analogue of
\Cref{prop:equivariant-invariant-scratch}: there, a deterministic
equivariant ODE plus an invariant prior gave an invariant marginal law at
every time $t$; here, the corresponding statement for the \emph{diffusion
SDE} (\Cref{chap:diffusion,chap:diffusion-so3}) is that an $SE(3)$-equivariant
score network, together with the $SE(3)$-invariant prior
$p_{\mathrm{init}}=\mathcal N(0,I)\otimes\mathcal U(SO(3)^m)$ established in
\Cref{sec:sigmadock-forward} below, is sufficient (via the reverse-time SDE
of \Cref{eq:reverse-sde-scratch,eq:reverse-so3-scratch}, whose drift and
diffusion coefficients depend on the score in an equivariant way by
construction) for stochastic equivariance of the generated law; we do not
reprove this SDE-level statement in full technical generality (it requires
checking that the reverse-time SDE's law, not merely a deterministic flow's
pushforward, transforms correctly, a short but somewhat more technical
argument than \Cref{prop:equivariant-invariant-scratch}'s proof), and refer
to the general treatment in De Bortoli et al.\ (2022, cited in
\Cref{sec:manifold-brownian}) for the manifold-SDE case.
\emph{Gauge invariance} of the loss and sampler, by contrast, is a
structural property directly checkable from
\Cref{sec:equiformer-readout}'s observation that the network's backbone
consumes only world-frame Cartesian coordinates
$\mathbf x=\varphi(\mathbf T)$ and never the fragment-local decomposition
$(\tilde x_F,p_F,R_F)$ separately --- so its output is, by construction,
manifestly independent of any particular choice of $R_s$, since $R_s$ never
appears anywhere in the network's input.

\section{The forward (noising) process}
\label{sec:sigmadock-forward}

\begin{remark}[SigmaDock's own internal time convention, relabelled]
\label{rem:time-convention-scratch}
SigmaDock's own source code (\file{r3\_diffuser.py}, \file{so3\_diffuser.py})
runs its internal time variable in the \emph{opposite} direction from this
document's global convention (\Cref{rem:time-convention-global}): its
$t=0$ is clean data ($\alpha_{t=0}=1$, no noise) and its $t=1$ is the
fully-noised end (confirmed directly, e.g.\ by
\pyfunc{SO3Diffuser.sample\_ref} literally evaluating the forward kernel at
$t=1$ to obtain an approximately-prior sample). Every formula in this
chapter is stated in \emph{this document's} convention
($t=0$ noise, $t=1$ data, matching \Cref{rem:time-convention-global} and,
as \Cref{sec:sigmaflow-time-convention} confirms, matching SigmaFlow's own
native convention exactly), obtained from SigmaDock's native formulas by
the substitution $s=1-t$ throughout --- i.e.\ $s$ below always denotes
SigmaDock's own noising-time variable, and every SigmaDock formula quoted
in this chapter has already been relabelled via this substitution before
being stated, rather than being transcribed directly in SigmaDock's own
time variable.
\end{remark}

We can now assemble \Cref{chap:diffusion,chap:diffusion-so3}'s two
constructions, componentwise, on the state space
$\mathcal M_m=SE(3)^m$ fixed by \Cref{sec:sigmadock-state-space}, per
\Cref{prop:se3-product}'s licence to do so. Writing $s=1-t$ for noising
time (\Cref{rem:time-convention-scratch}) and suppressing the fragment
index:

\textbf{Translation}: a variance-preserving process (\Cref{ex:vp-sde}),
$\beta(s)=\beta_{\min}+s(\beta_{\max}-\beta_{\min})$ (linear schedule), so
the conditional law derived in \Cref{ex:vp-sde} is
$p_{s\mid0}(p^{(s)}\mid p^{\mathrm{data}}) = \mathcal N\bigl(\alpha_sp^{\mathrm{data}},(1-\alpha_s^2)I\bigr)$,
$\alpha_s=\exp\bigl(-\tfrac12\int_0^s\beta(r)\,dr\bigr)$.

\textbf{Rotation}: a variance-exploding process (\Cref{ex:ve-sde}), with
the \emph{logarithmic} schedule
$\sigma(s)=\log\bigl(se^{\sigma_{\max}}+(1-s)e^{\sigma_{\min}}\bigr)$
(confirmed exactly against \file{so3\_diffuser.py} in
\Cref{sec:diff-so3-numerics}), so the conditional law is
$p_{s\mid0}(R^{(s)}\mid R^{\mathrm{data}}) = \mathrm{IGSO}(3)(R^{\mathrm{data}},\sigma(s)^2)$
(\Cref{thm:igso3-scratch}).

By \Cref{prop:se3-product}'s product structure, the forward kernel
factorises exactly across the two components (and, trivially, across
fragments, by \Cref{def:pose-manifold-scratch}), and the stationary
distribution --- the model's prior --- is
$p_{\mathrm{init}} = \mathcal N(0,I)\otimes\mathcal U(SO(3)^m)$: an isotropic
Gaussian for translations (\Cref{prop:linear-sde-kernel}'s VP schedule
drives the marginal variance to $1$ as $s\to1$, for a suitably long/steep
$\beta$ schedule) and the Haar-uniform measure for rotations
(\Cref{prop:so3-brownian-stationary}), matching exactly the invariant prior
\Cref{thm:sigmadock-symmetry} requires.

\begin{remark}[Coordinate rescaling]
\label{rem:dimensional-scale}
Before any of the above is applied, ligand and protein coordinates are
centred on the binding-pocket centre of mass and divided by a fixed
constant, $\texttt{dimensional\_scale}=2.7$ (\file{oracle.py}); this same
constant additionally rescales \emph{every} graph-construction cutoff
radius used throughout \Cref{chap:equiformer} (protein--ligand interaction,
Ca--Ca virtual-node, and every other edge type's cutoff is divided by
$2.7$ before use, per \file{oracle.py}::\pyfunc{get\_edge\_specs}) --- i.e.\
this is a single, global \r Angstrom-to-model-units rescaling applied
uniformly throughout the pipeline, \emph{not} solely a device for matching
the translational Gaussian prior's unit variance to the data's spread, as
a narrower reading might suggest. Its specific numerical value ($2.7$) is
reported, in the audited thesis draft, as the mean standard deviation (in
\r A) of fragment centres of mass relative to the pocket centre of mass
across the training set --- an empirical, dataset-derived constant, not a
theoretically derived one.
\end{remark}

Training minimises the Riemannian denoising score matching objective of
\Cref{thm:dsm-equivalence,thm:rdsm-scratch}, with per-component weights
$\lambda_s$ matching \Cref{sec:sde-perspective}'s
$\lambda(s)=1/\EE[\|\nabla\log p_{s\mid0}\|^2]$ recipe: for translation,
this reduces to the closed form of
\Cref{sec:sde-perspective}'s discussion ($\lambda_s\propto1/(1-\alpha_s^2)$,
matching \pyfunc{R3Diffuser.score\_weight} exactly); for rotation, this
expectation has no closed form and is computed numerically from the
tabulated $\partial_\omega\log f_\sigma$ table
(\Cref{sec:diff-so3-numerics}'s \pyfunc{score\_scaling} code
correspondence). This machinery --- the $\mathrm{IGSO}(3)$ kernel with its
truncated spectral series, the conditional score derived from it, and the
tabulated loss weight --- is the \emph{entirety} of what is
diffusion-specific in SigmaDock; \Cref{chap:sigmaflow} shows precisely
which of these pieces flow matching removes and which it retains
unchanged.

\section{Input graph, featurisation, and backbone}
\label{sec:sigmadock-graph}

The network never sees the abstract group element $z\in SE(3)^m$ directly;
it sees the geometric graph of \Cref{def:geometric-graph-scratch}, built
from the current Cartesian coordinates $\mathbf x^{(t)}=\varphi(z^{(t)})$
(the rigid-fragment reconstruction of \Cref{sec:fragment-reconstruction}
below) and the (fixed) protein structure. \Cref{chap:equiformer} already
derived this construction in full architectural detail --- the categorical
atom/edge features and their dimensions
(\Cref{sec:node-embedding}'s code correspondence), the radial basis and
smooth cutoffs (\Cref{sec:gnn-scratch}'s code correspondence), the
hierarchical virtual-node structure (fragment-level and protein-C$\alpha$-level,
\Cref{sec:node-embedding}), and the triangulation edges just derived
(\Cref{sec:triangulation-precise}) --- and none of it is repeated here. We
record only the one remaining, purely diffusion-side, distinction: the
graph's \emph{topology} splits into a \emph{global} part, immutable along
the noising/generation trajectory (intra-fragment bonds, the protein
structure, the triangulation and virtual-node edges), and a
\emph{transient} part, recomputed at every step from the current, moving
positions $\mathbf x^{(t)}$ (the protein--ligand interaction radius graph,
$4\,$\r A cutoff per \Cref{sec:gnn-scratch}'s code correspondence): exactly
the distinction that made the smooth radial-cutoff envelope
$\chi$ of \eqref{eq:rbf-scratch} necessary, since an edge in the transient
part can appear or disappear as $\mathbf x^{(t)}$ evolves during either
training's noising or sampling's ODE/SDE integration.

The backbone is EquiformerV2 exactly as reconstructed in
\Cref{chap:equiformer}, with the one modification already noted there
(\Cref{sec:so2-conv}'s code correspondence): the radial MLPs generating
$SO(2)$-convolution weights on transient (protein--ligand interaction)
edges have their bias terms removed, which is what makes the smooth-cutoff
argument of \Cref{sec:gnn-scratch} hold exactly rather than only
approximately.

\subsection{Fragment reconstruction: the map $\varphi$}
\label{sec:fragment-reconstruction}

We finally make precise the map $\varphi:\mathcal M_m\to\RR^{n\times3}$,
used silently throughout this chapter, that turns an abstract group
element $z=(R^{(1)},p^{(1)},\dots)\in SE(3)^m$ into the atomic Cartesian
coordinates the network actually consumes.

\begin{proposition}[Rigid per-fragment reconstruction]
\label{prop:fragment-reconstruction}
Fix a reference conformer $x^{\mathrm{ref}}\in\RR^{n\times3}$ (the
rigid-fragment template of \Cref{sec:conformational-manifold}, one fixed
geometry per training/sampling instance, never itself diffused or flowed),
and let $p_F^{\mathrm{ref}}$ denote fragment $F$'s centre of mass in that
reference geometry. For a target pose $z=(R_F,p_F)_{F}$, the reconstructed
atom position, for atom $i$ in fragment $F$, is
\begin{equation}
\label{eq:fragment-reconstruction-formula}
\varphi(z)_i = R_F\bigl(x_i^{\mathrm{ref}} - p_F^{\mathrm{ref}}\bigr) + p_F^{\mathrm{ref}} + \Delta p_F,
\qquad
\Delta p_F := p_F - p_F^{\mathrm{ref}} ,
\end{equation}
i.e.\ every atom of fragment $F$ is displaced from its reference position
by exactly the same rigid motion: rotate about the fragment's own
reference centre of mass by $R_F$, then translate by $\Delta p_F$. Protein
atoms, having no fragment structure, are passed through unchanged.
\end{proposition}

\begin{proof}[Why this is the right formula]
This is, by construction, the unique map with two required properties: (i)
\emph{rigidity} --- for two atoms $i,j$ in the same fragment $F$,
$\|\varphi(z)_i-\varphi(z)_j\| = \|R_F(x_i^{\mathrm{ref}}-x_j^{\mathrm{ref}})\| = \|x_i^{\mathrm{ref}}-x_j^{\mathrm{ref}}\|$
by \Cref{prop:so3-isometry}, i.e.\ every intra-fragment distance --- in
particular every bond length and bond angle --- is preserved \emph{exactly},
for \emph{any} $(R_F,p_F)$, which is exactly the content of
\Cref{sec:conformational-manifold}'s claim that rigidity is enforced
structurally, not learned; and (ii) \emph{correct fragment placement} ---
setting $R_F=I,\Delta p_F=0$ recovers the reference conformer exactly, and
in general the fragment's centre of mass is placed at
$p_F^{\mathrm{ref}}+\Delta p_F=p_F$, matching the $z$-coordinate's own
translational component by construction.
\end{proof}

\begin{codebox}
\textbf{Mathematical object:} $\varphi$ of
\Cref{prop:fragment-reconstruction} \\
\textbf{Implementation:} \file{denoiser.py}::\pyfunc{\_apply\_transformations}
(SigmaDock) and, identically in form,
\file{sigma\_flow\_generator.py}::\pyfunc{get\_transformations\_from\_rototranslations}
(SigmaFlow, \Cref{chap:sigmaflow}) \\
\textbf{Tensor shape:} \shape{N, 3} $\to$ \shape{N, 3} \\
\textbf{Interpretation:} both modules implement the identical formula
$\texttt{new\_pos}[i] = R_F\cdot(\texttt{pos\_old}[i]-p_F^{\mathrm{ref}}) + p_F^{\mathrm{ref}} + \Delta p_F$,
confirming, at the level of this specific, central piece of geometric
plumbing, the claim (recorded again in \Cref{chap:sigmaflow}) that
SigmaFlow's conversion from SigmaDock leaves this component completely
unchanged: only the process that determines \emph{which} $(R_F,\Delta p_F)$
to evaluate $\varphi$ at, at a given step $t$, differs between the two
methods (\Cref{chap:diffusion,chap:diffusion-so3} vs.\
\Cref{chap:flow-matching}), never the reconstruction map itself.
\end{codebox}

\section{The Newton--Euler prediction head}
\label{sec:newton-euler-precise}

\Cref{sec:equiformer-readout} established that the network's raw output is
a per-atom $3$-vector in the world frame, read from the $\ell=1$ channels.
This section derives how that raw output becomes the two score components
$\eqref{eq:network-signature-scratch}$ requires --- and, in doing so,
explains precisely \emph{why} the code names its intermediate quantities
``force'', ``torque'', ``mass'', and ``inertia'', since this is not
decorative language but the literal reuse of classical rigid-body mechanics
formulas as an architectural device.

\subsection{From network output to a physical analogy}

\begin{codebox}
\textbf{Mathematical object:} the raw per-atom network output, interpreted
as a noise-prediction (an $\epsilon$-parameterisation target, in the
language of \Cref{sec:score-matching-intro}: recall
$x=z+\sigma\varepsilon$, so predicting $\varepsilon$ and predicting the
score $-\varepsilon/\sigma$ are related by a fixed, known rescaling) \\
\textbf{Implementation:} \file{denoiser.py}::\pyfunc{\_compute\_forces};
the network call itself is
\texttt{epsilon\_theta = self.model(batch, t, **model\_kwargs)}, followed
by \texttt{lig\_forces = -epsilon\_theta.index\_select(0, forces\_idxs)} \\
\textbf{Tensor shape:} \shape{N\_ligand\_atoms, 3} (the code's own inline
type annotation states \shape{B x F x A, 3, 3}, i.e.\ a full $3\times3$
matrix per atom rather than a $3$-vector; every downstream use --- indexing
along a single leading dimension, elementwise operations consistent with a
$3$-vector --- is inconsistent with that annotation, and we record it, per
this document's stated policy, as a probable stale/incorrect docstring
rather than a $3\times3$-matrix-valued quantity actually being used) \\
\textbf{Interpretation:} the variable name \texttt{lig\_forces}, and the
sign flip (\texttt{-epsilon\_theta}), signal the reinterpretation about to
happen: \textbf{this per-atom vector is about to be treated as a literal
force acting on that atom}, and everything from here to
\Cref{eq:sd-head-precise} below is the consequence of taking that
interpretation seriously and applying textbook rigid-body mechanics.
\end{codebox}

\subsection{Net force, torque, and inertia: exact classical mechanics}

\begin{definition}[Fragment mass, per this codebase's specific convention]
\label{def:fragment-mass-precise}
For fragment $F$ with (non-virtual, non-dummy) atoms indexed by
$i\in F$, the \emph{mass} is
$M_F := |F| = \sum_{i\in F}1$ --- the \textbf{atom count}, with every atom
assigned unit mass regardless of its chemical identity. This is
\textbf{not} physical/chemical atomic mass (e.g.\ carbon is not weighted
$12\times$ more than hydrogen); it is a uniform, purely topological
quantity.
\end{definition}

\begin{codebox}
\textbf{Mathematical object:} $M_F$ of \Cref{def:fragment-mass-precise}, and
the fragment inertia tensor
$I_F=\sum_{i\in F}\bigl(\|r_i\|^2I_3 - r_ir_i^\top\bigr)$, $r_i=x_i-p_F$
(position relative to the fragment centre of mass) --- the classical
inertia tensor of a rigid body composed of point masses, specialised to
unit mass per point \\
\textbf{Implementation:} \file{denoiser.py}::\pyfunc{get\_fragment\_mass\_inertia} \\
\textbf{Interpretation:} the assertion \texttt{assert M.min() >= 2} (with
message \texttt{"Fragment mass below 2 detected. Check input data."}) is
now precisely explained: it fires whenever a fragment has \emph{fewer than
two atoms}, a purely topological/fragmentation degeneracy (a fragment
reduced, by FR3D's merging or by the underlying molecule's own structure,
to a single heavy atom plus dummies), not a statement about any atom's
chemical mass being unusually low --- the theory-document text of
\Cref{chap:sigmadock} must therefore never describe this condition as
``fragment mass'' in the chemical sense, only in the atom-count sense
\Cref{def:fragment-mass-precise} fixes. The inertia tensor's eigenvalues
are additionally floored (via an eigendecomposition, clamping to at least
$\max(10^{-8}\tr(I)/3,\,10^{-12})$) before use, protecting the subsequent
linear solve (\Cref{eq:newton-maruyama-precise} below) against
near-degenerate (e.g.\ near-linear, two-to-three-atom) fragments.
\end{codebox}

Given the per-atom pseudo-forces $f_i$ (the negated network output) and the
mass/inertia just defined, the exact classical-mechanics formulas for a
rigid body's net force, net torque about its centre of mass, are:
\begin{equation}
\label{eq:newton-euler-precise}
F_F = \sum_{i\in F}f_i,
\qquad
\tau_F = \sum_{i\in F}(x_i-p_F)\times f_i .
\end{equation}

\begin{codebox}
\textbf{Mathematical object:} \eqref{eq:newton-euler-precise} \\
\textbf{Implementation:} \file{denoiser.py}::\pyfunc{linear\_mechanics}
(a \texttt{scatter}/\texttt{index\_add}-based per-fragment reduction of a
per-atom force field, computed at the \emph{current}, time-$t$ fragment
centres of mass, i.e.\ using $\mathtt{coms}=\mathtt{trans\_t}$, not the
data-time centres) \\
\textbf{Interpretation:} this is, literally, the textbook definition of net
force ($\sum$ of per-particle forces) and net torque about a pivot
($\sum r_i\times f_i$) for a rigid body under a field of per-particle
forces --- the ``pseudo-force''/``Newton--Euler'' naming is earned by
direct, unmodified reuse of these formulas, with the network's noise
prediction standing in for the physical force field.
\end{codebox}

\subsection{Newton--Maruyama: one discretised step of rigid-body dynamics}

\begin{equation}
\label{eq:newton-maruyama-precise}
dT = \frac{F_F}{M_F}
\qquad\text{(force over mass = acceleration-like increment)},
\qquad
dW = I_F^{-1}\tau_F
\qquad\text{(Euler's rigid-body equation, solved for angular increment)} ,
\end{equation}
with $dW\in\RR^3$ then lifted to $\so(3)$ via $\widehat{dW}$
(\Cref{def:hatvee-scratch}).

\begin{codebox}
\textbf{Mathematical object:} \eqref{eq:newton-maruyama-precise} \\
\textbf{Implementation:} \file{denoiser.py}::\pyfunc{newton\_maruyama}
(\texttt{dT = force\_per\_fragment / mass}; \texttt{dW\_world = torch.linalg.solve(inertia, torque\_per\_fragment)},
with the inertia additionally Tikhonov-regularised by adding
$10^{-8}I$ on top of the eigenvalue floor of
\Cref{def:fragment-mass-precise}'s code correspondence; then
\texttt{omega = hat(dW\_world)}, clamped entrywise to $[-10^3,10^3]$ as a
loose numerical safety net) \\
\textbf{Interpretation:} the function name ``Newton--Maruyama'' fuses
Newton's second law (the $dT=F/M$ step) with the Euler--Maruyama
discretisation scheme referenced throughout
\Cref{chap:diffusion,chap:diffusion-so3} (this specific function itself is
purely deterministic given force/torque --- the stochastic Maruyama noise
increment is added \emph{separately}, at the reverse-SDE-integration step
of \Cref{sec:sigmadock-sampling} below, not here). The code's own inline
comments assert this quantity is expressed in the ``GLOBAL (left-invariant
--- WORLD) frame'', which we read, per
\Cref{sec:so3-trivialisation}'s resolved terminology, as the
\textbf{world-frame (left-trivialised)} angular velocity
$\omega^{\mathrm L}$ of \Cref{def:trivialisation-scratch}: no rotation by
$R_t$ is applied to $dW$ before or after the linear solve, consistent with
this reading, and consistent with $I_F$ itself having been computed at the
current \emph{world-frame} positions $x_i$ (not in any fragment-local
frame). This is worth flagging explicitly against
\Cref{sec:so3-trivialisation}'s resolution for the \emph{flow-matching}
side, which found the \emph{opposite} (right-trivialised/body-frame)
convention in \pyfunc{so3\_flow\_matcher.py}: SigmaDock's Newton--Euler
head and SigmaFlow's flow-matching update operate in different
trivialisation conventions at this specific point, a genuine architectural
difference (not necessarily an inconsistency, since each is internally
self-consistent within its own module, by \Cref{sec:so3-trivialisation}'s
verification for the flow-matching side and the world-frame reading just
given here for the diffusion side) that \Cref{chap:sigmaflow} discusses
further when it reconstructs the corresponding flow-matching quantity.
\end{codebox}

\subsection{Dressing the raw output into a score: the two branches}
\label{sec:sd-head-branches}

The pair $(dT,dW)\in\RR^3\times\RR^3$ from
\eqref{eq:newton-maruyama-precise} already has the correct \emph{shape} of
a tangent vector on $\mathcal M_m$ (\eqref{eq:network-signature-scratch}),
but it is not yet a score: \Cref{prop:igso3-score-scratch} and
\Cref{rem:gaussian-score} show the true conditional score carries a
noise-level-dependent scalar prefactor that this raw mechanical quantity
does not yet include. SigmaDock dresses the translational component in one
fixed way, and the rotational component via a code-configurable choice of
\emph{two} distinct branches.

\textbf{Translation} (single formula, no branching):
\begin{equation}
\label{eq:sd-head-precise}
s^p_\theta = r_3\text{-}\mathrm{scaling}(t)\cdot dT,
\qquad
r_3\text{-}\mathrm{scaling}(t) = \texttt{R3Diffuser.score\_scaling}(t) ,
\end{equation}
matching \Cref{rem:gaussian-score}'s Gaussian-score form up to the fixed,
schedule-determined scalar prefactor of \Cref{sec:sde-perspective}.

\textbf{Rotation}, branch \texttt{rot\_score\_method="space"} (constructor
default): treat $dW$ as an \emph{implied final-rotation update} rather than
a score directly --- form $\widehat{dW}$, rescale by the $SO(3)$-analogue
of the translational scaling (\texttt{so3\_scaling}, from
\pyfunc{SO3Diffuser.score\_scaling}), exponentiate
($\Delta R=\exp(\text{scaled }\widehat{dW})$, \Cref{sec:rodrigues-scratch}),
compose $\widehat R_0 := \Delta R\cdot R_t$ (an implied estimate of the
clean/data rotation), and \emph{then} evaluate the exact closed-form
conditional score \eqref{eq:igso3-score-scratch} of
\Cref{prop:igso3-score-scratch} at this implied $\widehat R_0$:
\begin{equation}
\label{eq:sd-head-rot-space}
s^R_\theta = \nabla_{R_t}\log p_{t\mid1}(R_t\mid\widehat R_0)
\qquad\text{(\Cref{eq:igso3-score-scratch}, evaluated at }R_1=\widehat R_0\text{).}
\end{equation}

\textbf{Rotation}, branch \texttt{rot\_score\_method="score"} (raises a
configuration error unless a scaling sub-choice is also specified): treat
$dW$ as \emph{directly proportional to} the score itself,
\begin{equation}
\label{eq:sd-head-rot-score}
s^R_\theta = -\alpha(t)\cdot\widehat{dW}\cdot R_t ,
\end{equation}
with the scalar $\alpha(t)$ chosen by a further sub-choice: either
\texttt{"rms"} (reuse the same \texttt{so3\_scaling} used by the
\texttt{"space"} branch) or \texttt{"true"} (evaluate
$\partial_\omega\log f_{\sigma(t)}(\omega)/\omega$ live, at the network's
own implied angle $\omega=\|dW\|$, from
\Cref{sec:diff-so3-numerics}'s \pyfunc{d\_log\_f\_d\_omega} routine ---
i.e.\ the \emph{exact} closed-form score-scaling of
\Cref{prop:igso3-score-scratch} evaluated at whatever angle the network
happens to have produced, rather than the RMS-averaged $\sigma(t)$-only
scaling of the \texttt{"rms"} option).

\begin{caution}[This document's confirmed configuration]
Per the code-research pass underlying this document, the constructor
default is \texttt{rot\_score\_method="space"},
\texttt{rot\_score\_scaling="rms"}. Which branch the specific trained
checkpoints referenced elsewhere in this project's broader development
history actually used at training time was not independently
cross-checked against those checkpoints' saved hyperparameters as part of
producing this document, and should be verified separately before drawing
any conclusion that depends on the specific branch in force.
\end{caution}

\section{The training loss}
\label{sec:sigmadock-loss}

\begin{codebox}
\textbf{Mathematical object:} the four-component SigmaDock training loss \\
\textbf{Implementation:} \file{denoiser.py}::\pyfunc{compute\_losses} \\
\textbf{Interpretation:} four terms, each a $\lambda$-weighted squared
error, matching \eqref{eq:dsm-continuous}'s continuous-time DSM objective
pattern:
\begin{align*}
\text{(score-matching, translation)}&\quad
\lambda^p_t\,\bigl\|s^p_\theta - s^p_{\mathrm{true}}\bigr\|_2^2,
\qquad
\lambda^p_t = 1/(\text{T\_score\_scaling})^2,\\
\text{(score-matching, rotation)}&\quad
\lambda^R_t\,\bigl\|s^R_\theta - s^R_{\mathrm{true}}\bigr\|_{\mathrm{Fro}}^2,
\qquad
\lambda^R_t = 1/(\text{R\_score\_scaling})^2,\\
\text{(data-space, translation)}&\quad
\lambda^{p,0}_t\,\|p_0-\widehat p_0\|_2^2,
\qquad
\lambda^{p,0}_t = (\alpha_t\cdot\text{T\_score\_scaling})^2,\\
\text{(data-space, rotation)}&\quad
\lambda^{R,0}_t\,\bigl\|\log(\widehat R_0^\top R_0)\bigr\|_{\mathrm{Fro}}^2,
\qquad
\lambda^{R,0}_t = (\text{R\_score\_scaling})^2 ,
\end{align*}
returned as a four-key dictionary $\{\texttt{T\_score},\texttt{R\_score},\texttt{T0},\texttt{R0}\}$.
\end{codebox}

\begin{codebox}
\textbf{Mathematical object:} per-molecule reduction of the (per-fragment)
losses above \\
\textbf{Implementation:} \file{denoiser.py}::\pyfunc{scaled\_fragmented\_loss}:
\texttt{scatter(x, batch\_idx, reduce="sum") / num\_fragments\textasciicircum{}fragment\_scaling} \\
\textbf{Interpretation:} \texttt{fragment\_scaling}$=1$ divides by the exact
fragment count, i.e.\ averages per molecule regardless of how many
fragments it has; \texttt{fragment\_scaling}$=-1$ (per one code comment
site's own docstring) applies \emph{no} normalisation, biasing the loss
towards molecules with more fragments; intermediate values interpolate.
\textbf{Note, per the code-research pass:} the default value of this
hyperparameter differs between two source locations in the codebase
(\texttt{1.0} in one configuration dataclass, \texttt{0.5} in the Lightning
module's own constructor default) --- almost certainly immaterial in
practice if the training script's CLI parsing always threads an explicit
value through, but this was not independently confirmed, and is recorded
as an open item rather than resolved by assumption.
\end{codebox}

\section{Sampling}
\label{sec:sigmadock-sampling}

Generation integrates the reverse-time SDEs of
\Cref{eq:reverse-sde-scratch,eq:reverse-so3-scratch}, componentwise, from a
near-$t=0$ starting time to $t=1$.

\begin{codebox}
\textbf{Mathematical object:} the timestep discretisation of the reverse
integration \\
\textbf{Implementation:} \file{sampling.py}::\pyfunc{sampler}, matching
\file{conf/sampling/base.yaml}'s \texttt{diffusion:} block \\
\textbf{Interpretation:} default \texttt{num\_steps}$=25$,
\texttt{t\_min}$=0.005$, \texttt{rho}$=3.0$, default
\texttt{discretization="edm"} (the Karras et al.\ power-law-in-$1/\rho$
schedule, $t_k=\bigl(t_{\max}^{1/\rho}+\tfrac k{N-1}(t_{\min}^{1/\rho}-t_{\max}^{1/\rho})\bigr)^\rho$,
rather than the simpler \texttt{"power"} schedule
$t_k=\mathrm{linspace}(t_{\max},t_{\min},N)^\rho$, which is available but
not the default here); the timestep array itself walks \emph{downward}, in
this module's own internal (noising-time-aligned) convention, from
$t_{\max}=1$ to $t_{\min}$, consistent with \Cref{rem:time-convention-global}'s
warning that SigmaDock's native time convention is reversed relative to
this document's global convention (\Cref{chap:sigmaflow} makes the exact
reversal explicit and code-confirmed).
\end{codebox}

\begin{codebox}
\textbf{Mathematical object:} one reverse-time step (Euler--Maruyama for
translation, geodesic random walk for rotation) \\
\textbf{Implementation:} \file{sampling.py}::\pyfunc{sampler}'s loop body,
composing \pyfunc{denoiser.\_compute\_forces} $\to$
\pyfunc{\_compute\_fragment\_dynamics} $\to$
\pyfunc{\_predict\_fragment\_updates} (Newton--Maruyama,
\Cref{eq:newton-maruyama-precise}) $\to$ \pyfunc{\_compute\_scores}
(\Cref{sec:sd-head-branches}'s branch) $\to$
\pyfunc{denoiser.diffuser.reverse} (dispatching componentwise to
\pyfunc{R3Diffuser.reverse}, \Cref{ex:vp-sde}'s code correspondence, and
\pyfunc{SO3Diffuser.reverse}, \Cref{sec:diff-so3}'s code correspondence)
$\to$ \pyfunc{\_apply\_transformations} (rigid-fragment position
reconstruction, \Cref{sec:fragment-reconstruction}) \\
\textbf{Interpretation:} exactly one EquiformerV2 forward pass per
integration step --- the same cost structure
\Cref{chap:flow-matching,chap:sigmaflow} will show flow matching shares,
so any difference in the number of network evaluations needed for
comparable sample quality between the two methods is a genuine, comparable
efficiency question, examined in this project's broader experimental
work rather than in this theory document. \textbf{Heun's method is not
implemented}, despite \texttt{solver:\ Literal["euler","heun"]} appearing
in the function signature: both this sampler and the flow-matching
sampler of \Cref{chap:sigmaflow} contain only a dead
\texttt{\# TODO Heun's method} comment block after the (Euler-only) loop;
selecting \texttt{solver="heun"} has, per the code-research pass
underlying this document, no effect on the executed computation.
\end{codebox}

\begin{codebox}
\textbf{Mathematical object:} the stochastic noise-annealing schedule \\
\textbf{Implementation:} quadratic decay
$\texttt{noise\_scales} = \mathrm{linspace}\bigl(\texttt{noise\_scale}^{1/\texttt{noise\_decay}},\,0,\,\texttt{num\_steps}\bigr)^{\texttt{noise\_decay}}$,
default \texttt{noise\_scale}$=0.0$, \texttt{noise\_decay}$=2.0$ (per
\file{conf/sampling/base.yaml}) \\
\textbf{Interpretation:} this array rescales \emph{only} the injected
Brownian-increment term of the Euler--Maruyama/geodesic-random-walk update
(the $\sqrt{dt}\,z$ term), independently of the score-drift term, and
decays from \texttt{noise\_scale} to exactly $0$ over the course of
sampling. \textbf{With the shipped default \texttt{noise\_scale=0.0},
every injected-noise term is zero at every step}: SigmaDock's
\emph{out-of-the-box default sampling configuration is already fully
deterministic}, i.e.\ it already runs, without any code change, the
probability-flow-ODE-equivalent trajectory of \Cref{thm:pf-ode-scratch}
rather than a genuinely stochastic SDE trajectory. This narrows the
``diffusion is stochastic, flow matching is deterministic'' framing
considerably: the distinction that is architecturally forced for flow
matching (\Cref{chap:flow-matching} never constructs a noising process at
all) is, for this specific configuration of SigmaDock, an \emph{opt-in}
choice (setting \texttt{noise\_scale}$>0$) rather than an always-active
default behaviour.
\end{codebox}

\section{Ranking generated poses}
\label{sec:sigmadock-ranking}

SigmaDock draws $N_{\mathrm{seeds}}$ (default $40$, per the audited thesis
draft) independent samples per complex and ranks them by a fixed, learned-model-free
heuristic,
\[
s_i = -b_i\,p_i^\beta,
\qquad
\beta=4,
\]
where $b_i$ is a Vinardo docking score (a classical, physics-based scoring
function, not part of the generative model itself) and
$p_i\in[0,1]$ is the fraction of PoseBusters checks
(\Cref{def:posebusters}) sample $i$ passes. SigmaDock deliberately
dispenses with both a learned confidence/ranking model and any post-hoc
coordinate-altering energy minimisation --- both standard in the
DiffDock lineage this document's introduction references --- specifically
so that reported accuracy can be attributed to the generative model itself
rather than to a downstream correction step; \Cref{sec:likelihood-confidence}
in \Cref{chap:sigmaflow} returns to whether flow matching's ODE structure
offers a principled alternative to this heuristic.

\clearpage

% ================================================================
\part{Flow Matching}
\label{part:flow-matching}
% ================================================================

\chapter{Continuous Normalizing Flows and Flow Matching in $\RR^d$}
\label{chap:flow-matching}

\Cref{sec:pf-ode-scratch} showed, as a \emph{consequence} of diffusion
theory, that every score-based diffusion model secretly transports
$p_{\mathrm{init}}$ to $p_{\mathrm{data}}$ along a deterministic ODE. This
chapter develops the same object --- a deterministic flow transporting one
distribution to another --- as a \emph{starting point} rather than a
derived fact, entirely independently of scores, noising processes, or
SDEs. The resulting framework, flow matching, is what \Cref{chap:sigmaflow}
substitutes for \Cref{chap:diffusion,chap:diffusion-so3} to obtain
SigmaFlow.

\section{Continuous normalizing flows}
\label{sec:cnf}

\begin{definition}[Time-dependent vector field, flow map]
\label{def:vector-field-flow}
A \emph{time-dependent vector field} is a map $u:\RR^d\times[0,1]\to\RR^d$.
Given $u$, the associated \emph{flow map} $\phi_t:\RR^d\to\RR^d$ is defined
by solving, for each starting point $x$,
\[
\frac{d}{dt}\phi_t(x) = u_t(\phi_t(x)),
\qquad
\phi_0(x)=x ,
\]
i.e.\ $\phi_t(x)$ is the position at time $t$ of the particle that started
at $x$ and has been carried along by the velocity field $u$ ever since. We
assume throughout (as is standard, and guaranteed by, e.g., $u_t$ globally
Lipschitz in $x$ uniformly in $t$) that this ODE has a unique solution for
every $x$ and every $t\in[0,1]$, so that $\phi_t$ is well defined and, since
distinct trajectories of an autonomous-in-space-at-each-fixed-$t$ ODE
cannot cross without violating uniqueness, a diffeomorphism of $\RR^d$ for
each fixed $t$.
\end{definition}

\begin{definition}[Continuous normalizing flow]
\label{def:cnf}
Given $u$ and a source density $p_{\mathrm{init}}$, the
\emph{continuous normalizing flow (CNF)} associated with $u$ is the
probability path $p_t := (\phi_t)_\#p_{\mathrm{init}}$, the pushforward
(\Cref{def:pushforward}) of $p_{\mathrm{init}}$ under the flow map.
\end{definition}

Equivalently, and operationally: draw $X_0\sim p_{\mathrm{init}}$, solve
$\dot X_t=u_t(X_t)$ forward to $t=1$; then $X_1\sim p_1$. This is precisely
\Cref{def:vector-field-flow}'s flow map restated with full definitional
care, and it is worth being explicit about what has \emph{not} yet been
claimed:
nothing so far says $p_1=p_{\mathrm{data}}$ for any \emph{particular}
choice of $u$, only that \emph{some} probability path $p_t$ is generated by
\emph{any} choice of (well-behaved) $u_t$. The next result is the tool that
lets us go the other way: given a \emph{desired} path $p_t$, characterise
which vector fields generate it, so that we may search for $u$ within that
characterisation rather than by trial and error.

\subsection{The continuity equation}

\begin{theorem}[Continuity equation]
\label{thm:continuity}
Let $u$ be a time-dependent vector field with flow $\phi_t$
(\Cref{def:vector-field-flow}), and let $p_t=(\phi_t)_\#p_{\mathrm{init}}$
(\Cref{def:cnf}). Then, wherever $p_t$ is differentiable,
\begin{equation}
\label{eq:continuity}
\partial_tp_t(x) + \nabla\cdot\bigl(u_t(x)p_t(x)\bigr) = 0 ,
\qquad
\nabla\cdot v := \sum_{k=1}^d\partial_{x_k}v_k \ (\text{divergence}) .
\end{equation}
Conversely, if a smooth, compactly-supported-density path $(p_t)$ and
vector field $u_t$ satisfy \eqref{eq:continuity}, then $u_t$ generates
$(p_t)$ in the sense of \Cref{def:cnf} (given standard regularity
sufficient for existence/uniqueness of the flow, \Cref{def:vector-field-flow}).
\end{theorem}

\begin{proof}
(Forward direction.) By \Cref{thm:change-of-variables} applied to the
diffeomorphism $\phi_t$,
\[
p_t(x) = p_{\mathrm{init}}\bigl(\phi_t^{-1}(x)\bigr)\,\bigl|\det D\phi_t^{-1}(x)\bigr| .
\]
Write $J_t(x) := \det D\phi_t(\phi_t^{-1}(x))$, so
$\bigl|\det D\phi_t^{-1}(x)\bigr| = 1/J_t(x)$ (inverse function theorem, and
$J_t>0$ since $\phi_t$ is orientation-preserving as the flow of a
continuous vector field starting at the identity, $\phi_0=\mathrm{id}$,
$\det D\phi_0=1>0$, and $J_t$ cannot cross zero continuously without the
flow ceasing to be a diffeomorphism). We use \emph{Liouville's formula} for
the rate of change of this Jacobian along the flow,
\begin{equation}
\label{eq:liouville}
\frac{d}{dt}\log J_t(\phi_t(x_0)) = (\nabla\cdot u_t)\bigl(\phi_t(x_0)\bigr) ,
\end{equation}
which we justify as follows: differentiating $J_t(\phi_t(x_0))=\det D\phi_t(x_0)$
in $t$, using the matrix identity $\frac{d}{dt}\det A(t)=\det A(t)\,\tr\bigl(A(t)^{-1}\dot A(t)\bigr)$
(Jacobi's formula, a standard fact of matrix calculus we take as given)
with $A(t)=D\phi_t(x_0)$, and $\dot A(t) = D\bigl(\partial_t\phi_t(x_0)\bigr) = D\bigl(u_t(\phi_t(x_0))\bigr) = Du_t(\phi_t(x_0))\cdot D\phi_t(x_0) = Du_t(\phi_t(x_0))\,A(t)$
(chain rule, differentiating the defining ODE of \Cref{def:vector-field-flow}
with respect to the initial condition $x_0$, and using that spatial and
initial-condition derivatives commute for a smooth flow), gives
$\tr(A(t)^{-1}\dot A(t)) = \tr\bigl(Du_t(\phi_t(x_0))\bigr) = (\nabla\cdot u_t)(\phi_t(x_0))$
(the trace of the Jacobian matrix of a vector field is, by definition, its
divergence), yielding \eqref{eq:liouville} directly.

Now compute $\partial_tp_t(x)$ along a \emph{fixed} $x$ by relating it to
the material (Lagrangian) derivative along the flow. Write
$q(t) := p_t(\phi_t(x_0))\,J_t(\phi_t(x_0))$ for fixed $x_0$; since
$p_t(\phi_t(x_0)) = p_{\mathrm{init}}(x_0)/J_t(\phi_t(x_0))$
(the change-of-variables formula above, evaluated at $x=\phi_t(x_0)$, using
$\phi_t^{-1}(\phi_t(x_0))=x_0$), we get $q(t)=p_{\mathrm{init}}(x_0)$,
constant in $t$. Differentiating $q(t)=p_t(\phi_t(x_0))J_t(\phi_t(x_0))$
using the product rule and the chain rule (treating $p_t(\phi_t(x_0))$'s
time-derivative as $\partial_tp_t(\phi_t(x_0)) + \nabla p_t(\phi_t(x_0))\cdot u_t(\phi_t(x_0))$,
the material derivative), and using \eqref{eq:liouville} for
$\dot J_t(\phi_t(x_0))=J_t(\phi_t(x_0))(\nabla\cdot u_t)(\phi_t(x_0))$, gives,
after dividing by $J_t(\phi_t(x_0))\ne0$ and evaluating at the (arbitrary)
point $x=\phi_t(x_0)$,
\[
0 = \partial_tp_t(x) + \nabla p_t(x)\cdot u_t(x) + p_t(x)\,(\nabla\cdot u_t)(x)
= \partial_tp_t(x) + \nabla\cdot\bigl(p_t(x)u_t(x)\bigr) ,
\]
using the product rule for divergence,
$\nabla\cdot(pu)=\nabla p\cdot u+p\,\nabla\cdot u$, in the last step. This is
\eqref{eq:continuity}.

(Converse direction.) We do not prove the converse in full generality here
(it is a uniqueness statement for the linear transport PDE
\eqref{eq:continuity}, following from the method of characteristics: any
two solutions of \eqref{eq:continuity} with the same initial condition
$p_0$ agree, because \eqref{eq:continuity} is exactly the statement that
mass is conserved along the characteristics $\dot x=u_t(x)$, so the unique
solution is obtained by transporting $p_0$'s mass along those
characteristics, which is precisely $(\phi_t)_\#p_0$); see, e.g., Villani,
\emph{Optimal Transport: Old and New}, 2009, \S4, or Chen et al.,
\emph{Neural Ordinary Differential Equations}, NeurIPS 2018, \S2 for the
specific CNF framing.
\end{proof}

\Cref{thm:continuity} is the ODE analogue of a Fokker--Planck equation
(indeed \Cref{thm:pf-ode-scratch}'s proof sketch used exactly this fact,
citing this theorem for the diffusionless special case that appears when
the SDE's diffusion term is cancelled by the specific drift construction
there), and it is the tool that converts ``I want to transport
$p_{\mathrm{init}}$ to $p_{\mathrm{data}}$'' into a concrete, checkable
condition on $u_t$: any $(p_t,u_t)$ pair satisfying
\eqref{eq:continuity}, with $p_0=p_{\mathrm{init}}$, generates a CNF
ending, at $t=1$, at whatever $p_1$ the pair specifies.

\subsection{CNF likelihood}
\label{sec:cnf-likelihood}

A byproduct of \Cref{thm:continuity}'s proof is worth recording separately,
since \Cref{sec:likelihood-confidence} in \Cref{chap:sigmaflow} needs it.

\begin{corollary}[Instantaneous change of variables]
\label{cor:instantaneous-cov}
Along a trajectory $X_t=\phi_t(X_0)$ of the flow, the log-density evolves by
\[
\frac{d}{dt}\log p_t(X_t) = -(\nabla\cdot u_t)(X_t) .
\]
\end{corollary}

\begin{proof}
This is exactly \eqref{eq:liouville} combined with
$p_t(\phi_t(x_0))=p_{\mathrm{init}}(x_0)/J_t(\phi_t(x_0))$ from the proof of
\Cref{thm:continuity}: taking $\log$ of both sides,
$\log p_t(\phi_t(x_0)) = \log p_{\mathrm{init}}(x_0) - \log J_t(\phi_t(x_0))$,
and differentiating in $t$ using \eqref{eq:liouville} for the second term
gives the claim directly.
\end{proof}

\Cref{cor:instantaneous-cov} says that, given the vector field $u_t$
\emph{and} the initial log-density $\log p_{\mathrm{init}}(X_0)$, the
log-density of the sample $X_1=\phi_1(X_0)$ can be recovered \emph{exactly}
by integrating the (scalar!) divergence $-\nabla\cdot u_t$ along the same
trajectory the sample itself follows --- no separate, learned density
network is needed, only the already-learned vector field and one
additional scalar ODE run alongside the original one. Evaluating
$\nabla\cdot u_t$ exactly requires the trace of a $d\times d$ Jacobian,
$\mathcal O(d)$ backward passes if done naively; in practice this is
estimated stochastically via the \emph{Hutchinson trace estimator},
\[
\nabla\cdot u_t(x) = \tr\bigl(Du_t(x)\bigr) = \EE_{\varepsilon}\bigl[\varepsilon^\top Du_t(x)\varepsilon\bigr]
\]
for any $\varepsilon$ with $\EE[\varepsilon\varepsilon^\top]=I$ (e.g.\
$\varepsilon\sim\mathcal N(0,I)$ or Rademacher), computable with a single
Jacobian-vector product per sample of $\varepsilon$ rather than a full
Jacobian; we record this identity (a direct consequence of linearity of
trace and expectation, $\EE[\varepsilon^\top A\varepsilon]=\EE[\tr(A\varepsilon\varepsilon^\top)]=\tr(A\EE[\varepsilon\varepsilon^\top])=\tr(A)$)
without further development, since \Cref{sec:likelihood-confidence}
records only whether this machinery is implemented for SigmaFlow, not an
independent numerical study of it.

\section{Conditional flow matching}
\label{sec:cfm-scratch}

\Cref{thm:continuity} characterises which $u_t$ generate a \emph{given}
path $(p_t)$, but a path connecting an entire distribution
$p_{\mathrm{init}}$ to an entire distribution $p_{\mathrm{data}}$ is
generally just as intractable to write down directly as the marginal score
of \Cref{sec:score-matching-intro} was. The resolution is, again,
\Cref{prop:marginalisation}'s conditioning trick, now applied to
probability paths instead of static densities --- and the reader will
notice the entire structure of this section mirrors
\Cref{sec:score-matching-intro}'s denoising-score-matching argument almost
verbatim, which is not a coincidence but the content of \Cref{sec:pf-ode-scratch}'s
``diffusion and flow matching are two parameterisations of one object''
claim, now demonstrated by the parallel structure of the two derivations
rather than merely asserted.

\begin{definition}[Conditional probability path, conditional vector field]
\label{def:cond-path-cfm}
A \emph{conditional probability path} is a family $p_t(\cdot\mid z)$,
$z\in\RR^d$, with $p_0(\cdot\mid z)=p_{\mathrm{init}}$,
$p_1(\cdot\mid z)=\delta_z$ (a Dirac mass at $z$) for every $z$. A
\emph{conditional vector field} $u_t(\cdot\mid z)$ generates
$p_t(\cdot\mid z)$ in the sense of \Cref{thm:continuity}, for each fixed
$z$.
\end{definition}

By \Cref{prop:marginalisation}, the \emph{marginal} path
$p_t(x) := \int p_t(x\mid z)p_{\mathrm{data}}(z)\,dz$ automatically
satisfies $p_0=p_{\mathrm{init}}$ (each conditional starts at
$p_{\mathrm{init}}$, so their $p_{\mathrm{data}}$-weighted mixture does
too) and $p_1=p_{\mathrm{data}}$ ($p_1(x)=\int\delta_z(x)p_{\mathrm{data}}(z)dz=p_{\mathrm{data}}(x)$).
The following is the flow-matching analogue of
\Cref{thm:marginalisation-vf-scratch}, now proved (main.tex, the audited draft,
stated this without proof; we supply one, since it is short and instructive).

\begin{theorem}[Marginalisation of vector fields]
\label{thm:marginalisation-vf-scratch}
Let $u_t(\cdot\mid z)$ generate the conditional path $p_t(\cdot\mid z)$ for
every $z$. Then the marginal vector field
\begin{equation}
\label{eq:marginal-vf-scratch}
u_t(x) := \int u_t(x\mid z)\,\frac{p_t(x\mid z)\,p_{\mathrm{data}}(z)}{p_t(x)}\,dz
= \EE\bigl[u_t(x\mid Z)\mid X_t=x\bigr]
\end{equation}
(a $p_t(x)$-weighted average of the conditional fields, i.e.\ the
posterior expectation of $u_t(x\mid Z)$ given that the marginal path is at
$x$ at time $t$, treating $Z\sim p_{\mathrm{data}}$ as a latent variable and
$X_t\sim p_t(\cdot\mid Z)$) generates the marginal path $(p_t)$.
\end{theorem}

\begin{proof}
It suffices, by \Cref{thm:continuity}, to check that $(p_t,u_t)$ satisfies
the continuity equation \eqref{eq:continuity}. Each conditional pair
satisfies it: $\partial_tp_t(x\mid z) + \nabla\cdot(u_t(x\mid z)p_t(x\mid z)) = 0$
for every $z$. Multiply by $p_{\mathrm{data}}(z)$ and integrate over $z$:
\[
\partial_t\underbrace{\int p_t(x\mid z)p_{\mathrm{data}}(z)\,dz}_{=p_t(x)}
+
\nabla\cdot\underbrace{\int u_t(x\mid z)p_t(x\mid z)\,p_{\mathrm{data}}(z)\,dz}_{=:V(x)}
=0
\]
(differentiation and the divergence operator both commute with integration
over $z$, a standard application of dominated convergence, which we do not
belabour). It remains to check $V(x)=u_t(x)p_t(x)$ with $u_t$ as defined in
\eqref{eq:marginal-vf-scratch}: indeed,
$V(x) = \int u_t(x\mid z)\,p_t(x\mid z)p_{\mathrm{data}}(z)\,dz = p_t(x)\int u_t(x\mid z)\frac{p_t(x\mid z)p_{\mathrm{data}}(z)}{p_t(x)}\,dz = p_t(x)u_t(x)$,
using the definition of $u_t(x)$ verbatim. Hence
$\partial_tp_t(x)+\nabla\cdot(u_t(x)p_t(x))=0$, i.e.\ $(p_t,u_t)$ satisfies
\eqref{eq:continuity}, and \Cref{thm:continuity}'s converse direction gives
the claim.
\end{proof}

The marginal field \eqref{eq:marginal-vf-scratch}, like the marginal score
of \Cref{sec:score-matching-intro}, is an intractable integral (it
requires the unknown $p_t(x)$ in its own denominator). The following is
the flow-matching analogue of \Cref{thm:dsm-equivalence}, and, again, its
proof is structurally identical.

\begin{theorem}[Conditional flow matching]
\label{thm:cfm-equivalence-scratch}
Define
\[
\mathcal L_{\mathrm{FM}}(\theta) = \EE_{t\sim\mathcal U(0,1),\,x\sim p_t}\bigl[\|u_\theta(x,t)-u_t(x)\|^2\bigr],
\qquad
\mathcal L_{\mathrm{CFM}}(\theta) = \EE_{t,\,z\sim p_{\mathrm{data}},\,x\sim p_t(\cdot\mid z)}\bigl[\|u_\theta(x,t)-u_t(x\mid z)\|^2\bigr] .
\]
Then $\mathcal L_{\mathrm{FM}}(\theta) = \mathcal L_{\mathrm{CFM}}(\theta) + C$,
$C$ independent of $\theta$; the two losses have identical gradients and
minimisers.
\end{theorem}

\begin{proof}
Identical in structure to \Cref{thm:dsm-equivalence}'s proof, with
$\nabla_x\log p_\sigma(\cdot)$ replaced throughout by $u_t(\cdot)$, and the
Gaussian perturbation kernel replaced by the general conditional path
$p_t(\cdot\mid z)$: expand the square under $x\sim p_t$, the first
(quadratic-in-$\theta$) term already matches by
\Cref{prop:marginalisation} applied to the joint law of $(Z,X_t)$; the
third ($\theta$-independent) term is absorbed into $C$; and the cross term
is handled by the identity
$p_t(x)u_t(x) = \int u_t(x\mid z)p_t(x\mid z)p_{\mathrm{data}}(z)\,dz$
established in \Cref{thm:marginalisation-vf-scratch}'s proof (specifically
the definition of $V(x)$ there), which lets the cross term
$\int\langle u_\theta(x,t),u_t(x)\rangle p_t(x)\,dx$ be rewritten as an
expectation over $(Z,X_t)$ exactly as in \Cref{thm:dsm-equivalence}'s proof.
We do not repeat every line.
\end{proof}

\Cref{thm:cfm-equivalence-scratch} is the entire flow-matching recipe: pick
\emph{any} conditional path $p_t(\cdot\mid z)$ whose conditional vector
field $u_t(\cdot\mid z)$ can be written down in closed form, train
$u_\theta$ by regressing against that closed-form target on samples
$(t,z,x)$ drawn exactly as in $\mathcal L_{\mathrm{CFM}}$, and the trained
$u_\theta$ approximates the marginal field that transports
$p_{\mathrm{init}}$ to $p_{\mathrm{data}}$ --- without ever needing to
evaluate the intractable integral \eqref{eq:marginal-vf-scratch}.

\subsection{The linear (conditional optimal transport) path}
\label{sec:linear-path-scratch}

\begin{proposition}[The linear path and its vector field]
\label{prop:linear-path-scratch}
The path $p_t(\cdot\mid z) = \mathcal N\bigl((1-t)x_0+tz\ \big|_{x_0\sim p_{\mathrm{init}}}\bigr)$,
realised by $X_t = (1-t)X_0+tz$, $X_0\sim p_{\mathrm{init}}$ (for
$p_{\mathrm{init}}=\mathcal N(0,I)$, an ordinary linear interpolation
between an independently drawn noise sample and the data point $z$),
satisfies \Cref{def:cond-path-cfm}, and its conditional vector field is
\begin{equation}
\label{eq:linear-path-vf-scratch}
u_t(x\mid z) = \frac{z-x}{1-t} = z - x_0
\qquad\text{(the second equality valid exactly on the interpolant itself, }x=(1-t)x_0+tz\text{).}
\end{equation}
\end{proposition}

\begin{proof}
$p_0(\cdot\mid z)$: at $t=0$, $X_0\sim p_{\mathrm{init}}$ by construction.
$p_1(\cdot\mid z)$: at $t=1$, $X_1=z$ deterministically, i.e.\
$p_1(\cdot\mid z)=\delta_z$. For the vector field, differentiate the
interpolant directly: $X_t=(1-t)X_0+tz\ \Rightarrow\ \dot X_t = z-X_0$,
constant in $t$ along a fixed sample path; substituting
$X_0 = (X_t-tz)/(1-t)$ (inverting the interpolation formula) gives
$\dot X_t = z-(X_t-tz)/(1-t) = \bigl((1-t)z-X_t+tz\bigr)/(1-t) = (z-X_t)/(1-t)$,
recovering \eqref{eq:linear-path-vf-scratch}'s first form as a function of
the current point $X_t=x$ alone (as required by
\Cref{def:cond-path-cfm}, which needs $u_t(\cdot\mid z)$ as a genuine
function on $\RR^d$, not merely along one sample path).
\end{proof}

\Cref{eq:linear-path-vf-scratch}'s two forms are algebraically
equal (exactly, not merely approximately) when evaluated on the
interpolant, but computationally distinct: the first,
$(z-x)/(1-t)$, is a function of $(x,t,z)$ alone and is what a network,
given only the current state and time, is asked to match at training time;
the second, $z-x_0$, is \emph{constant in $t$} along a fixed pair
$(x_0,z)$ and is the more numerically convenient quantity to compute the
regression \emph{target} from during training (no division, no
singularity as $t\to1$), since $x_0,z$ are both known exactly at the moment
a training example is constructed. This exact distinction --- ``current-point
form'', singular at $t\to1$, versus ``constant-in-$t$ endpoint form'',
used for the actual training target --- reappears, verified directly
against the code, as precisely the two forms
\pyfunc{conditional\_probability\_path} and
\pyfunc{calc\_trans\_vector\_field} realise in SigmaFlow's
\file{r3\_flow\_matcher.py}, reconstructed exactly in
\Cref{sec:sigmaflow-translation}.

\begin{remark}[On the name ``conditional optimal transport'']
\label{rem:cot-name}
The linear path is frequently called the ``conditional optimal transport''
(CondOT) path, since for fixed $x_0,z$ the straight-line interpolant
$(1-t)x_0+tz$ is exactly the (constant-speed) geodesic of the flat
Euclidean metric between them, and constant-speed geodesics are the
building block of Euclidean optimal-transport (Wasserstein) interpolation.
We do not develop optimal transport theory in this document; the name is
recorded for consistency with the literature this monograph audits and
with the terminology used when this path is generalised to a Riemannian
manifold in \Cref{chap:riemannian-fm}, where ``geodesic interpolant''
(a notion this document \emph{has} developed, \Cref{def:metric-scratch}) is the
more directly justified term and the one used throughout the rest of this
document.
\end{remark}

\section{Sampling: integrating the learned ODE}
\label{sec:cfm-sampling-scratch}

Given a trained $u_\theta\approx u_t$, generation solves
$\dot X_t=u_\theta(X_t,t)$, $X_0\sim p_{\mathrm{init}}$, forward to $t=1$,
by any numerical ODE solver; the simplest, and the one this document's
code correspondences confirm is what both SigmaDock's (nominally deterministic
default, \Cref{sec:sigmadock-sampling}) and SigmaFlow's samplers actually
execute despite a \texttt{"heun"} option appearing in each function's type
signature, is the forward Euler step,
\[
X_{t+\Delta t} = X_t + \Delta t\,u_\theta(X_t,t) .
\]
Unlike the diffusion sampler of \Cref{eq:reverse-so3-scratch}, there is no
injected noise anywhere in this recipe: the entire stochasticity of the
generative model is confined to the single draw $X_0\sim p_{\mathrm{init}}$,
after which the trajectory is completely deterministic given the trained
network --- the structural fact, not an implementation simplification, that
makes the ``deterministic vs.\ stochastic'' framing of
\Cref{sec:sigmadock-sampling}'s noise-annealing discussion an
\emph{architecturally forced} property for flow matching specifically,
where for SigmaDock (as that section showed) it was, in the shipped
default configuration, already true as well, just not forced by the
method's structure.

\clearpage

% ================================================================
\chapter{Riemannian Flow Matching}
\label{chap:riemannian-fm}
% ================================================================

\Cref{chap:flow-matching} developed flow matching entirely in $\RR^d$; by
\Cref{prop:se3-product} this already suffices for translation. This
chapter lifts every construction to a Riemannian manifold $\mathcal M$,
following exactly the same two-step pattern
\Cref{chap:diffusion-so3} used for diffusion --- state the general
manifold construction, then specialise concretely to $SO(3)$ --- and, as
that chapter did, ends by comparing directly against SigmaFlow's actual
source code.

\section{Vector fields and flows on a manifold}
\label{sec:manifold-vf}

\begin{definition}[Vector field on a manifold, flow]
\label{def:manifold-vf}
A (time-dependent) \emph{vector field} on $\mathcal M$ is
$u:\mathcal M\times[0,1]\to T\mathcal M$ with $u_t(x)\in T_x\mathcal M$ for
every $x$ (i.e.\ a section of the tangent bundle at each time). Its
\emph{flow} $\phi_t:\mathcal M\to\mathcal M$ solves
$\frac{d}{dt}\phi_t(x)=u_t(\phi_t(x))$, $\phi_0(x)=x$, exactly as in
\Cref{def:vector-field-flow}, now understood as a curve \emph{on the
manifold} whose velocity at each instant is the prescribed tangent vector.
\end{definition}

Everything in \Cref{sec:cnf} generalises verbatim once ``vector field'' and
``flow'' are understood in the sense of \Cref{def:manifold-vf}: a CNF on
$\mathcal M$ is $(\phi_t)_\#p_{\mathrm{init}}$ for $p_{\mathrm{init}}$ a
distribution on $\mathcal M$; a manifold continuity equation
$\partial_tp_t+\mathrm{div}_g(u_tp_t)=0$ holds, with $\mathrm{div}_g$ the
Riemannian divergence (the manifold generalisation of $\nabla\cdot$,
defined so that \Cref{thm:continuity}'s proof --- which used only the
change-of-variables formula, itself valid on a manifold with the
Riemannian volume form in place of Lebesgue measure, and Liouville's
formula, which likewise generalises --- carries over without conceptual
change); and \Cref{thm:marginalisation-vf-scratch,thm:cfm-equivalence-scratch}'s
proofs use nothing beyond linearity of integration and the (now manifold)
continuity equation, so they too generalise verbatim, with $\|\cdot\|$
replaced by the Riemannian norm $\|\cdot\|_g$ of \Cref{def:metric-scratch}.
We record the two generalised results once, for reference, without
repeating their (unchanged) proofs.

\begin{theorem}[Riemannian conditional flow matching]
\label{thm:rcfm-general}
Let $p_t(\cdot\mid z)$ be a conditional probability path on $\mathcal M$
($p_0(\cdot\mid z)=p_{\mathrm{init}}$, $p_1(\cdot\mid z)=\delta_z$) with
conditional vector field $u_t(\cdot\mid z)$. Then
\[
\mathcal L_{\mathrm{RCFM}}(\theta)
=
\EE_{t,\,z\sim p_{\mathrm{data}},\,x\sim p_t(\cdot\mid z)}
\Bigl[\bigl\|u_\theta(x,t)-u_t(x\mid z)\bigr\|_g^2\Bigr]
\]
has the same gradients and minimisers, in $\theta$, as the (intractable)
marginal loss $\EE_{t,x\sim p_t}[\|u_\theta(x,t)-u_t(x)\|_g^2]$, where
$u_t(x)=\EE[u_t(x\mid Z)\mid X_t=x]$ as in
\eqref{eq:marginal-vf-scratch}.
\end{theorem}

What remains, exactly as in \Cref{chap:diffusion-so3}'s treatment of
diffusion, is to exhibit a \emph{specific}, computable conditional path and
its vector field on $SO(3)$ --- and this is where the Riemannian
generalisation of $\exp$/$\log$ (\Cref{def:explog-scratch}) does all the
work, in a construction dramatically simpler than
\Cref{chap:diffusion-so3}'s heat-kernel machinery.

\section{The geodesic conditional path}
\label{sec:geodesic-path-general}

\begin{definition}[Geodesic interpolant]
\label{def:geodesic-interpolant}
For $x_0,z\in\mathcal M$ (assuming, for simplicity, that $z$ lies within
the domain where $\log_{x_0}$ is defined --- always true when
$\mathcal M=SO(3)$ except at the single problematic value
$\theta=\pi$ discussed already in \Cref{sec:so3-log-scratch}, and handled
there), the \emph{geodesic interpolant} between $x_0$ and $z$ is
\[
\gamma_t(x_0,z) := \exp_{x_0}\bigl(t\,\log_{x_0}(z)\bigr) .
\]
\end{definition}

By \Cref{def:metric-scratch},\Cref{def:explog-scratch}, this traces exactly the
(constant-speed) geodesic from $x_0$ ($t=0$) to $z$ ($t=1$), the direct
manifold generalisation of $\RR^d$'s straight line $(1-t)x_0+tz$ (indeed,
on $\mathcal M=\RR^d$ with the flat metric, $\exp_x(v)=x+v$,
$\log_x(y)=y-x$ exactly, so \Cref{def:geodesic-interpolant} reduces to
$\gamma_t=x_0+t(z-x_0)=(1-t)x_0+tz$, recovering
\Cref{prop:linear-path-scratch}'s linear path precisely as the flat special
case).

\begin{proposition}[Conditional vector field of the geodesic interpolant]
\label{prop:geodesic-vf-general}
Fixing $x_0\sim p_{\mathrm{init}}$ and setting $X_t=\gamma_t(x_0,z)$
(\Cref{def:geodesic-interpolant}), the tangent vector
$u_t(X_t\mid z) := \dot X_t$ satisfies
\[
u_t(X_t\mid z) = \frac{\log_{X_t}(z)}{1-t} ,
\]
generalising \Cref{prop:linear-path-scratch}'s Euclidean formula exactly by
replacing $z-x$ with $\log_x(z)$.
\end{proposition}

We give this proposition's proof concretely for $\mathcal M=SO(3)$ in
\Cref{sec:fm-so3-scratch} below (the general Riemannian proof requires
facts about geodesics under reparameterisation and Jacobi fields that are
more naturally handled case-by-case for the one manifold this document
actually needs than developed in full generality; the $SO(3)$ computation
that follows is a complete, self-contained proof for the case that
matters).

\section{Flow matching on $SO(3)$}
\label{sec:fm-so3-scratch}

\begin{definition}[Geodesic interpolant on $SO(3)$, restated concretely]
\label{def:so3-interp-detail}
For $R_0,R_1\in SO(3)$, $t\in[0,1]$,
\begin{equation}
\label{eq:so3-interp-detail}
R_t := \exp_{R_0}\bigl(t\log_{R_0}(R_1)\bigr) = R_0\exp\bigl(t\log(R_0^\top R_1)\bigr) ,
\end{equation}
using the left-trivialised exponential and logarithm of \Cref{thm:rodrigues-scratch} (equivalently, per
\Cref{sec:so3-trivialisation}'s notation, this expresses the interpolant
via the group-log $\log(R_0^\top R_1)\in\so(3)$, a coordinate-free object
not itself labelled left or right, whose trivialisation only becomes a
meaningful distinction once we ask for the interpolant's \emph{velocity},
next).
\end{definition}

In quaternion language, \eqref{eq:so3-interp-detail} is exactly spherical
linear interpolation (SLERP) between $R_0$ and $R_1$, and it is the direct
$SO(3)$-analogue of the linear path
\Cref{prop:linear-path-scratch}, playing the role
\Cref{chap:diffusion-so3}'s $\mathrm{IGSO}(3)$ kernel played for diffusion
--- except that, unlike $\mathrm{IGSO}(3)$'s infinite spectral series,
\eqref{eq:so3-interp-detail} is a two-line, exactly-closed-form expression.

\begin{proposition}[Constant right-trivialised velocity of the geodesic]
\label{prop:so3-constant-velocity}
Let $\Xi:=\log(R_0^\top R_1)\in\so(3)$ and $R_t=R_0\exp(t\Xi)$ as in
\eqref{eq:so3-interp-detail}. Then $\dot R_t = R_t\,\Xi$, i.e.\ the curve
has \textbf{constant right-trivialised (body-frame)} angular velocity
$\omega^{\mathrm R}(t)\equiv\Xi^\vee$ (\Cref{def:trivialisation-scratch}),
for every $t$.
\end{proposition}

This is exactly \Cref{prop:geodesic-right-triv} from \Cref{chap:groups},
restated here in the flow-matching notation; we do not reprove it (the
proof --- differentiate $R_0\exp(t\Xi)$ directly, using that $\Xi$
commutes with its own exponential --- is identical).

\begin{proposition}[The current-point form: $SO(3)$ specialisation of \Cref{prop:geodesic-vf-general}]
\label{prop:so3-current-point}
$\log(R_t^\top R_1) = \exp\bigl((1-t)\Xi\bigr)$, hence
\begin{equation}
\label{eq:so3-cvf-scratch}
\frac{\log(R_t^\top R_1)^\vee}{1-t} = \Xi^\vee ,
\end{equation}
i.e.\ the ``current-point'' form of the conditional vector field
(\Cref{prop:geodesic-vf-general}'s $\log_{X_t}(z)/(1-t)$, here with
$X_t=R_t$, $z=R_1$) and the ``constant-in-$t$, endpoint'' form
(\Cref{prop:so3-constant-velocity}'s $\Xi^\vee$) are, exactly, the
\emph{same} tangent vector at every $t$ where both are defined --- not two
competing formulas, but one quantity computed two different ways.
\end{proposition}

\begin{proof}
$R_t^\top R_1 = \exp(-t\Xi)R_0^\top R_1 = \exp(-t\Xi)\exp(\Xi) = \exp\bigl((1-t)\Xi\bigr)$,
using $R_0^\top R_1=\exp(\Xi)$ (definition of $\Xi$) and that $\exp(-t\Xi)$
and $\exp(\Xi)$ commute (both are power series in the same matrix $\Xi$, so
$\exp(-t\Xi)\exp(\Xi)=\exp((1-t)\Xi)$ by the standard exponential addition
law for commuting matrices, itself a direct consequence of the binomial
theorem applied term-by-term to the two power series --- valid here
specifically because $-t\Xi$ and $\Xi$ commute, not in general for two
arbitrary elements of $\so(3)$). Taking $\log$ of both sides (valid for
$t\in[0,1)$, away from the $\theta=\pi$ subtlety of
\Cref{sec:so3-log-scratch}) gives $\log(R_t^\top R_1)=(1-t)\Xi$, and dividing
by $1-t$ gives \eqref{eq:so3-cvf-scratch}.
\end{proof}

This is the exact algebraic identity underlying the code's use of
\emph{two} superficially different formulas for what is provably the same
mathematical object; \Cref{sec:sigmaflow-rotation} records precisely where
each of the two forms of \Cref{prop:so3-current-point} is used
computationally, and why the distinction matters numerically (one form is
singular as $t\to1$, the other is not) despite the two being exactly equal
wherever both are defined.

\begin{theorem}[Riemannian CFM loss on $SO(3)$]
\label{thm:rcfm-so3}
\[
\mathcal L_{\mathrm{RCFM}}(\theta)
=
\EE_{t\sim\mathcal U(0,1),\,R_0\sim p_{\mathrm{init}},\,R_1\sim p_{\mathrm{data}}}
\Bigl[\bigl\|u_\theta(R_t,t) - \Xi^\vee\bigr\|_2^2\Bigr],
\qquad
R_t = R_0\exp(t\Xi),\ \Xi=\log(R_0^\top R_1) ,
\]
has the same minimiser, $u_\theta(R,t)\to u_t(R)$ (the true marginal
field), as the intractable marginal loss, by
\Cref{thm:rcfm-general} specialised to $\mathcal M=SO(3)$ (using
\Cref{prop:so3-distance}'s bi-invariant metric, under which $\|\cdot\|_g$
on $T_RSO(3)\cong\RR^3$ is the ordinary Euclidean norm, per
\Cref{prop:so3-distance}'s proof).
\end{theorem}

\begin{remark}[Prior distribution]
\label{rem:so3-fm-prior}
The natural (and, per \Cref{chap:sigmaflow}'s code correspondence, actually
used) source distribution for $SO(3)$ flow matching is the exact
Haar-uniform measure $\mathcal U(SO(3))$ of \Cref{prop:haar-so3}, sampled
directly via \Cref{prop:haar-so3}'s code correspondence
(\pyfunc{sample\_uniform}) --- \emph{not} the diffusion side's
$t=1$-of-a-VE-schedule approximation to it
(\Cref{sec:diff-so3-numerics}'s \pyfunc{sample\_ref}). This is possible
\emph{because} flow matching never needs a noising process to define its
prior at all: any fixed, sampleable distribution with the right invariance
properties (\Cref{prop:equivariant-invariant-scratch}) will do, and the
exact Haar measure is simply the natural, no-approximation-needed choice,
unlike diffusion's necessarily-approximate large-noise limit.
\end{remark}

Sampling integrates
\begin{equation}
\label{eq:so3-flow-euler-scratch}
R_{t+\Delta t} = \exp_{R_t}\bigl(\Delta t\,u_\theta(R_t,t)\bigr) = R_t\exp\bigl(\Delta t\,\widehat{u_\theta(R_t,t)}\bigr) ,
\end{equation}
exactly \Cref{prop:euler-right-triv}'s right-trivialised Euler update,
which --- by construction, since $\exp$ always lands back on $SO(3)$
regardless of step size (\Cref{sec:rodrigues-scratch}) --- keeps every
iterate exactly on $SO(3)$, for any number of steps, not merely
approximately for small steps.

Comparing \Cref{sec:diff-so3,sec:fm-so3-scratch}: flow matching on $SO(3)$
needs no Brownian motion, no $\mathrm{IGSO}(3)$ heat kernel or its
truncated spectral series, no numerical tabulation/caching
(\Cref{sec:diff-so3-numerics}), no score estimation, and no reverse-time
SDE --- only the closed-form geodesic interpolant
\eqref{eq:so3-interp-detail}, its (equally closed-form) velocity
\Cref{eq:so3-cvf-scratch}, and a deterministic ODE. This elimination of an
entire, nontrivial piece of numerical machinery is the first of the two
arguments motivating the SigmaDock-to-SigmaFlow substitution examined in
this project's broader work; the second, a reduction in the number of
network evaluations needed at sampling time, is an empirical question
addressed in that broader experimental work, not re-derived here.

\begin{remark}[Numerical instability near $t=1$]
\label{rem:so3-singularity-remark}
The current-point form \eqref{eq:so3-cvf-scratch} carries a $(1-t)^{-1}$
factor that is exactly cancelled, algebraically, by
\Cref{prop:so3-current-point}'s identity --- but only when evaluated
\emph{exactly}, in exact arithmetic. In floating-point arithmetic, the
numerator $\log(R_t^\top R_1)$ and the denominator $1-t$ both individually
approach $0$ as $t\to1$, so this specific closed form is numerically
fragile in exactly this limit, even though the underlying tangent vector
it computes is perfectly well-behaved (constant, in fact, by
\Cref{prop:so3-constant-velocity}). This is a genuine, verified numerical
consideration, examined empirically (via oracle-vector-field ablations
that isolate this exact mechanism from network-quality effects) in this
project's broader experimental investigation, referenced but not
re-conducted in this theory document.
\end{remark}

\section{Flow matching on the pose manifold $SE(3)^N$}
\label{sec:fm-se3n}

By \Cref{prop:se3-product}'s Riemannian-product structure, applied to flow
matching exactly as it was applied to diffusion in
\Cref{sec:sigmadock-forward}: the translational component uses
\Cref{prop:linear-path-scratch}'s linear path independently on each
fragment's centre of mass, the rotational component uses
\Cref{sec:fm-so3-scratch}'s geodesic interpolant independently on each
fragment's orientation, and no cross term couples them --- the geometry
factorises completely, with all inter-fragment and rotation--translation
coupling introduced \emph{only} by the shared neural network
(\Cref{chap:equiformer}) that predicts both components jointly from the
same input graph, exactly as \Cref{sec:sigmadock-state-space} established
for SigmaDock and as \Cref{chap:sigmaflow} now confirms, component by
component and directly against the code, holds equally for SigmaFlow.

\clearpage

% ================================================================
\part{SigmaFlow}
\label{part:sigmaflow}
% ================================================================

\chapter{SigmaFlow: Full Reconstruction}
\label{chap:sigmaflow}

We now perform, for SigmaFlow, exactly what \Cref{chap:sigmadock} did for
SigmaDock: assemble the already-derived theory
(\Cref{chap:flow-matching,chap:riemannian-fm}) with the specific network of
\Cref{chap:equiformer} into a complete, code-verified reconstruction.
Because \Cref{chap:sigmadock} already derived the state space
(\Cref{sec:sigmadock-state-space}), the fragmentation and triangulation
machinery (\Cref{sec:sigmadock-fr3d,sec:triangulation-precise}), the
symmetry argument (\Cref{sec:sigmadock-symmetry-precise}), the input graph
and backbone (\Cref{sec:sigmadock-graph}), and the fragment-reconstruction
map $\varphi$ (\Cref{sec:fragment-reconstruction}) --- all of which
\Cref{tab:sigmadock-vs-sigmaflow-precise} below confirms are unchanged ---
this chapter focuses entirely on what \emph{does} change: the probability
path, the training target, the loss, and the sampler.

\section{What changes, and what does not}
\label{sec:sigmaflow-overview}

\begin{longtable}{p{3.6cm}p{5.3cm}p{5.3cm}}
\caption{Component-wise comparison, verified against the source in both
trees. Rows above the double rule are unchanged; rows below are what this
chapter derives.}
\label{tab:sigmadock-vs-sigmaflow-precise}\\
\toprule
Component & SigmaDock (\Cref{chap:sigmadock}) & SigmaFlow \\
\midrule
\endfirsthead
\toprule
Component & SigmaDock & SigmaFlow \\
\midrule
\endhead
State space $\mathcal M_m=SE(3)^m$, product metric & \multicolumn{2}{c}{identical (\Cref{sec:sigmadock-state-space,prop:se3-product})} \\
Fragmentation (FR3D), triangulation & \multicolumn{2}{c}{identical (\Cref{sec:sigmadock-fr3d,sec:triangulation-precise}; code diff'd byte-identical, per this project's own audit)} \\
Input graph, featurisation, backbone & \multicolumn{2}{c}{identical (\Cref{chap:equiformer,sec:sigmadock-graph})} \\
Fragment reconstruction $\varphi$ & \multicolumn{2}{c}{identical formula (\Cref{prop:fragment-reconstruction})} \\
\midrule
Prior $p_{\mathrm{init}}$, translation & $\mathcal N(0,I)$ (VP schedule's $s\to1$ limit) & $\mathcal N(0,I)$ (exact, by construction) \\
Prior $p_{\mathrm{init}}$, rotation & $\approx\mathcal U(SO(3))$ (VE schedule's $s\to1$ limit, \Cref{sec:diff-so3-numerics}) & $\mathcal U(SO(3))$ (exact Haar sampler, \Cref{rem:so3-fm-prior}) \\
Probability path & VP Gaussian $\oplus$ $\mathrm{IGSO}(3)$ heat kernel & linear $\oplus$ geodesic (SLERP) interpolant \\
Regression target & conditional score $\nabla\log p_{t\mid1}$ & conditional velocity $u_t(\cdot\mid z)$ \\
Numerical machinery & truncated $\mathrm{IGSO}(3)$ series, tabulation, caching (\Cref{sec:diff-so3-numerics}) & none --- closed form throughout \\
Output dressing & \Cref{sec:sd-head-branches}'s $t$-dependent scalar prefactors & \textbf{none} (\Cref{sec:sigmaflow-nodressing}) \\
Loss terms & four (\texttt{T\_score,R\_score,T0,R0}, \Cref{sec:sigmadock-loss}) & \textbf{two} (\texttt{loss\_trans,loss\_R}, \Cref{sec:sigmaflow-loss}) \\
Sampler & reverse-time SDE (deterministic by default configuration, \Cref{sec:sigmadock-sampling}) & deterministic ODE (architecturally, not by configuration) \\
Time convention (internal) & $t=0$ data, $t=1$ noise & $t=0$ noise, $t=1$ data (\textbf{reversed}, \Cref{sec:sigmaflow-time-convention}) \\
\bottomrule
\end{longtable}

\section{A time-convention warning, confirmed at the source level}
\label{sec:sigmaflow-time-convention}

\Cref{rem:time-convention-global} fixed a global convention for this
document ($t=0$ noise, $t=1$ data) and warned that SigmaDock's own native
convention runs the opposite way internally
(\Cref{rem:time-convention-scratch} in \Cref{chap:sigmadock}, and
\Cref{sec:sigmadock-sampling}'s observation that the diffusion sampler's
own timestep array decreases from $t_{\max}=1$ to $t_{\min}$). We record
here, as a fact independently confirmed against \emph{both} source trees
during the preparation of this document (not merely inferred from
exposition), that this is not a convention chosen for pedagogical
convenience in this monograph but a genuine, load-bearing difference
between the two codebases' own internal time variables:

\begin{itemize}
\item \pyfunc{R3Diffuser}/\pyfunc{SO3Diffuser} (SigmaDock): $t=0$ is clean
data ($\alpha_{t=0}=1$, no noise); $t=1$ is the fully-noised end
(\pyfunc{sample\_ref} literally evaluates the forward kernel at $t=1$).
\item \pyfunc{R3\_FlowMatcher}/\pyfunc{SO3\_FlowMatcher} (SigmaFlow): $t=0$
is the noise/source end ($X_0\sim p_{\mathrm{init}}$ drawn internally by
\pyfunc{conditional\_probability\_path}); $t=1$ is the data end
($X_1=z$, the training target).
\end{itemize}

This document's global convention (\Cref{rem:time-convention-global})
matches SigmaFlow's own native convention exactly, and every formula
attributed to SigmaDock throughout \Cref{chap:sigmadock} was, correspondingly,
already relabelled via $s=1-t$ before being stated (\Cref{rem:time-convention-scratch}
there). Any reader cross-referencing a formula directly against SigmaDock's
own source should apply this relabelling explicitly; failing to do so is,
per this document's own code-research process, a genuine and easy source of
sign errors, exactly as this document's introduction anticipated.

\section{Translation flow matching, exactly as implemented}
\label{sec:sigmaflow-translation}

\begin{codebox}
\textbf{Mathematical object:} the source distribution, conditional path,
and both forms of the conditional vector field of
\Cref{prop:linear-path-scratch} \\
\textbf{Implementation:} \file{r3\_flow\_matcher.py}, class
\pyfunc{R3\_FlowMatcher} (56 lines total) \\
\textbf{Interpretation:}
\pyfunc{sample\_init}$(n)=\texttt{randn(n,3)}$ ($X_0\sim\mathcal N(0,I_3)$
per fragment, no rescaling by any $\sigma_{\min}$-type parameter);
\pyfunc{conditional\_probability\_path}$(x_1,t)$ draws $x_0$
\emph{internally} (unlike the diffusion side's
\pyfunc{forward\_marginal}, which takes $x_0$ as an argument --- a
structural difference reflecting that flow matching's source is drawn once
per training example and never needs to be revisited, whereas diffusion's
$x_0$ is the fixed data point being noised), computes
$x_t=(1-t)x_0+tx_1$ (\Cref{prop:linear-path-scratch}'s interpolant) and
returns $u_t=x_1-x_0$, the \emph{constant-in-$t$, endpoint} form --- this
is the quantity actually regressed against during training
(\Cref{sec:sigmaflow-loss}). \pyfunc{calc\_trans\_vector\_field}$(x_t,x_1,t)=(x_1-x_t)/(1-t)$
implements the \emph{current-point} form
\eqref{eq:linear-path-vf-scratch}, singular as $t\to1$ by construction
(the docstring itself notes $t$ must lie in $(0,1)$); this form is
\textbf{not} used during training (confirmed by tracing every call site in
\file{sigma\_flow\_generator.py}), only at sampling time
(\Cref{sec:sigmaflow-sampling}) and for the \texttt{use\_true\_vector\_field}
oracle debugging path, where $x_0$ is not naturally available but the
current point $(x_t,t)$ is. \pyfunc{euler\_step}$(x_t,v_t,dt)=x_t+v_t\,dt$,
the plain (non-geodesic, since $\RR^3$ is flat) Euler step of
\Cref{sec:cfm-sampling-scratch}, and the module's own inline comment notes
this step is \emph{exact} (zero discretisation error, for any step size)
specifically because $u_t$ is constant along the true conditional path ---
a special property of the linear path, not shared by the rotational
component's geodesic (\Cref{sec:sigmaflow-rotation}) once the network's
prediction, rather than the true field, is substituted in. The constructor
parameter \texttt{sigma\_min} is accepted but \textbf{provably unused}: it
is stored (\texttt{self.sigma\_min = sigma\_min}) and never read anywhere
else in the file, with an explicit docstring noting it is ``reserved for
future noised OT path, currently unused --- the linear path below is
deterministic''.
\end{codebox}

\section{Rotation flow matching, exactly as implemented}
\label{sec:sigmaflow-rotation}

\begin{codebox}
\textbf{Mathematical object:} the source distribution, conditional path,
and both forms of the conditional vector field of
\Cref{sec:fm-so3-scratch} \\
\textbf{Implementation:} \file{so3\_flow\_matcher.py}, class
\pyfunc{SO3\_FlowMatcher} (63 lines total) \\
\textbf{Interpretation:}
\pyfunc{sample\_init}$(n)=$\pyfunc{so3\_utils.sample\_uniform}$(n)$, the
exact Haar sampler of \Cref{prop:haar-so3}'s code correspondence
(\Cref{rem:so3-fm-prior}'s claim, confirmed directly). The class's own
module-level comment states the convention adopted is
``right-trivialisation''; per \Cref{sec:so3-trivialisation}'s resolved
terminology, this matches --- not merely by assumption but by direct
inspection of \pyfunc{euler\_step}$(R_t,v_t,dt)=R_t\,\exp(v_t\,dt)$, which
right-multiplies the exponential onto $R_t$, exactly
\Cref{prop:euler-right-triv}'s right-trivialised update formula.
\pyfunc{conditional\_probability\_path}$(R_1,t)$ draws $R_0$ internally,
computes $\Delta=R_0^\top R_1$, $\Xi=\log(\Delta)\in\so(3)$
(\Cref{def:so3-interp-detail}'s notation), sets $u_t=\Xi$ (a $[n,3,3]$
skew-symmetric-matrix-valued, i.e.\ $\so(3)$-valued, not yet
vee-mapped-to-$\RR^3$ quantity --- the \emph{constant-in-$t$} form of
\Cref{prop:so3-constant-velocity}, and the training target,
\Cref{sec:sigmaflow-loss}), and $R_t=R_0\exp(t\Xi)$
(\Cref{eq:so3-interp-detail}, confirmed exactly).
\pyfunc{calc\_rot\_vector\_field}$(R_t,R_1,t)=\log(R_t^\top R_1)/(1-t)$
implements the \emph{current-point} form
\eqref{eq:so3-cvf-scratch}, algebraically equal to $\Xi$ by
\Cref{prop:so3-current-point} but singular in floating-point arithmetic as
$t\to1$ (\Cref{rem:so3-singularity-remark}) --- again not used during
training, only at sampling time and inside the
\texttt{use\_true\_vector\_field} oracle path. The module additionally
carries an explicit inline comment recording exactly
\Cref{caution:omega-clamp}'s finding in this document's own words: the
original SigmaDock \pyfunc{Omega} clamp broke exact
$\exp\circ\log=\mathrm{id}$ round trips needed by
\eqref{eq:so3-interp-detail}, motivating the narrower clamp confirmed in
\Cref{sec:so3-log-scratch} --- i.e.\ the bugfix's rationale is documented,
first-party, directly in the file that needed it.
\end{codebox}

\begin{remark}[The hard-coded $R_1\equiv I$ simplification]
\label{rem:r1-identity-simplification}
A significant, easily-missed simplification, confirmed directly against
\file{sigma\_flow\_generator.py}::\pyfunc{get\_fragment\_com\_and\_rot}: the
function that computes each fragment's data-time state
$(\mathrm{trans}_1,R_1)$ from the reference conformer sets the rotational
component to the identity, \emph{always}, for every fragment, with an
explicit inline comment: ``Rotation matrix is identity: it is (and should
be) independent in the vector field.'' Consequently, throughout the actual
training and sampling code, $R_1\equiv I$ is not a general fragment
orientation but a fixed constant, and formulas that look like they involve
a genuine $R_1$ --- for instance the delta-rotation
$\Delta R := R_t\,R_1^\top$ appearing inside
\Cref{prop:fragment-reconstruction}'s reconstruction map when specialised
to this codebase --- collapse to $\Delta R=R_t$ exactly. This is
consistent with, and is the concrete mechanism realising,
\Cref{sec:sigmadock-symmetry-precise}'s gauge-invariance discussion: since
the reference conformer's own orientation is arbitrary
(\Cref{thm:sigmadock-symmetry}'s gauge freedom), fixing it at the identity
is simply one particular, legitimate choice of gauge, not a loss of
generality --- but a reader auditing the code without this fact stated
explicitly could easily mistake $R_1$ for a nontrivial, data-derived
quantity, which it is not.
\end{remark}

\section{No output dressing: the entire diffusion-specific scalar machinery is gone}
\label{sec:sigmaflow-nodressing}

\begin{codebox}
\textbf{Mathematical object:} the network's predicted vector field,
$(\mathrm{pred\_u\_t\_trans},\mathrm{pred\_u\_t\_R})$ \\
\textbf{Implementation:} \file{sigma\_flow\_generator.py}::\pyfunc{\_compute\_vector\_field},
quoted here in full (it is short enough to reproduce verbatim, and its
brevity is itself the point):
\begin{verbatim}
def _compute_vector_field(self, sampled, updates, t_batch):
    """
    Interprets the pooled per-fragment force/torque directly as the
    predicted conditional vector field. No time dependent scaling
    needed, unlike the diffusion score.
    """
    pred_u_t_trans = updates["total_force"]
    pred_u_t_R = updates["omega"]
    return {"pred_u_t_trans": pred_u_t_trans, "pred_u_t_R": pred_u_t_R}
\end{verbatim}
\textbf{Interpretation:} \texttt{updates["total\_force"]} and
\texttt{updates["omega"]} are exactly the raw Newton--Maruyama outputs
$dT=F_F/M_F$ and $\widehat{dW}=\widehat{I_F^{-1}\tau_F}$ of
\eqref{eq:newton-maruyama-precise} --- \emph{unchanged} from
\Cref{sec:newton-euler-precise}'s construction, since the physical
Newton--Euler reduction (net force, torque, mass, inertia) is one of the
components \Cref{tab:sigmadock-vs-sigmaflow-precise} lists as unmodified.
What is different is that this raw mechanical output is returned
\emph{directly} as the predicted vector field, with \textbf{no}
$t$-dependent rescaling of the kind
\Cref{eq:sd-head-precise,eq:sd-head-rot-space,eq:sd-head-rot-score}
apply: no \texttt{r3\_scaling}, no \texttt{so3\_scaling}, no
$\partial_\omega\log f_\sigma/\omega$ dressing, no branch selection
(\texttt{rot\_score\_method}/\texttt{rot\_score\_scaling}, whose
corresponding config fields were removed entirely from
\file{config.py} in an earlier development pass of this project, since
they had become dead code once this function stopped consulting them).
The comment's own justification --- ``no time dependent scaling needed,
unlike the diffusion score'' --- is a genuine mathematical observation, not
an oversight: \Cref{prop:so3-constant-velocity,prop:linear-path-scratch}
both establish that the flow-matching conditional targets are
\emph{constant in $t$} (unlike the diffusion score, whose scale, per
\Cref{sec:sde-perspective}'s discussion, varies enormously across noise
levels and specifically requires the $\lambda(s)$/scalar-prefactor
machinery to keep training well-conditioned) --- \textbf{but} this
observation concerns the \emph{training target's} scale, not necessarily
whether the specific quantity a fixed-capacity network can most easily
regress onto (raw Newton--Euler output, versus some other, still
$t$-independent, reparameterisation) is optimally chosen; this project's
broader experimental investigation (the ``time-weighting'' variant referenced
in this project's development history) tested, and found no benefit from,
reintroducing a $t$-dependent multiplicative reweighting on top of this
already-unweighted target, discussed further in
\Cref{sec:variant-time-weighting}.
\end{codebox}

\section{The loss function}
\label{sec:sigmaflow-loss}

\begin{codebox}
\textbf{Mathematical object:} $\mathcal L_{\mathrm{trans}}$,
$\mathcal L_R$ of \Cref{thm:cfm-equivalence-scratch,thm:rcfm-so3},
evaluated at their sampled-path targets \\
\textbf{Implementation:} \file{sigma\_flow\_generator.py}::\pyfunc{compute\_losses} \\
\textbf{Interpretation:}
\begin{align*}
\texttt{loss\_trans} &= \|\mathrm{pred\_u\_t\_trans} - u_t^{\mathrm{trans}}\|_2^2
&&\text{(sum over the 3 coordinates, per fragment; }u_t^{\mathrm{trans}}=x_1-x_0\text{ of \Cref{prop:linear-path-scratch})}\\
\texttt{loss\_R} &= \|\mathrm{pred\_u\_t\_R} - u_t^R\|_{\mathrm{Fro}}^2
&&\text{(Frobenius norm}^2\text{ of the }3\times3\text{ difference; }u_t^R=\Xi\text{ of \Cref{def:so3-interp-detail})}
\end{align*}
Both are raw per-fragment sums with \textbf{no weighting and no clamping}
inside \pyfunc{compute\_losses} itself --- confirming, precisely,
\Cref{tab:sigmadock-vs-sigmaflow-precise}'s claim that SigmaFlow's loss has
\emph{two} terms where SigmaDock's \pyfunc{compute\_losses}
(\Cref{sec:sigmadock-loss}) has \emph{four}: the two ``data-space''
($\hat x_0$-reconstruction) terms \texttt{T0}/\texttt{R0} that SigmaDock's
diffusion parameterisation computes alongside its score-matching terms have
\textbf{no analogue at all} in SigmaFlow's loss, confirmed by direct
tracing of \file{trainer.py}::\pyfunc{SigmaLightningModule.\_shared\_step},
whose total-loss combination
$\texttt{total\_loss} = \texttt{trans\_score\_weight}\cdot\texttt{loss\_trans} + \texttt{rot\_score\_weight}\cdot\texttt{loss\_R}$
consumes exactly these two keys and no others (the CLI flag names
\texttt{trans\_score\_weight}/\texttt{rot\_score\_weight} are themselves a
naming holdover from the diffusion codebase this file is shared with;
functionally, for the flow-matching path, they are ordinary fixed loss
weights on \eqref{eq:linear-path-vf-scratch}/\eqref{eq:so3-cvf-scratch}'s
regression losses, not score-matching weights in
\Cref{sec:sde-perspective}'s $\lambda(s)$ sense). Per-molecule reduction
uses \file{trainer.py}::\pyfunc{scaled\_fragmented\_loss}, structurally
identical to \Cref{sec:sigmadock-loss}'s code correspondence (shared
infrastructure, per \Cref{tab:sigmadock-vs-sigmaflow-precise}).
\end{codebox}

\section{The full training step}
\label{sec:sigmaflow-training-pipeline}

We now trace \pyfunc{SigmaFlowGenerator.forward} end to end, with every
intermediate tensor's shape as commented in the source, assembling every
piece derived above into the single computational pipeline a training
batch actually executes.

\begin{enumerate}
\item \textbf{Prepare batch, sample time.}
$t=\pyfunc{\_sample\_time}(\mathrm{batch})\in[\eps_t,1]^B$, one scalar
\emph{per molecule} in the batch (\emph{not} per fragment), via
$t=\eps_t+(1-\eps_t)\cdot\mathrm{Unif}(0,1)$, $\eps_t=\texttt{HPARAMS.general.epsilon\_t}=0.01$
--- matching \Cref{sec:sde-perspective}'s $t\sim\mathcal U(0,1)$ up to the
training-time floor $\eps_t$, whose purpose (per an explicit code comment
quoted verbatim, judged authoritative and reproduced here rather than
re-derived) is distributional consistency between training and sampling,
not any numerical-singularity avoidance: ``training's \texttt{sample\_time}
never samples $t$ below this; sampling's \texttt{t\_min} should match it
$\dots$ not related to the separate $1/(1-t)$ singularity near $t=1$'' ---
i.e.\ the network should never be \emph{queried} at a $t$ value it was
never trained on, which is a distributional/generalisation concern
entirely distinct from \Cref{rem:so3-singularity-remark}'s numerical
concern about the current-point vector-field formula near $t=1$.
\item \textbf{Get initial (data-time) states.}
$\mathrm{pos}_0,\ \mathrm{trans}_1,\ R_1,\ \mathrm{num\_fragments}
= \pyfunc{\_get\_initial\_states}(\mathrm{batch})$: despite the
subscript ``$0$'', $\mathrm{pos}_0$ is the (pocket-centred,
\texttt{dimensional\_scale}-rescaled, \Cref{rem:dimensional-scale}) ground
truth/reference atom geometry --- the fixed per-fragment \emph{template}
$x^{\mathrm{ref}}$ of \Cref{prop:fragment-reconstruction}, reused at every
$t$, not a draw from the source distribution --- while $\mathrm{trans}_1,R_1$
are the per-fragment data-time (centre-of-mass, and, per
\Cref{rem:r1-identity-simplification}, identically-$I$) rigid-motion state.
$t$ is broadcast from one value per molecule to one value per fragment,
$t_{\mathrm{batch}}\in[\eps_t,1]^{\Sigma F}$ (\emph{all fragments of a
given molecule share the same sampled $t$}, a design choice with no
counterpart requirement in the theory of
\Cref{chap:flow-matching,chap:riemannian-fm}, which permit --- but do not
require --- an independent $t$ per component; the code's own inline
comment flags this exact choice as a documented, unresolved limitation:
``better sampling. Uniform makes noises (sigmas) that are too
redundant'').
\item \textbf{Sample the conditional path.}
$(\mathrm{trans}_t,R_t,u_t^{\mathrm{trans}},u_t^R) = \pyfunc{\_sample\_flow}(\mathrm{trans}_1,R_1,t_{\mathrm{batch}})$,
i.e.\ \Cref{sec:sigmaflow-translation,sec:sigmaflow-rotation}'s
\pyfunc{conditional\_probability\_path}, dispatched componentwise per
\Cref{prop:se3-product}.
\item \textbf{Reconstruct atom positions at time $t$.}
$\mathrm{pos}_t = \varphi(\mathrm{trans}_t,R_t)$ via
\Cref{prop:fragment-reconstruction}'s formula
(\pyfunc{\_apply\_transformations}), and the graph's transient (interaction)
edges are recomputed at these new positions
(\pyfunc{\_update\_batch}, \Cref{sec:sigmadock-graph}'s discussion of the
global/transient topology split).
\item \textbf{Network forward pass.}
$\mathrm{lig\_pseudoforces},\mathrm{forces\_idxs} = \pyfunc{\_compute\_forces}(\mathrm{batch},t)$
--- exactly one EquiformerV2 evaluation
(\Cref{chap:equiformer,sec:newton-euler-precise}'s code correspondence),
identical in cost and mechanism to SigmaDock's own single-network-call
Newton--Euler pipeline.
\item \textbf{Newton--Euler reduction and Newton--Maruyama update.}
$(F_F,\tau_F,M_F,I_F) = \pyfunc{\_compute\_fragment\_dynamics}(\dots)$
(\eqref{eq:newton-euler-precise}); $(dT,\widehat{dW})=\pyfunc{\_predict\_fragment\_updates}(\dots)$
(\eqref{eq:newton-maruyama-precise}) --- structurally byte-for-byte the
same construction as \Cref{sec:newton-euler-precise}, confirming
\Cref{tab:sigmadock-vs-sigmaflow-precise}'s claim that this component is
unchanged.
\item \textbf{Predicted vector field.}
$(\mathrm{pred\_u\_t\_trans},\mathrm{pred\_u\_t\_R}) = \pyfunc{\_compute\_vector\_field}(\dots)$,
\Cref{sec:sigmaflow-nodressing}'s undressed pass-through.
\item \textbf{Loss.} \Cref{sec:sigmaflow-loss}'s two-term
\pyfunc{compute\_losses}, weighted and summed by
\pyfunc{trainer.py}::\pyfunc{SigmaLightningModule.\_shared\_step}, wrapped
in the try/except-\texttt{AssertionError} and finite-value guards this
project's own development history documents in detail (catching
\Cref{def:fragment-mass-precise}'s ``fragment mass $<2$'' assertion
specifically, among other numerical failure modes).
\end{enumerate}

Every step above, except step 8's loss formula and steps 2--3's use of
the flow-matching (rather than diffusion) conditional path, is verified,
in \Cref{tab:sigmadock-vs-sigmaflow-precise}, to be identical to
SigmaDock's corresponding pipeline stage --- the precise, exhaustive sense
in which this project's conversion of SigmaDock into SigmaFlow is, as
claimed throughout this document's audited source material, a
\emph{surgical}, minimally-invasive substitution of the probability path
and training target alone.

\section{Sampling}
\label{sec:sigmaflow-sampling}

\begin{codebox}
\textbf{Mathematical object:} the forward ODE integration loop of
\Cref{sec:cfm-sampling-scratch,sec:fm-so3-scratch} \\
\textbf{Implementation:} \file{sampling.py}::\pyfunc{sampler} (in
\file{SigmaFlow\_Development/src/sigmadock/diff/}) \\
\textbf{Interpretation:} default \texttt{num\_steps}$=25$,
\texttt{t\_min}$=\eps_t=0.01$ (matching training's floor, per this
chapter's Step 1 discussion above), \texttt{rho}$=3.0$, default
\texttt{discretization="power"} (\emph{not} SigmaDock's default
\texttt{"edm"} --- a genuine, code-confirmed default-configuration
asymmetry between the two samplers, flagged here rather than assumed
identical: any timing/quality comparison between the two must either
override one to match the other or explicitly account for this
difference), \texttt{solver="euler"} (Heun again unimplemented, per
\Cref{sec:sigmadock-sampling}'s finding, verbatim-identical dead
\texttt{\# TODO} comment block in this file too). The timestep array walks
\emph{upward}, $t_{\min}\to t_{\max}=1$, matching this document's global
convention directly (no relabelling needed, unlike SigmaDock's sampler).
Each loop iteration: (1) evaluate the closed-form current-point field
$u_t(\cdot\mid z)$ (\pyfunc{\_compute\_true\_vector\_field}, i.e.\
\Cref{eq:linear-path-vf-scratch,eq:so3-cvf-scratch}'s singular-but-exact
forms, using the \emph{known} $\mathrm{trans}_1,R_1$ --- available here
because sampling is being run, in this project's own evaluation
methodology, on complexes whose true pose is known, purely for
diagnostic loss logging at each step, not because a real blind-docking
deployment would have this available); (2) if
\texttt{use\_true\_vector\_field=False} (the normal, non-diagnostic case),
run the network (Steps 5--7 of \Cref{sec:sigmaflow-training-pipeline}'s
list) to get the predicted field instead; (3) Euler step
(\pyfunc{SE3\_FlowMatcher.euler\_step}, dispatching to
\Cref{sec:sigmaflow-translation,sec:sigmaflow-rotation}'s componentwise
\pyfunc{euler\_step} implementations); (4) reconstruct positions via
$\varphi$ (\Cref{prop:fragment-reconstruction}) at the new
$(\mathrm{trans}_{t+dt},R_{t+dt})$. Exactly one network evaluation per
step, matching SigmaDock's sampler exactly in this specific respect
(\Cref{sec:sigmadock-sampling}).
\end{codebox}

\begin{remark}[The \texttt{use\_true\_vector\_field} oracle: what Test~1 established]
\label{rem:oracle-test-recap}
Substituting the closed-form field for the network's prediction at
\emph{every} step (\texttt{use\_true\_vector\_field=True}) isolates a
question this chapter's derivations make precise: does the ODE
integration/reconstruction machinery itself (Euler stepping, the
$\varphi$ reconstruction, batching and edge recomputation) introduce error,
independent of network quality? By \Cref{prop:so3-current-point}'s exact
algebraic identity and \Cref{prop:linear-path-scratch}'s exact linearity
(both established rigorously in this chapter, not merely observed
empirically), the answer, \emph{in exact arithmetic}, must be no: the
oracle field reconstructs the true trajectory exactly, for any number of
steps, since both conditional vector fields are (by
\Cref{sec:cfm-sampling-scratch}'s remark on the linear path, and
\Cref{prop:so3-current-point} for the geodesic path) either exactly
integrable by a single Euler step (translation) or exactly consistent with
the constant-velocity geodesic the Euler step is discretising (rotation,
up to floating-point round-off). This is the theoretical content
underlying this project's empirical ``Test 1'' finding (oracle-vector-field
sampling reconstructs the true structure to within floating-point
precision, $\le0.0002$\,\r A deviation, across a sample of real generated
structures) --- the theorem, derived here, is what the experiment
confirmed, not a separate, independent fact.
\end{remark}

\section{Likelihood and confidence: what is implemented, and what is not}
\label{sec:likelihood-confidence}

\Cref{sec:cnf-likelihood} derived, from first principles, that a CNF's
log-density along any trajectory evolves by
$\frac{d}{dt}\log p_t(X_t) = -(\nabla\cdot u_t)(X_t)$
(\Cref{cor:instantaneous-cov}), estimable via the Hutchinson trace
estimator without a separate density network. This is a genuine,
well-established capability of the flow-matching framework in general (and
one diffusion models do not share as directly, since the diffusion
side's analogous likelihood requires the full probability-flow-ODE
construction of \Cref{thm:pf-ode-scratch}, an extra derived object rather
than something immediately available from the sampler already being run).

\begin{caution}[Not implemented in this codebase]
\label{caution:likelihood-not-implemented}
Per the code-research pass underlying this document, \textbf{no code path
in either \file{sampling.py} or \file{sigma\_flow\_generator.py} computes
$\nabla\cdot u_\theta$, a Hutchinson trace estimate, or any accumulated
log-density along a sampled trajectory}; no likelihood or confidence score
derived from \Cref{cor:instantaneous-cov} is produced anywhere in the
sampling pipeline reconstructed in \Cref{sec:sigmaflow-sampling}. SigmaFlow's
actual ranking mechanism, where used, is inherited unchanged from
SigmaDock's non-learned heuristic (\Cref{sec:sigmadock-ranking}), per
\Cref{tab:sigmadock-vs-sigmaflow-precise}. \Cref{cor:instantaneous-cov}'s
construction is recorded in this document because it is a genuine,
theoretically well-founded \emph{capability} that the flow-matching
formulation --- unlike the diffusion formulation --- would make relatively
straightforward to add (a scalar ODE run alongside the existing vector
field, using an already-trained $u_\theta$ with no retraining required),
and because \Cref{chap:docking}'s discussion of confidence estimation as a
distinct task from pose generation
(\Cref{sec:docking-tasks}) specifically named this as an open, unimplemented
direction; it should not be read as a description of current SigmaFlow
behaviour.
\end{caution}

\section{A tested, and rejected, loss variant}
\label{sec:variant-time-weighting}

We record, briefly and for completeness given
\Cref{sec:sigmaflow-nodressing}'s discussion of the ``no time-dependent
scaling needed'' design choice, one concrete alternative this project's
broader experimental work tested empirically, since a reader who has just
followed \Cref{sec:sde-perspective}'s discussion of diffusion's
$\lambda(s)$ weighting might reasonably ask whether an analogous
reweighting should be added to SigmaFlow's (currently unweighted, per
\Cref{sec:sigmaflow-loss}) flow-matching loss.

\begin{caution}[Empirically tested, no benefit found; not part of this document's core theory]
\label{caution:time-weighting-empirical}
A variant multiplying \Cref{sec:sigmaflow-loss}'s per-fragment squared
errors by a factor $1/(1-t)^2$ (capped near $t=1$ for numerical stability)
was implemented, trained, and evaluated as part of this project's broader
experimental investigation. Two findings from that investigation are
relevant here, stated precisely rather than left vague: \textbf{(i)}
this reweighting has \emph{no} grounding in the flow-matching theory
developed in \Cref{chap:flow-matching,chap:riemannian-fm} of this document
--- unlike diffusion's $\lambda(s)$, which \Cref{sec:sde-perspective} derives
as compensating for the conditional score's genuinely $t$-varying scale,
\Cref{sec:sigmaflow-nodressing}'s constant-in-$t$ conditional targets have
no analogous scale variation to compensate for, so this variant is, by this
document's own theory, an unmotivated modification rather than a
theoretically indicated correction; \textbf{(ii)} consistent with point
(i), it was found, empirically, to produce no measurable improvement (and
if anything a marginal degradation) in the fragment-boundary bond-length
metric this project's broader work uses to quantify exactly the failure
mode \Cref{sec:triangulation-precise}'s discussion anticipates. We record
this only to close the loop on a natural question the theory of this
chapter raises, not as part of SigmaFlow's actual, shipped training
recipe.
\end{caution}

\clearpage

% ================================================================
\part{Synthesis}
\label{part:synthesis}
% ================================================================

\chapter{Implementation Details, Numerical Stability, and Consistency Checks}
\label{chap:numerics}

This chapter consolidates, in one place, the numerical-implementation
details that \Cref{chap:groups,chap:diffusion,chap:diffusion-so3,chap:equiformer,chap:sigmadock,chap:sigmaflow}
introduced individually, at the point each became relevant, and adds the
handful that are more naturally stated only once every other piece of
machinery is available. None of the content here is new mathematics; it is
an index and a systematic check, in the spirit of
\Cref{sec:dimension-checks}'s closing consistency pass.

\section{Every clamp, epsilon, and numerical safeguard in one table}
\label{sec:clamp-table}

\begin{center}
\begin{longtable}{p{3.4cm}p{2.6cm}p{7.3cm}}
\toprule
Location & Guard & Purpose (cross-reference) \\
\midrule
\endfirsthead
\toprule
Location & Guard & Purpose \\
\midrule
\endhead
\pyfunc{Omega} (angle from trace) & $\mathrm{clamp}(\cdot,-1{+}10^{-7},1{-}10^{-7})$ (SigmaFlow); $(-0.99,0.99)$ (original SigmaDock, now recognised as a bug for the flow-matching use case) & \pyfunc{arccos} domain safety; the wider original clamp additionally, and incorrectly, truncated true angles $>172^\circ$ (\Cref{caution:omega-clamp}) \\
\pyfunc{rotation\_vector\_from\_matrix} & soft blend (\pyfunc{isclose}, $\mathrm{atol}{=}10^{-8}$) near $\theta{=}0$; hard/loose blend ($\mathrm{atol}{=}10^{-2}$) near $\theta{=}\pi$ & avoids $\sin\theta{=}0$ division in $\log(R) = \tfrac{\theta}{2\sin\theta}(R{-}R^\top)$ (\Cref{sec:so3-log-scratch}) \\
\pyfunc{d\_log\_f\_d\_omega} (IGSO(3) score derivative) & hard branch at $\sigma^2<0.3$: literal ratio vs.\ approximation $-\omega/\sigma^2$ & literal $\partial_\omega f/f$ numerically ill-conditioned at small noise (\Cref{sec:diff-so3-numerics}); small-noise limit asserted by the code, not independently re-derived here \\
\pyfunc{get\_fragment\_mass\_inertia} & eigenvalue floor $\max(10^{-8}\tr(I)/3,10^{-12})$, plus $+10^{-8}I$ Tikhonov term before the linear solve & protects \pyfunc{torch.linalg.solve} against near-degenerate (near-linear, 2--3-atom) fragment inertia tensors (\Cref{sec:newton-euler-precise}) \\
\pyfunc{newton\_maruyama} & \texttt{omega}$=\mathrm{hat}(dW)$ clamped entrywise to $[-10^3,10^3]$ & loose safety net against a diverging angular-velocity solve, not a tight physical bound \\
\pyfunc{d\_f\_igso3\_d\_omega}, series truncation & fixed truncation at $L{=}1000$ terms; angle grid excludes $\omega{=}0$ exactly & \eqref{eq:igso3-scratch} is an infinite series (\Cref{rem:igso3-truncation}); $\sin(\omega/2){=}0$ singularity at $\omega{=}0$ avoided by grid construction, not a runtime guard \\
\pyfunc{BesselBasis} (radial embedding) & $+\eps$ in denominator $d_{\mathrm{scaled}}+\eps$ & avoids $0/0$ as an edge length $\to0$ (\Cref{sec:gnn-scratch}) \\
\pyfunc{PolynomialCutoff} & indicator $\ind[d<r_{\max}]$, degree-8 envelope vanishing (with 2 derivatives) at $r_{\max}$ & continuity of messages as an edge crosses the cutoff radius (\Cref{sec:gnn-scratch}) \\
Assorted \pyfunc{vee}, \pyfunc{expmap} & \emph{print}-only warnings (not raised exceptions) on a non-skew-symmetric input & a soft, non-blocking sanity check --- callers can silently pass a malformed matrix without the check stopping execution (\Cref{def:hatvee-scratch}'s code correspondence) \\
\texttt{assert M.min()} $\ge 2$ & hard assertion, caught upstream & fragment mass (atom count, \Cref{def:fragment-mass-precise}) degeneracy, not a chemical-mass concern \\
\bottomrule
\end{longtable}
\end{center}

\section{Masking, batching, and scatter conventions}
\label{sec:masking-batching}

Every per-fragment or per-molecule reduction encountered in
\Cref{chap:sigmadock,chap:sigmaflow} (net force/torque, mass/inertia,
per-molecule loss aggregation) uses the same underlying mechanism, worth
stating once explicitly since it is assumed silently throughout: a flat
tensor over \emph{all} atoms (or fragments) in a batch, together with an
integer index tensor recording which molecule/fragment each row belongs to,
reduced via a \texttt{scatter}/\texttt{index\_add} operation keyed on that
index. This is the standard \texttt{torch\_geometric} convention for
representing a batch of variable-sized graphs as one large disjoint-union
graph rather than as a padded, fixed-size tensor, and it is why, throughout
\Cref{chap:sigmadock,chap:sigmaflow}, a formula written ``per fragment''
(e.g.\ \eqref{eq:newton-euler-precise}) is, in the actual tensor
implementation, a single vectorised scatter-sum over the entire batch's
flattened atom list, not a Python-level loop over fragments. Masking (e.g.\
excluding virtual/dummy atoms from the physical mass/inertia computation,
\Cref{sec:newton-euler-precise}) is implemented via boolean indexing keyed
on the categorical node-entity type (\pyfunc{HPARAMS.get\_node\_idx}), not
via a separate, redundant tensor.

\section{Floating-point precision}
\label{sec:precision}

Several routines identified in earlier chapters explicitly upcast to
\texttt{float64} (or, on Apple-Silicon/MPS backends where
\texttt{float64} support is limited, round-trip through the CPU) before a
numerically sensitive computation and cast back afterward: \pyfunc{Log},
\pyfunc{Omega} (\Cref{sec:so3-scratch}'s code correspondences). This
reflects that the angle/logarithm computation, involving \pyfunc{arccos}
near its domain boundary and ratios like $\theta/\sin\theta$ near $\theta=0$,
is more sensitive to rounding error than the surrounding
\texttt{float32}-precision network computation that dominates the rest of
the pipeline; training and sampling otherwise run at \texttt{float32}
(\file{conf/sampling/base.yaml}'s \texttt{hardware.precision:32} default).

\section{Dimensional and consistency checks}
\label{sec:dimension-checks}

As a final, explicit sanity pass --- in the spirit this document has
maintained throughout, of tying every formula to a concrete tensor shape
--- we record the shape/transformation-behaviour consistency of the
central quantities of \Cref{chap:sigmaflow}, gathered in one place.

\begin{center}
\begin{tabular}{@{}lll@{}}
\toprule
Quantity & Shape & Transformation under global $Q\in SO(3)$ \\
\midrule
$\mathrm{pos}_0,\mathrm{pos}_t$ (atom coords) & \shape{N,3} & $x\mapsto Qx$ (world-frame, polar) \\
$\mathrm{trans}_1,\mathrm{trans}_t$ (fragment COM) & \shape{$\Sigma F$,3} & $x\mapsto Qx$ (world-frame, polar) \\
$R_1,R_t$ (fragment orientation) & \shape{$\Sigma F$,3,3} & $R\mapsto QR$ (world-frame) \\
$u_t^{\mathrm{trans}}$, $\mathrm{pred\_u\_t\_trans}$ & \shape{$\Sigma F$,3} & $v\mapsto Qv$ (polar, \Cref{eq:equivariance-concrete}) \\
$u_t^R$, $\mathrm{pred\_u\_t\_R}$ & \shape{$\Sigma F$,3,3} & world-frame $\so(3)$-valued (per \Cref{sec:newton-euler-precise}'s reading); axial parity (\Cref{rem:parity-scratch}) \\
$h_i$ (steerable node feature) & \shape{N, 16, 256} ($L{=}3$, \texttt{sphere\_channels}$=256$) & $h^{\ell,c}\mapsto D^\ell(Q)h^{\ell,c}$, degreewise \\
Network output (pre-slice) & \shape{N\_atoms, 16} & type-1 slice, indices 1--3, is the pseudo-force \\
Loss $\texttt{loss\_trans},\texttt{loss\_R}$ & \shape{$\Sigma F$} & scalar (invariant), per fragment, before batch reduction \\
\bottomrule
\end{tabular}
\end{center}

Every row is either a direct restatement of a code correspondence already
given in full in an earlier chapter (cross-referenced there) or an
immediate consequence of \Cref{prop:wigner-properties}(iv)'s identification
$D^1(Q)=Q$; we do not introduce any new claim in this table, only collect
what has already been derived, as a single point of reference for
verifying a specific tensor's shape or transformation law without
re-reading the chapter it first appeared in.

\section{Summary of open items}
\label{sec:open-items-summary}

In keeping with this document's stated policy of flagging rather than
guessing, we collect, in one place, every point flagged elsewhere in this
document as not fully, independently verified from the code (each is
discussed, with full context, at the cross-referenced location):
\begin{itemize}
\item The exact FR3D stochastic merge algorithm and stopping criterion
(\Cref{sec:sigmadock-fr3d}'s caution).
\item The precise semantic identity of each of the ten categorical atom
feature slots (\Cref{sec:node-embedding}'s code correspondence).
\item The exact formula of the \texttt{rms\_norm\_sh} normalisation layer
and the sinusoidal time-embedding (\Cref{sec:equiformer-linear-norm},
\Cref{sec:node-embedding}).
\item Whether the small-noise limit $-\omega/\sigma^2$ used by
\pyfunc{d\_log\_f\_d\_omega} has been independently re-derived from
\eqref{eq:igso3-scratch} (it has not, in this document; it is reported as
the code's own documented claim, \Cref{sec:diff-so3-numerics}).
\item Which specific value of \texttt{rot\_score\_method}/\texttt{rot\_score\_scaling}
(\Cref{sec:sd-head-branches}) was used to train the specific checkpoints
referenced in this project's broader experimental work, as opposed to the
constructor defaults recorded here.
\item The \texttt{fragment\_scaling} default-value discrepancy between two
source locations (\Cref{sec:sigmadock-loss}).
\end{itemize}
None of these open items affects the mathematical derivations of
\Cref{chap:groups,chap:steerable,chap:diffusion,chap:diffusion-so3,chap:flow-matching,chap:riemannian-fm},
which are self-contained proofs independent of any specific code detail;
they affect only the precision with which a small number of secondary
implementation facts are pinned down.

\chapter*{Closing Remarks}
\addcontentsline{toc}{chapter}{Closing Remarks}

This document set out to make one substitution precise:
\Cref{tab:sigmadock-vs-sigmaflow-precise}'s claim that converting SigmaDock
into SigmaFlow replaces a probability path and a regression target while
leaving the state space, the fragmentation and triangulation geometry, the
input graph, and the entire EquiformerV2 backbone untouched. Making that
claim \emph{precise} required building, from elementary analysis and linear
algebra upward, the full apparatus of Riemannian manifolds and Lie groups
(\Cref{chap:groups}), the representation theory of $SO(3)$
(\Cref{chap:steerable}), score-based diffusion in Euclidean space and on
$SO(3)$ (\Cref{chap:diffusion,chap:diffusion-so3}), the equivariant
message-passing architecture that both methods share
(\Cref{chap:pre-equiformer,chap:equiformer}), and continuous-normalizing-flow
theory in its Euclidean and Riemannian forms
(\Cref{chap:flow-matching,chap:riemannian-fm}) --- at which point the
substitution itself (\Cref{chap:sigmaflow}) could be stated as a small,
precisely delimited set of changes against the common baseline
(\Cref{chap:sigmadock}) that the preceding chapters established.

A reader who has worked through this document in order should now be able
to: state and prove why a fragment-based parameterisation, rather than a
torsional one, gives a decoupled noising/flow measure
(\Cref{thm:entanglement-scratch}); derive the $\mathrm{IGSO}(3)$ heat
kernel and its score from the Peter--Weyl decomposition of the
Laplace--Beltrami operator on $SO(3)$ (\Cref{sec:igso3-derivation}); derive
the geodesic flow-matching interpolant on $SO(3)$ and prove the exact
algebraic identity between its two computational forms
(\Cref{prop:so3-current-point}); explain, from Schur's lemma, exactly which
operations an equivariant neural network layer is and is not permitted to
perform (\Cref{prop:schur-scratch,prop:multichannel-linear}); and, for
every one of these constructions, point to the specific file, function,
tensor shape, and line of reasoning connecting the mathematics to the
program that actually runs on real protein--ligand complexes.

Where this document's own derivations and the codebase's own comments
disagreed, or where a detail could not be pinned down from a static reading
of the source, this document has said so explicitly rather than resolving
the disagreement by assumption --- most notably the corrected
$SO(3)$-angle clamp (\Cref{caution:omega-clamp}), the reversed internal
time conventions between the two methods
(\Cref{sec:sigmaflow-time-convention}), the hybrid V1/V2 attention
mechanism actually shipped (\Cref{caution:hybrid-architecture}), and the
always-identity fragment reference rotation
(\Cref{rem:r1-identity-simplification}). These are exactly the kind of
facts a purely paper-level account of either method would not surface, and
locating them is, alongside the from-scratch mathematical development
itself, the primary contribution of treating this material as a full
monograph rather than a summary.

\end{document}
